\chaptertitle{A Mu Offering}

                                 A Mu Offering
   The Tortoise and Achilles have just been to hear a lecture on the origins of the
       Genetic Code, and are now drinking some tea at Achilles' home.

Achilles: I have something terrible to confess, Mr. T.
Tortoise: What is it, Achilles?
Achilles: Despite the fascinating subject matter of that lecture, I drifter to sleep a time or
   two. But in my drowsy state, I still was semi-awake aware of the words coming into
   my ears. One strange image that floated up from my lower levels was that `A' and `T',
   instead of standing "adenine" and "thymine", stood for my name and yours-and
   double-strands of DNA had tiny copies of me and you along backbones, always
   paired up, just as adenine and thymine always Isn't that a strange symbolic image?
Tortoise: Phooey! Who believes in that silly kind of stuff? Anyway, about `C' and `G'?
Achilles: Well, I suppose `C' could stand for Mr. Crab, instead o cytosine. I'm not sure
   about `G', but I'm sure one could thin something. Anyway, it was amusing to imagine
   my DNA being with minuscule copies of you-as well as tiny copies of myself, for
   matter. Just think of the infinite regress THAT leads to!
Tortoise: I can see you were not paying too much attention to the lecture.
Achilles: No, you're wrong. I was doing my best, only I had a hard keeping fancy
   separated from fact. After all, it is such a strange netherworld that those molecular
   biologists are exploring.
Tortoise: How do you mean?
Achilles: Molecular biology is filled with peculiar convoluted loops which I can't quite
   understand, such as the way that folded proteins, which are coded for in DNA, can
   loop back and manipulate the DNA which came from, possibly even destroying it.
   Such strange loops always confuse the daylights out of me. They're eerie, in a way.
Tortoise: I find them quite appealing.
Achilles: You would, of course-they're just down your alley. But me, sometimes I like to
   retreat from all this analytic thought any meditate a little, as an antidote. It clears my
   mind of all those conf loops and incredible complexities which we were hearing about
   tonight.
Tortoise: Fancy that. I wouldn't have guessed that you were a meditator
Achilles: Did I never tell you that I am studying Zen Buddhism?
Tortoise: Heavens, how did you come upon that?
Achilles: I have always had a yen for the yin and yang, you know – the




A Mu Offering                                                                             239

   whole Oriental mysticism trip, with the I Ching, gurus, and whatnot. So one day I'm
   thinking to myself, "Why not Zen too?" And that's how it all began.
Tortoise: Oh, splendid. Then perhaps I can finally become enlightened. Achilles: Whoa,
   now. Enlightenment is not the first step on the road to Zen; if anything, it'. the last
   one! Enlightenment is not for novices like you, Mr. T!
Tortoise: I see we have had a misunderstanding. By "enlightenment", I hardly meant
   something so weighty as is meant in Zen. All I meant is that I can perhaps become
   enlightened as to what Zen is all about. Achilles: For Pete's sake, why didn't you say
   so? Well, I'd be only too happy to tell you what I know of Zen. Perhaps you might
   even be tempted to become a student of it, like me.
Tortoise: Well, nothing's impossible.
Achilles: You could study with me under my master, Okanisama-the seventh patriarch.
Tortoise: Now what in the world does that mean?
Achilles: You have to know the history of Zen to understand that.
Tortoise: Would you tell me a little of the history of Zen, then?
Achilles: An excellent idea. Zen is a kind of Buddhism which was founded by a monk
   named Bodhidharma, who left India and went to China around the sixth century.
   Bodhidharma was the first patriarch. The sixth one was Eno. (I've finally got it
   straight now!)
Tortoise: The sixth patriarch was Zeno, eh? I find it strange that he, of all people, would
   get mixed up in this business.
Achilles: I daresay you underestimate the value of Zen. Listen just a little more, and
   maybe you'll come to appreciate it. As I was saying, about five hundred years later,
   Zen was brought to Japan, and it took hold very well there. Since that time it has been
   one of the principal religions in Japan.
Tortoise: Who is this Okanisama, the "seventh patriarch"?
Achilles: He is my master, and his teachings descend directly from those of the sixth
   patriarch. He has taught me that reality is one, immutable, and unchanging; all
   plurality, change, and motion are mere illusions of the senses.
Tortoise: Sure enough, that's Zeno, a mile away. But how ever did he come to be tangled
   up in Zen? Poor fellow!
Achilles: Whaaat? I wouldn't put it that way. If ANYONE is tangled up, it's ... But that's
   another matter. Anyway, I don't know the answer to your question. Instead, let me tell
   you something of the teachings of my master. I have learned that in Zen, one seeks
   enlightenment, or SATORI-the state of "No-mind". In this state, one does not think
   about the world-one just is. I have also learned that a student of Zen is not supposed
   to "attach" to any object or thought or person-which is to say, he must not believe in,
   or depend on, any absolute-not even this philosophy of nonattachment.




A Mu Offering                                                                          240

Tortoise: Hmm ... Now THERE'S something I could like about Achilles: I had a hunch
   you'd get attached to it.
Tortoise: But tell me: if Zen rejects intellectual activity, does it make sense to
   intellectualize about Zen, to study it rigorously?
Achilles: That matter has troubled me quite a bit. But I think I have finally worked out an
   answer. It seems to me that you may begin approaching Zen through any path you
   know-even if it is completely antithetical to Zen. As you approach it, you gradually
   learn to stray from that path. The more you stray from the path, the closer you get to
   Zen.
Tortoise: Oh, it all begins to sound so clear now.
Achilles: My favorite path to Zen is through the short, fascinating and weird Zen parables
   called "koans".
Tortoise: What is a koan?
Achilles: A koan is a story about Zen masters and their student times it is like a riddle;
   other times like a fable; and other ti nothing you've ever heard before.
Tortoise: Sounds rather intriguing. Would you say that to read al koans is to practice
   Zen?
Achilles: I doubt it. However, in my opinion, a delight in koans million times closer to
   real Zen than reading volume after about Zen, written in heavy philosophical jargon.
Tortoise: I would like to hear a koan or two.
Achilles: And I would like to tell you one-or a few. Perhaps begin with the most famous
   one of all. Many centuries ago, the Zen master named Joshu, who lived to be 119
   years old.
Tortoise: A mere youngster!
Achilles: By your standards, yes. Now one day while Joshu and monk were standing
   together in the monastery, a dog wand The monk asked Joshu, "Does a dog have
   Buddha-nature,
Tortoise: Whatever that is. So tell me-what did Joshu reply?
Achilles: 'MU'.
Tortoise: 'MU? What is this 'MU'? What about the dog? What about Buddha-nature?
   What's the answer?
Achilles: Oh, but 'MU' is Joshu's answer. By saying 'MU', Joshu let the other monk know
   that only by not asking such questions can one know the answer to them.
Tortoise: Joshu "unasked" the question.
Achilles: Exactly!
Tortoise: 'MU' sounds like a handy thing to have around. I'd like unask a question or two,
   sometimes. I guess I'm beginning to get the hang of Zen. Do you know any other
   koans, Achilles? I would like to hear some more.
Achilles: My pleasure. I can tell you a pair of koans which go together
   Only ...
Tortoise: What's the matter?




A Mu Offering                                                                          241

Achilles: Well, there is one problem. Although both are widely told koans, my master has
   cautioned me that only one of them is genuine. And what is more, he does not know
   which one is genuine, and which one is a fraud.
Tortoise: Crazy! Why don't you tell them both to me and we can speculate to our hearts'
   content!
Achilles: All right. One of the alleged koans goes like this:

                        A monk asked Baso: "What is Buddha?"
                          Baso said: "This mind is Buddha."

Tortoise: Hmm ... "This mind is Buddha"? Sometimes I don't quite understand what these
   Zen people are getting at. Achilles: You might prefer the other alleged koan then.
Tortoise: How does it run? Achilles: Like this:

                        A monk asked Baso: "What is Buddha?"
                         Baso said: "This mind is not Buddha."

Tortoise: My, my! If my shell isn't green and not green! I like that! Achilles: Now, Mr. T-
   you're not supposed to just "like" koans.
Tortoise: Very well, then-I don't like it.
Achilles: That's better. Now as I was saying, my master believes only one of the two is
   genuine.
Tortoise: I can't imagine what led him to such a belief. But anyway, I suppose it's all
   academic, since there's no way to know if a koan is genuine or phony.
Achilles: Oh, but there you are mistaken. My master has shown us how to do it.
Tortoise: Is that so? A decision procedure for genuineness of koans? I should very much
   like to hear about THAT.
Achilles: It is a fairly complex ritual, involving two stages. In the first stage, you must
   TRANSLATE the koan in question into a piece of string, folded all around in three
   dimensions.
Tortoise: That's a curious thing to do. And what is the second stage?
Achilles: Oh, that's easy-all you need to do is determine whether the string has Buddha-
   nature, or not! If it does, then the koan is genuine-if not, the koan is a fraud.
Tortoise: Hmm ... It sounds as if all you've done is transfer the need for a decision
   procedure to another domain. Now it's a decision procedure for Buddha-nature that
   you need. What next? After all, if you can't even tell whether a Do(; has Buddha-
   nature or not, how can you expect to do so for every possible folded string?
Achilles: Well, my master explained to me that shifting between domains can help. It's
   like switching your point of view. Things sometimes look complicated from one
   angle, but simple from another. He gave the example of an orchard, in which from
   one direction no order is




A Mu Offering                                                                          242

        FIGURE 45. La Mezquita, by M. C. Escher (black and white chalk, 1936

    apparent, but from special angles, beautiful regularity em, You've reordered the same
   information by changing your way of looking at it.
Tortoise: I see. So perhaps the genuineness of a koan is concealed how very deeply inside
   it, but if you translate it into a string it ma in some way to float to the surface?
Achilles: That's what my master has discovered.
Tortoise: Then I would very much like to learn about the techniqu first, tell me: how can
   you turn a koan (a sequence of words) folded string (a three-dimensional object)?
   They are rather dif kinds of entities.
Achilles: That is one of the most mysterious things I have learned i There are two steps:
   "transcription" and "translation". TRANSCF a koan involves writing it in a phonetic
   alphabet, which contain four geometric symbols. This phonetic rendition of the koan
   is called the MESSENGER.
Tortoise: What do the geometric symbols look like?
Achilles: They are made of hexagons and pentagons. Here is what they




A Mu Offering                                                                        243

   look like (picks up a nearby napkin, and draws for the Tortoise these four figures):




Tortoise: They are mysterious-looking.
Achilles: Only to the uninitiated. Now once you have made the messenger, you rub your
   hands in some ribo, and
Tortoise: Some ribo? Is that a kind of ritual anointment?
Achilles: Not exactly. It is a special sticky preparation which makes the string hold its
   shape, when folded up. Tortoise: What is it made of?
Achilles: I don't know, exactly. But it feels sort of gluey, and it works exceedingly well.
   Anyway, once you have some ribo on your hands, you can TRANSLATE the
   sequence of symbols in the messenger into certain kinds of folds in the string. It's as
   simple as that. Tortoise: Hold on! Not so fast! How do you do that?
Achilles: You begin with the string entirely straight. Then you go to one end and start
   making folds of various types, according to the geometric symbols in the messenger.
Tortoise: So each of those geometric symbols stands for a different way to curl the string
   up?
Achilles: Not in isolation. You take them three at a time, instead of one at a time. You
   begin at one end of the string, and one end of the messenger. What to do with the first
   inch of the string is determined by the first three geometric symbols. The next three
   symbols tell you how to fold the second inch of string. And so you inch your way
   along the string and simultaneously along the messenger, folding each little segment
   of string until you have exhausted the messenger. If you have properly applied some
   ribo, the string will keep its folded shape, and what you thereby produce is the
   translation of the koan into a string.
Tortoise: The procedure has a certain elegance to it. You must get some wild-looking
   strings that way.
Achilles: That's for sure. The longer koans translate into quite bizarre shapes.
Tortoise: I can imagine. But in order to carry out the translation of the messenger into the
   string, you need to know what kind of fold each triplet of geometric symbols in the
   messenger stands for. How do you know this? Do you have a dictionary?
Achilles: Yes-there is a venerated book which lists the "Geometric Code". If you don't
   have a copy of this book, of course, you can't translate a koan into a string.




A Mu Offering                                                                           244

Tortoise: Evidently not. What is the origin of the Geometric Code Achilles: It came from
   an ancient master known as "Great Tutor" who my master says is the only one ever to
   attain the Enlightenment ' Enlightenment.
Tortoise: Good gravy! As if one level of the stuff weren't enough. But then there are
   gluttons of every sort-why not gluttons for enlighten] Achilles: Do you suppose that
   "Enlightenment 'Yond Enlighten] stands for "EYE"?
Tortoise: In my opinion, it's rather doubtful that it stands for you, Ac More likely, it
   stands for "Meta-Enlightenment"-"ME", that is
Achilles: For you? Why would it stand for you? You haven't even re; the FIRST stage of
   enlightenment, let alone the
Tortoise: You never know, Achilles. Perhaps those who have learn( lowdown on
   enlightenment return to their state before enlighten I've always held that "twice
   enlightened is unenlightened." But le back to the Grand Tortue-uh, I mean the Great
   Tutor.
Achilles: Little is known of him, except that he also invented the Art of Zen Strings.
Tortoise: What is that?
Achilles: It is an art on which the decision procedure for Buddha-nature is based. I shall
   tell you about it.
Tortoise: I would be fascinated. There is so much for novices like absorb!
Achilles: There is even reputed to be a koan which tells how the Art Strings began. But
   unfortunately, all this has long since been lost sands of time, and is no doubt gone
   forever. Which may be just a for otherwise there would be imitators who would take
   on the m~ name, and copy him in other ways.
Tortoise: But wouldn't it be a good thing if all students of Zen copied that most
   enlightened master of all, the Great Tutor?
Achilles: Let me tell you a koan about an imitator.

   Zen master Gutei raised his finger whenever he was asked a question about Zen. A
   young novice began to irritate him in this way. When Gut was told about the
   novice's imitation, he sent for him and asked him if were true. The novice
   admitted it was so. Gutei asked him if he understood. In reply the novice held up
   his index finger. Gutei promptly cut off. The novice ran from the room, howling in
   pain. As he reached it threshold, Gutei called, "Boy!" When the novice turned,
   Gutei raised h index finger. At that instant the novice vas enlightened.

Tortoise: Well, what do you know! Just when I thought Zen was all about Joshu and his
   shenanigans, now I find out that Gutei is in on the merriment too. He seems to have
   quite a sense of humor.

Achilles: That koan is very serious. I don't know how you got the idea that it is
   humorous.
Tortoise: Perhaps Zen is instructive because it is humorous. I would guess




A Mu Offering                                                                         245

   that if you took all such stories entirely seriously, you would miss the point as often as
   you would get it.
Achilles: Maybe there's something to your Tortoise-Zen.
Tortoise: Can you answer just one question for me? I would like to know this: Why did
   Bodhidharma come from India into China?
Achilles: Oho! Shall I tell you what Joshu said when he was asked that very question?
Tortoise: Please do.
Achilles: He replied, "That oak tree in the garden."
Tortoise: Of course; that's just what I would have said. Except that I would have said it in
   answer to a different question-namely, "Where can I find some shade from the
   midday sun?"
Achilles: Without knowing it, you have inadvertently hit upon one of the basic questions
   of all Zen. That question, innocent though it sounds, actually means, "What is the
   basic principle of Zen?"
Tortoise: How extraordinary. I hadn't the slightest idea that the central aim of Zen was to
   find some shade.
Achilles: Oh, no-you've misunderstood me entirely. I wasn't referring to THAT question.
   I meant your question about why Bodhidharma came from India into China.
Tortoise: I see. Well, I had no idea that I was getting into such deep waters. But let's
   come back to this curious mapping. I gather that any koan can be turned into a folded
   string by following the method you outlined. Now what about the reverse process?
   Can any folded string be read in such a way as to yield a koan?
Achilles: Well, in a way. However .. .
Tortoise: What's wrong?
Achilles: You're just not supposed to do it that way 'round. It would violate the Central
   Dogma of Zen strings, you see, which is contained in this picture (picks up a napkin
   and draws):

                        koan    =>      messenger     folded string
                                 transcription translation

   You're not supposed to go against the arrows-especially not the second one.
Tortoise: Tell me, does this Dogma have Buddha-nature, or not? Come to think of it, I
   think I'll unask the question. Is that all right?
Achilles: I am glad you unasked the question. But-I'll let you in on a secret. Promise you
   won't tell anyone?
Tortoise: Tortoise's honor.
Achilles: Well, once in a while, I actually do go against the arrows. I get sort of an illicit
   thrill out of it, I guess.
Tortoise: Why, Achilles! I had no idea you would do something so irreverent!
Achilles: I've never confessed it to anyone before-not even Okanisama.




A Mu Offering                                                                             246

Tortoise: So tell me, what happens when you go against the arrows i Central Dogma?
   Does that mean you begin with a string and m koan?
Achilles: Sometimes-but some weirder things can happen.
Tortoise: Weirder than producing koans?
Achilles: Yes ... When you untranslate and untranscribe, you get THING, but not always
   a koan. Some strings, when read out Ion way, only give nonsense.
Tortoise: Isn't that just another name for koans?
Achilles: You clearly don't have the true spirit of Zen yet.
Tortoise: Do you always get stories, at least?
Achilles: Not always-sometimes you get nonsense syllables, other you get ungrammatical
   sentences. But once in a while you get seems to be a koan.
Tortoise: It only SEEMS to be one?
Achilles: Well, it might be fraudulent. you see.
Tortoise: Oh, of course.
Achilles: I call those strings which yield apparent koans "well-foi strings.
Tortoise: Why don't you tell me about the decision procedure which allows you to
   distinguish phony koans from the genuine article?
Achilles: That's what I was heading towards. Given the koan, or non• as the case may be,
   the first thing is to translate it into the dimensional string. All that's left is to find out
   if the strip Buddha-nature or not.
Tortoise: But how do you do THAT?
Achilles: Well, my master has said that the Great Tutor was able, I glancing at a string, to
   tell if it had Buddha-nature or not.
Tortoise: But what if you have not reached the stage of the Enlightenment: 'Yond
   Enlightenment? Is there no other way to tell if a string hasi Buddha-nature?
Achilles: Yes, there is. And this is where the Art of Zen Strings come is a technique for
   making innumerably many strings, all of whit Buddha-nature.
Tortoise: You don't say! And is there a corresponding way of n strings which DON'T
   have Buddha-nature?
Achilles: Why would you want to do that?
Tortoise: Oh, I just thought it might be useful.
Achilles: You have the strangest taste. Imagine! Being more intere things that DON'T
   have Buddha-nature than things that DO!
Tortoise: Just chalk it up to my unenlightened state. But go on. T how to make a string
   which DOES have Buddha-nature.
Achilles: Well, you must begin by draping a loop of string over your in one of five legal
   starting positions, such as this one ... (Picks up a string and drapes it in a simple loop
   between a finger on each hand.:)




A Mu Offering                                                                               247

Tortoise: What are the other four legal starting positions?
Achilles: Each one is a position considered to be a SELF-EVIDENT manner of picking
   up a string. Even novices often pick up strings in those positions. And these five
   strings all have Buddha-nature. Tortoise: Of course.
Achilles: Then there are some String Manipulation Rules, by which you can make more
   complex string figures. In particular, you are allowed to modify your string by doing
   certain basic motions of the hands. For instance, you can reach across like this-and
   pull like this-and twist like this. With each operation you are changing the overall
   configuration of the string draped over your hands.
Tortoise: Why, it looks just like making cat's-cradles and such string figures!
Achilles: That's right. Now as you watch, you'll see that some of these rules make the
   string more complex; some simplify it. But whichever way you go, as long as you
   follow the String Manipulation Rules, every string you produce will have Buddha-
   nature.
Tortoise: That is truly marvelous. Now what about the koan concealed inside this string
   you've just made? Would it be genuine?
Achilles: Why, according to what I've learned, it must. Since I made it according to the
   Rules, and began in one of the five self-evident positions, the string must have
   Buddha-nature, and consequently it must correspond to a genuine koan.
Tortoise: Do you know what the koan is?
Achilles: Are you asking me to violate the Central Dogma? Oh, you naughty fellow!

   (And with furrowed brow and code book in hand, Achilles points along the string
   inch by inch, recording each fold by a triplet of geometric symbols of the strange
   phonetic alphabet for koan, until he has nearly a napkinful.)

   Done!

Tortoise: Terrific. Now let's hear it.
Achilles: All right.

   A traveling monk asked an old woman the road to Taizan, a popular temple
   supposed to give wisdom to the one who worships there. The old woman said:
   "Go straight ahead." After the monk had proceeded a few steps, she said to herself,
   "He also is a common church-goer." Someone told this incident to Joshu, who
   said: "Wait until I investigate." The next day he went and asked the same question,
   and the old woman gave the same answer. Joshu remarked: "I have investigated
   that old woman."

Tortoise: Why, with his flair for investigations, it's a shame that Joshu
never was hired by the FBI. Now tell me-what you did, I could also
do, if I followed the Rules from the Art of Zen Strings, right?
Achilles: Right.
Tortoise: Now would I have to perform the operations in just the same ORDER as you
    did?



A Mu Offering                                                                        248

.Achilles: No, any old order will do.
Tortoise: Of course, then I would get a different string, and consequently a different
   koan. Now would I have to perform the same NUMBER of steps as you did?
Achilles: By no means. Any number of steps is fine.
Tortoise: Well, then there are an infinite number of strings with Buddha nature-and
   consequently an infinite number of genuine koans Howdo you know there is any
   string which CAN "I- be made by your Achilles: Oh, yes-back to things which lack
   Buddha-nature. It just so happens that once you know how to make strings WITH
   Buddha nature, you can also make strings WITHOUT Buddha-nature. That is
   something which my master drilled into me right at the beg Tortoise: Wonderful!
   How does it work?
Achilles: Easy. Here, for example-I'll make a string which lacks Buddha-nature .. .

   (He picks up the string out of which the preceding koan was "pulled", ties a little
   teeny knot at one end of it, pulling it tight with his thumb forefinger.)

This is it -- no Buddha-nature here.

Tortoise: Very illuminating. All it takes is adding a knot? How know that the new string
   lacks Buddha-nature?
Achilles: Because of this fundamental property of Buddha-nature; when two well-formed
   strings are identical but for a knot at one end, then only ONE of them can have
   Buddha-nature. It's a rule of thumb which my master taught me.
Tortoise: I'm just wondering about something. Are there some strings with Buddha-
   nature which you CAN'T reach by following the Rules of Zen Strings, no matter in
   what order?
Achilles: I hate to admit it, but I am a little confused on this point myself. At first my
   master gave the strongest impression that Buddha in a string was DEFINED by
   starting in one of the five legal positions, and then developing the string according to
   the Rules. But then later, he said something about somebody-o "Theorem". I never
   got it straight. Maybe I even misheard said. But whatever he said, it put some doubt in
   my mind as to this method hits ALL strings with Buddha-nature. To the be
   knowledge, at least, it does. But Buddha-nature is a pretty elusive thing, you know.
Tortoise: I gathered as much, from Joshu's 'MU'. I wonder ...
Achilles: What is it?
Tortoise: I was just wondering about those two koans-I mean t and its un-koan-the ones
   which say "This mind is Buddha" at mind is not Buddha"-what do they look like,
   when turned int via the Geometric Code?
Achilles: I'd be glad to show you.




A Mu Offering                                                                          249

   (He writes down the phonetic transcriptions, and then pulls from his pocket a
   couple of pieces of string, which he carefully folds inch by inch, following the
   triplets of symbols written in the strange alphabet. Then he places the finished
   strings side by side.)

You see, here is the difference.

Tortoise: They are very similar, indeed. Why, I do believe there is only one difference
   between them: it's that one of them has a little knot on its end!
Achilles: By Joshu, you're right.
Tortoise: Aha! Now I understand why your master is suspicious.
Achilles: You do?
Tortoise: According to your rule of thumb, AT MOST ONE of such a pair can have
   Buddha-nature, so you know right away that one of the koans must be phony.
Achilles: But that doesn't tell which one is phony. I've worked, and so has my master, at
   trying to produce these two strings by following the String Manipulation Rules, but to
   no avail. Neither one ever turns up. It's quite frustrating. Sometimes you begin to
   wonder ...
Tortoise: You mean, to wonder if either one has Buddha-nature? Perhaps neither of them
   has Buddha-nature-and neither koan is genuine!
Achilles: I never carried my thoughts as far as that-but you're right-it's possible, I guess.
   But I think you should not ask so many questions about Buddha-nature. The Zen
   master Mumon always warned his pupils of the danger of too many questions.
Tortoise: All right-no more questions. Instead, I have a sort of hankering to make a string
   myself. It would be amusing to see if what I come up with is well-formed or not.
Achilles: That could be interesting. Here's a piece of string. (He passes one to the
   Tortoise.)
Tortoise: Now you realize that I don't have the slightest idea what to do.
We'll just have to take potluck with my awkward production, which will follow no rules
   and will probably wind up being completely undecipherable. (Grasps the loop
   between his feet and, with a few simple manipulations, creates a complex string which
   he proffers wordlessly to Achilles. At that moment, Achilles' face lights up.)
Achilles: Jeepers creepers! I'll have to try out your method myself. I have never seen a
   string like this!
Tortoise: I hope it is well-formed. Achilles: I see it's got a knot at one end.
Tortoise: Oh just a moment! May I have it back? I want to do one thing to it.
Achilles: Why, certainly. Here you are.

   (Hands it back to the Tortoise, who ties another knot at the same end. Then the
   Tortoise gives a sharp tug, and suddenly both knots disappear!)




A Mu Offering                                                                            250

Achilles: What happened?
Tortoise: I wanted to get rid of that knot.
Achilles: But instead of untying it, you tied another one, and then BOTH disappeared!
   Where did they go?
Tortoise: Tumbolia, of course. That's the Law of Double Nodulation

   (Suddenly, the two knots reappear from out of nowhere-that is to say, Tumbolia.)

Achilles: Amazing. They must lie in a fairly accessible layer of Tumbol they can pop into
   it and out of it so easily. Or is all of Tumbolia equally inaccessible?
Tortoise: I couldn't say. However, it does occur to me that burning string would make it
   quite improbable for the knots to come back such a case, you could think of them as
   being trapped in a deeper la of Tumbolia. Perhaps there are layers and layers of
   Tumbolia. that's neither here nor there. What I would like to know is how string
   sounds, if you turn it back into phonetic symbols. (As he hauls it back, once again, the
   knots pop into oblivion.)
Achilles: I always feel so guilty about violating the Central Dogma (Takes out his pen
   and code book, and carefully jots down the many sym triplets which correspond to the
   curvy involutions of the Tortoise's string; when he is finished, he clears his voice.)
   Ahem. Are you ready to hear w you have wrought?
Tortoise: I'm willing if you're willing.
Achilles: All right. It goes like this:
A certain monk had a habit of pestering the Grand Tortue (the only one who had ever
   reached the Enlightenment 'Yond Enlightenment), by asking whether various objects
   had Buddha-nature or not. To such questions Tortue invariably sat silent. The monk
   had already asked about a bean, a lake, and a moonlit night. One day, he brought to
   Tortue a piece of string, and asked the same question. In reply, the Grand Tortue
   grasped the loop between his feet and
Tortoise: Between his feet? How odd! Achilles: Why should you find that odd?
Tortoise: Well, ah ... you've got a point there. But please go on!
Achilles: All right.

   The Grand Tortue grasped the loop between     his feet and, with a few simple
   manipulations, created a complex string which he proffered wordlessly to the
   monk. At that moment, the monk was enlightened.

Tortoise: I'd rather be twice-enlightened, personally.
Achilles: Then it tells how to make the Grand Tortue's string, if you be, with a string
   draped over your feet. I'll skip those boring details concludes this way:

   From then on, the monk did not bother Tortue. Instead, he made string after string
   by Tortue's method; and he passed the method on to his own disciples, who passed
   it on to theirs.




A Mu Offering                                                                          251

Tortoise: Quite a yarn. It's hard to believe it was really hidden inside my string.
Achilles: Yet it was. Astonishingly, you seem to have created a well-formed string right
   off the bat.
Tortoise: But what did the Grand Tortue's string look like? That's the main point of this
   koan, I'd suppose.
Achilles: I doubt it. One shouldn't "attach" to small details like that inside koans. It's the
   spirit of the whole koan that counts, not little parts of it. Say, do you know what I just
   realized? I think, crazy though it sounds, that you may have hit upon that long-lost
   koan which describes the very origin of the Art of Zen Strings!
Tortoise: Oh, that would almost be too good to have Buddha-nature.
Achilles: But that means that the great master-the only one who ever reached the mystical
   state of the Enlightenment 'Yond Enlightenment-was named "Tortue", not "Tutor".
   What a droll name!
Tortoise: I don't agree. I think it's a handsome name. I still want to know how Tortue's
   string looked. Can you possibly recreate it from the description given in the koan?
Achilles: I could try ... Of course, I'll have to use my feet, too, since it's described in
   terms of foot motions. That's pretty unusual. But I think I can manage it. Let me give
   it a go. (He picks up the koan and a piece of string, and for a few minutes twists and
   bends the string in arcane ways until he has the finished product.) Well, here it is.
   Odd, how familiar it looks.
Tortoise: Yes, isn't that so? I wonder where I saw it before? Achilles: I know! Why, this
   is YOUR string, Mr. T! Or is Tortoise: Certainly not.
Achilles: Of course not-it's the string which you first handed to me, before you took it
   back to tie an extra knot in it.
Tortoise: Oh, yes-indeed it is. Fancy that. I wonder what that implies.
Achilles: It's strange, to say the least.
Tortoise: Do you suppose my koan is genuine?
Achilles: Wait just a moment ...
Tortoise: Or that my string has Buddha-nature?
Achilles: Something about your string is beginning. to trouble me, Mr.Tortoise .
Tortoise (looking most pleased with himself and paying no attention to Achilles): And
   what about Tortue's string? Does it have Buddha nature? There are a host of questions
   to ask!
Achilles: I would be scared to ask such questions, Mr. T. There is something mighty
   funny going on here, and I'm not sure I like it. Tortoise: I'm sorry to hear it. I can't
   imagine what's troubling you. Achilles: Well, the best way I know to explain it is to
   quote the words of another old Zen master, Kyogen.




A Mu Offering                                                                             252

   Kyogen said: Zen is like a man hanging in a tree by his teeth over a precipice. His
   har grasp no branch, his feet rest on no limb, and under the tree anotl person asks
   him: "Why did Bodhidharma come to China from India?" the man in the tree does
   not answer, he fails; and if he does answer, falls and loses his life. Now what shall
   he do?

Tortoise: That's clear; he should give up Zen, and take up molecular biology.




A Mu Offering                                                                          253


                            Mumon and Gödel

                                    What Is Zen?
I'M NOT SURE I know what Zen is. In a way, I think I understand it very well; but in a
way, I also think I can never understand it at all. Ever since my freshman English teacher
in college read Joshu's MU out loud to our class, I have struggled with Zen aspects of
life, and probably I will never cease doing so. To me, Zen is intellectual quicksand-
anarchy, darkness, meaninglessness, chaos. It is tantalizing and infuriating. And yet it is
humorous, refreshing, enticing. Zen has its own special kind of meaning, brightness, and
clarity. I hope that in this Chapter, I can get some of this cluster of reactions across to
you. And then, strange though it may seem, that will lead us directly to Godelian matters.
         One of the basic tenets of Zen Buddhism is that there is no way to characterize
what Zen is. No matter what verbal space you try to enclose Zen in, it resists, and spills
over. It might seem, then, that all efforts to explain Zen are complete wastes of time. But
that is not the attitude of Zen masters and students. For instance, Zen koans are a central
part of Zen study, verbal though they are. Koans are supposed to be "triggers" which,
though they do not contain enough information in themselves to impart enlightenment,
may possibly be sufficient to unlock the mechanisms inside one's mind that lead to
enlightenment. But in general, the Zen attitude is that words and truth are incompatible,
or at least that no words can capture truth.

Zen Master Mumon
Possibly in order to point this out in an extreme way, the monk Mumon ("No-gate"), in
the thirteenth century, compiled forty-eight koans, following each with a commentary
and a small "poem". This work is called "The Gateless Gate" or the Mumonkan ("No-
gate barrier"). It is interesting to note that the lives of Mumon and Fibonacci coincided
almost exactly: Mumon living from 1183 to 1260 in China, Fibonacci from 1180 to 1250
in Italy. To those who would look to the Mumonkan in hopes of making sense of, or
"understanding", the koans, the Mumonkan may come as a rude shock, for the comments
and poems are entirely as opaque as the koans which they are supposed to clarify. Take
this, for example:' -




Mumon and Gödel                                                                        254

          FIGURE 46. Three Worlds by M. C. Escher (lithograph, 1955)




Mumon and Gödel                                                        255

Koan:
Hogen of Seiryo monastery was about to lecture before dinner when he noticed that the bamboo screen,
lowered for meditation, had not been rolled up. He pointed to it. Two monks arose wordlessly from the
audience and rolled it up. Hogen, observing the physical moment, said, "The state of the first monk is good,
not that of the second."

Mumon's Commentary:
I want to ask you: which of those two monks gained and which lost? If any of you has one eye, he will see
the failure on the teacher's part. However, I am not discussing gain and loss.

Mumon's Poem:
        When the screen is rolled up the great sky opens,
        Yet the sky is not attuned to Zen.
        It is best to forget the great sky
        And to retire from every wind.

Or then again, there is this one:2
        Koan:
        Goso said: "When a buffalo goes out of his enclosure to the edge of the abyss, his horns and his
        head and his hoofs all pass through, but why can't the tail also pass?"

Mumon's Commentary:
If anyone can open one eye at this point and say a word of Zen, he is qualified to repay
the four gratifications, and, not only that, he can save all sentient beings under him. But if
he cannot say such a word of Zen, he should turn back to his tail.
Mumon's Poem:
        If the buffalo runs, he will fall into the trench;
        If he returns, he will be butchered.
        That little tail
        Is a very strange thing.

I think you will have to admit that Mumon does not exactly clear everything up. One
might say that the metalanguage (in which Mumon writes) is not very different from the
object language (the language of the koan). According to some, Mumon's comments are
intentionally idiotic, perhaps meant to show how useless it is to spend one's time in
chattering about Zen. How ever, Mumon's comments can be taken on more than one
level. For instance, consider this :3
Koan:
A monk asked Nansen: "Is there a teaching no master ever taught before?"
Nansen said: "Yes, there is."
"What is it?" asked the monk.
Nansen replied: "It is not mind, it is not Buddha, it is not things."




Mumon and Gödel                                                                                        256

               FIGURE 47. Dewdrop, by M. C. Escher (mezzotint, 1948).

Mumon's Commentary:

Old Nansen gave away his treasure-words. He must have been greatly upset

Mumon's Poem:

       Nansen was too kind and lost his treasure.
       Truly, words have no power.
       Even though the mountain becomes the sea,
       Words cannot open another's mind.

In this poem Mumon seems to be saying something very central to Zen, not making
idiotic statements. Curiously, however, the poem is referential, and thus it is a comment
not only on Nansen’s words, but on its own ineffectiveness. This type of paradox is quite
characteristic of Zen. It is an attempt to "break the mind of logic". You see this paradox
quality in the koan, as well. Concerning Mumon's commentary, do think that Nansen was
really so sure of his answer? Or did the "correctness of his answer matter at all? Or does
correctness play any role in Zen? is the difference between correctness and truth, or is
there any? What if Nansen had said, "No, there is not any such teaching"? Would it have
r any difference? Would his remark have been immortalized in a koan?



Mumon and Gödel                                                                       257

         FIGURE 48. Another World, by M. C. Escher (wood-engraving, 1947).

Here is another koan which aims to break the mind of logic:'

The student Doko came to a Zen master, and said: "I am seeking the truth. In what state
of mind should I train myself, so as to find it?"
        Said the master, "There is no mind, so you cannot put it in any state. There is no
truth, so you cannot train yourself for it."
        "If there is no mind to train, and no truth\_ to find, why do you have these monks
gather before you every day to study- Zen and train themselves for this study?"
        "But I haven't an inch of room here," said the master, "so how could the monks
gather? I have no tongue, so how could I call them together or teach them?"




Mumon and Gödel                                                                       258

       "Oh, how can you lie like this?" asked Doko.
       "But if I have no tongue to talk to others, how can I lie to you?" asked the master.
       Then Doko said sadly, "I cannot follow you. I cannot understand you
       "I cannot understand myself," said the master.

If any koan serves to bewilder, this one does. And most likely, can bewilderment is its
precise purpose, for when one is in a bewildered s one's mind does begin to operate
nonlogically, to some extent. Only by stepping outside of logic, so the theory goes, can
one make the lea enlightenment. But what is so bad about logic? Why does it prevent the
to enlightenment?

                          Zen's Struggle Against Dualism
To answer that, one needs to understand something about what enlightenment is. Perhaps
the most concise summary of enlightenment w be: transcending dualism. Now what is
dualism? Dualism is the conceptual division of the world into categories. Is it possible to
transcend this natural tendency? By prefixing the word "division" by the word
"conceptual", I may have made it seem that this is an intellectual or cons effort, and
perhaps thereby given the impression that dualism could overcome simply by suppressing
thought (as if to suppress thinking act were simple!). But the breaking of the world into
categories takes plat below the upper strata of thought; in fact, dualism is just as a
perceptual division of the world into categories as it is a conceptual division In other
words, human perception is by nature a dualistic phenomenon which makes the quest for
enlightenment an uphill struggle, to say the least.
        At the core of dualism, according to Zen, are words just plain w The use of words
is inherently dualistic, since each word represents, obviously, a conceptual category.
Therefore, a major part of Zen is the against reliance on words. To combat the use of
words, one of the devices is the koan, where words are so deeply abused that one's mi
practically left reeling, if one takes the koans seriously. Therefore perhaps wrong to say
that the enemy of enlightenment is logic; rather dualistic, verbal thinking. In fact, it is
even more basic than that: perception. As soon as you perceive an object, you draw a line
between it and the rest of the world; you divide the world, artificially, into parts you
thereby miss the Way.

Here is a koan which demonstrates the struggle against words:
Koan:

Shuzan held out his short staff and said: "If you call this a short staff, you oppose its
reality. If you do not call it a short staff, you ignore the fact. N, what do you wish to call
this?"




Mumon and Gödel                                                                           259

              FIGURE 49. Day and Night, by M. C. Escher (woodcut, 1938).

Mumon's Commentary:

If you call this a short staff, you oppose its reality. If you do not call it a short staff, you
ignore the fact. It cannot be expressed with words and it cannot be expressed without
words. Now say quickly what it is.

Mumon's Poem:
.
     Holding out the short staff,
     He gave an order of life or death.
     Positive and negative interwoven,
     Even Buddhas and patriarchs cannot escape this attack.

("Patriarchs" refers to six venerated founders of Zen Buddhism, of whom Bodhidharma is
the first, and Eno is the sixth.)
         Why is calling it a short staff opposing its reality? Probably because such a
categorization gives the appearance of capturing reality, whereas the surface has not even
been scratched by such a statement. It could be compared to saying "5 is a prime
number". There is so much more-an infinity of facts-that has been omitted. On the other
hand, not to call it a staff is, indeed, to ignore that particular fact, minuscule as it may be.
Thus words lead to some truth-some falsehood, perhaps, as well-but certainly not to all
truth. Relying on words to lead you to the truth is like relying on an incomplete formal
system to lead you to the truth. A formal system will give you some truths, but as we
shall soon see, a formal system-no matter how powerful-cannot lead to all truths. The
dilemma of mathematicians is: what else is there to rely on, but formal systems? And the
dilemma of

Mumon and Gödel                                                                            260

Zen people is, what else is there to rely on, but words? Mumon states t dilemma very
clearly: "It cannot be expressed with words and it cannot
expressed without words."

      Here is Nansen, once again:'

   Joshu asked the teacher Nansen, "What is the true Way?"
   Nansen answered, "Everyday way is the true Way.' Joshu asked, "Can I study it?"
   Nansen answered, "The more you study, the further from the Way." Joshu asked, "If I
   don't study it, how can I know it?"
   Nansen answered, "The Way does not belong to things seen: nor to thing: unseen. It
   does not belong to things known: nor to things unknown. Do not seek it, study it, or
   name it. To find yourself on it, open yourself wide as the sky." [See Fig. 50.]

FIGURE 50. Rind, by M. C. Escher (wood-engraving, 1955).




Mumon and Gödel                                                                    261

        This curious statement seems to abound with paradox. It is a little reminiscent of
this surefire cure for hiccups: "Run around the house three times without thinking of the
word `wolf'." Zen is a philosophy which seems to have embraced the notion that the road
to ultimate truth, like the only surefire cure for hiccups, may bristle with paradoxes.

                          Ism, The Un-Mode, and Unmon
If words are bad, and thinking is bad, what is good? Of course, to ask this is already
horribly dualistic, but we are making no pretense of being faithful to Zen in discussing
Zen-so we can try to answer the question seriously. I have a name for what Zen strives
for: ism. Ism is an antiphilosophy, a way of being without thinking. The masters of ism
are rocks, trees, clams; but it is the fate of higher animal species to have to strive for ism,
without ever being able to attain it fully. Still, one is occasionally granted glimpses of
ism. Perhaps the following koan offers such a glimpse :7

   Hyakujo wished to send a monk to open a new monastery. He told his pupils that
   whoever answered a question most ably would be appointed. Placing a water vase on
   the ground, he asked: "Who can say what this is without calling its name?"
     The chief monk said: "No one can call it a wooden shoe."
     Isan, the cooking monk, tipped over the vase with his foot and went out. Hyakujo
     smiled and said: "The chief monk loses." And Isan became the
   master of the new monastery.

To suppress perception, to suppress logical, verbal, dualistic thinking-this is the essence
of Zen, the essence of ism. This is the Unmode-not Intelligent, not Mechanical, just "Un".
Joshu was in the Unmode, and that is why his 'MU' unasks the question. The Un-mode
came naturally to Zen Master Unmon:8

   One day Unmon said to his disciples, "This staff of mine has transformed itself into a
   dragon and has swallowed up the universe! Oh, where are the rivers and mountains
   and the great earth?"

Zen is holism, carried to its logical extreme. If holism claims that things can only be
understood as wholes, not as sums of their parts, Zen goes one further, in maintaining that
the world cannot be broken into parts at all. To divide the world into parts is to be
deluded, and to miss enlightenment.

   A master was asked the question, "What is the Way?" by a curious monk. "
     It is right before your eyes," said the master. "Why do I not see it for myself?"
     "Because you are thinking of yourself."
     "What about you: do you see it?"
     "So long as you see double, saying `I don't', and `you do', and so on, your
   eyes are clouded," said the master.
     "When there is neither 'I' nor 'You', can one see it?"
     "When there is neither `I' nor `You', who is the one that wants to see it?"9


Mumon and Gödel                                                                            262

      Apparently the master wants to get across the idea that an enlighte state is one
where the borderlines between the self and the rest of universe are dissolved. This would
truly be the end of dualism, for a says, there is no system left which has any desire for
perception. But what is that state, if not death? How can a live human being dissolve the
borderlines between himself and the outside world?

                                  Zen and Tumbolia
The Zen monk Bassui wrote a letter to one of his disciples who was about to die, and in it
he said: "Your end which is endless is as a snowflake dissolving in the pure air." The
snowflake, which was once very much a discernible subsystem of the universe, now
dissolves into the larger system which 4 held it. Though it is no longer present as a
distinct subsystem, its essence somehow still present, and will remain so. It floats in
Tumbolia, along hiccups that are not being hiccupped and characters in stories that are
being read . . . That is how I understand Bassui's message.
        Zen recognizes its own limitations, just as mathematicians have lea: to recognize
the limitations of the axiomatic method as a method attaining truth. This does not mean
that Zen has an answer to what beyond Zen any more than mathematicians have a clear
understanding the forms of valid reasoning which lie outside of formalization. One o1
clearest Zen statements about the borderlines of Zen is given in the fol ing strange koan,
very much in the spirit of Nansen:10

   Tozan said to his monks, "You monks should know there is an even high
   understanding in Buddhism." A monk stepped forward and asked, "What the higher
   Buddhism?" Tozan answered, "It is not Buddha."

There is always further to go; enlightenment is not the end-all of And there is no recipe
which tells how to transcend Zen; the only thing can rely on for sure is that Buddha is not
the way. Zen is a system cannot be its own metasystem; there is always something
outside of which cannot be fully understood or described within Zen.

                                   Escher and Zen
In questioning perception and posing absurd answerless riddles, Zen company, in the
person of M. C. Escher. Consider Day and Night (Fig. 4 masterpiece of "positive and
negative interwoven" (in the words of Mumoni). One might ask, "Are those really birds,
or are they really field it really night, or day?" Yet we all know there is no point to such
questions The picture, like a Zen koan, is trying to break the mind of logic. Es4 also
delights in setting up contradictory pictures, such as Another World




Mumon and Gödel                                                                         263

                 FIGURE 51. Puddle, by M. C. Escher (woodcut, 1952).

(Fig. 48)-pictures that play with reality and unreality the same way as Zen plays with
reality and unreality. Should one take Escher seriously? Should one take Zen seriously?
        There is a delicate, haiku-like study of reflections in Dewdrop (Fig. 47); and then
there are two tranquil images of the moon reflected in still waters: Puddle (Fig. 51), and
Rippled Surface (Fig. 52). The reflected moon is a theme which recurs in various koans.
Here is an example:'

   Chiyono studied Zen for many years under Bukko of Engaku. Still, she could not
   attain the fruits of meditation. At last one moonlit night she was carrying water in an
   old wooden pail girded with bamboo. The bamboo broke, and the bottom fell out of
   the pail. At that moment, she was set free. Chiyono said, "No more water in the pail,
   no more moon in the water."

Three Worlds: an Escher picture (Fig. 46), and the subject of a Zen koan:12

   A monk asked Ganto, "When the three worlds threaten me, what shall I do?" Ganto
   answered, "Sit down." "I do not understand," said the monk. Canto said, "Pick up the
   mountain and bring it to me. Then I will tell you."




Mumon and Gödel                                                                        264

                               Hemiolia and Escher
In Verbum (Fig. 149), oppositions are made into unities on several I Going around we see
gradual transitions from black birds to white birds to black fish to white fish to black
frogs to white frogs to black birds ... six steps, back where we started! Is this a
reconciliation of the dichotomy of black and white? Or of the trichotomy of birds, fish,
and frogs? Or sixfold unity made from the opposition of the evenness of 2 an oddness of
3? In music, six notes of equal time value create a rhythmic ambiguity-are they 2 groups
of 3, or 3 groups of 2? This ambiguity has a name: hemiolia. Chopin was a master of
hemiolia: see his Waltz op. his Etude op. 25, no. 2. In Bach, there is the Tempo di
Menuetto from the keyboard Partita no. 5, or the incredible Finale of the first Sonata
unaccompanied violin, in G Minor.
        As one glides inward toward the center of Verbum, the distinctions gradually blur,
so that in the end there remains not three, not two, but one single essence: "VERBUM",
which glows with brilliancy-perhaps a symbol of enlightenment. Ironically, ` verbum"
not only is a word, but "word"-not exactly the most compatible notion with Zen. On the
hand, "verbum" is the only word in the picture. And Zen master 1 once said, "The
complete Tripitaka can be expressed in one character ("Tripitaka", meaning "three
baskets", refers to the complete texts c original Buddhist writings.) What kind of
decoding-mechanism, I wonder would it take to suck the three baskets out of one
character? Perhaps one with two hemispheres.


            FIGURE 52. Rippled Surface, by M. C. Escher (lino-cut, 1950).




Mumon and Gödel                                                                       265

           FIGURE 53. Three Spheres II, by M. C. Escher (lithograph, 1946),

                                       Indra's Net

Finally, consider Three Spheres II (Fig. 53), in which every part of the world seems to
contain, and be contained in, every other part: the writing table reflects the spheres on top
of it, the spheres reflect each other, as well as the writing table, the drawing of them, and
the artist drawing it. The endless connections which all things have to each other is only
hinted at here, yet the hint is enough. The Buddhist allegory of "Indra's Net" tells of an
endless net of threads throughout the universe, the horizontal threads running through
space, the vertical ones through time. At every crossing of threads is an individual, and
every individual is a crystal bead. The great light of "Absolute Being" illuminates and
penetrates every crystal bead; moreover, every crystal bead reflects not only the light
from every other crystal in the net-but also every reflection of every reflection throughout
the universe.
         To my mind, this brings forth an image of renormalized particles: in every
electron, there are virtual photons, positrons, neutrinos, muons ... ; in every photon, there
are virtual electrons, protons, neutrons, pions ... ; in every pion, there are ...
         But then another image rises: that of people, each one reflected in the minds of
many others, who in turn are mirrored in yet others, and so on.
         Both of these images could be represented in a concise, elegant way by using
Augmented Transition Networks. In the case of particles, there would be one network for
each category of particle; in the case of people,




Mumon and Gödel                                                                          266

one for each person. Each one would contain calls to many others, t creating a virtual
cloud of ATN's around each ATN. Calling one we create calls on others, and this process
might cascade arbitrarily far, un~ bottomed out.

                                  Mumon on MU
Let us conclude this brief excursion into Zen by returning to Mumon. H is his comment
on Joshu's MU:13

  To realize Zen one has to pass through the barrier of the patriarchs. Enlightenment
  always comes after the road of thinking is blocked. If you do nc pass the barrier of the
  patriarchs or if your thinking road is not blocked whatever you think, whatever you
  do, is like a tangling ghost. You may ask "What is a barrier of a patriarch?" This one
  word, 'MU', is it.
       This is the barrier of Zen. If you pass through it, you will see Joshu face t face.
  Then you can work hand in hand with the whole line of patriarchs. I this not a pleasant
  thing to do?
       If you want to pass this barrier, you must work through every bone in you body,
  through every pore of your skin, filled with this question: "What `MU'?" and carry it
  day and night. Do not believe it is the common negative symbol meaning nothing. It is
  not nothingness, the opposite of existence. I you really want to pass this barrier, you
  should feel like drinking a hot iro ball that you can neither swallow nor spit out.
       Then your previous lesser knowledge disappears. As a fruit ripening i season,
  your subjectivity and objectivity naturally become one. It is like dumb man who has
  had a dream. He knows about it but he cannot tell i
       When he enters this condition his ego-shell is crushed and he can shake th heaven
  and move the earth. He is like a great warrior with a sharp sword. If Buddha stands in
  his way, he will cut him down; if a patriarch offers him an obstacle, he will kill him;
  and he will be free in his way of birth and death. H can enter any world as if it were
  his own playground. I will tell you how to d this with this koan:
       Just concentrate your whole energy into this MU, and do not allow an
  discontinuation. When you enter this MU and there is no discontinuation -- your
  attainment will be as a candle burning and illuminating the who] universe.

                        From Mumon to the MU-puzzle
From the ethereal heights of Joshu's MU, we now descend to the private lowlinesses of
Hofstadter's MU . . . I know that you have already concentrated your whole energy into
this MU (when you read Chapter 1). So n wish to answer the question which was posed
there:

                            Has MU theorem-nature, or not?

The answer to this question is not an evasive MU; rather, it is a resounding NO. In order
to show this, we will take advantage of dualistic, logical thinking.


Mumon and Gödel                                                                       267

We made two crucial observations in Chapter I:

   (1) that the MU-puzzle has depth largely because it involves the interplay of
       lengthening and shortening rules;

   (2) that hope nevertheless exists for cracking the problem by employing a tool which
        is in some sense of adequate depth to handle matters of that complexity: the
        theory of numbers.

We did not analyze the MU-puzzle in those terms very carefully in Chapter I; we shall do
so now. And we will see how the second observation (when generalized beyond the
insignificant MIU-system) is one of the most fruitful realizations of all mathematics, and
how it changed mathematicians' view of their own discipline.
        For your ease of reference, here is a recapitulation of the MIU-system:

SYMBOLS: M, I, U

Axiom: MI

RULES:

   I. If xl is a theorem, so is xIU.
   II. If Mx is a theorem, so is Mxx.
   III. In any theorem, III can be replaced by U.
   IV. UU can be dropped from any theorem.

                Mumon Shows Us How to Solve the MU-puzzle
According to the observations above, then, the MU-puzzle is merely a puzzle about
natural numbers in typographical disguise. If we could only find a way to transfer it to the
domain of number theory, we might be able to solve it. Let us ponder the words of
Mumon, who said, "If any of you has one eye, he will see the failure on the teacher's
part." But why should it matter to have one eye?
         If you try counting the number of l's contained in theorems, you will soon notice
that it seems never to be 0. In other words, it seems that no matter how much lengthening
and shortening is involved, we can never work in such a way that all I's are eliminated.
Let us call the number of I's in any string the I-count of that string. Note that the I-count
of the axiom MI is 1. We can do more than show that the I-count can't be 0-we can show
that the I-count can never be any multiple of 3.
         To begin with, notice that rules I and IV leave the I-count totally undisturbed.
Therefore we need only think about rules II and III. As far as rule III is concerned, it
diminishes the I-count by exactly 3. After an application of this rule, the I-count of the
output might conceivably be a multiple of 3-but only if the I-count of the input was also.
Rule III, in short, never creates a multiple of 3 from scratch. It can only create one when
it began with one. The same holds for rule II, which doubles the


Mumon and Gödel                                                                          268

I-count. The reason is that if 3 divides 2n, then-because 3 does not dig 2-it must divide n
(a simple fact from the theory of numbers). Neither rule II nor rule III can create a
multiple of 3 from scratch.
       But this is the key to the MU-puzzle! Here is what we know:

   (1) The I-count begins at 1 (not a multiple of 3);

   (2) Two of the rules do not affect the I-count at all; (3)

   (3) The two remaining rules which do affect the I-count do so in such a way as never
   to create a multiple of 3 unless given one initially.

 The conclusion-and a typically hereditary one it is, too-is that I-count can never become
  any multiple of 3. In particular, 0 is a forbid value of the I-count. Hence, MU is not a
                               theorem of the MIU-system.
        Notice that, even as a puzzle about I-counts, this problem was plagued by the
crossfire of lengthening and shortening rules. Zero became the goal; I-counts could
increase (rule II), could decrease (rule III). 1 we analyzed the situation, we might have
thought that, with enough switching back and forth between the rules, we might
eventually hit 0. IS thanks to a simple number-theoretical argument, we know that the
impossible.

                       Gödel-Numbering the MIU-System
Not all problems of the the type which the MU-puzzle symbolizes at easy to solve as this
one. But we have seen that at least one such pr could be embedded within, and solved
within, number theory. We are going to see that there is a way to embed all problems
about any for system, in number theory. This can happen thanks to the discovery Gödel,
of a special kind of isomorphism. To illustrate it, I will use MIU-system.

We begin by considering the notation of the MIU-system. We map each symbol onto a
new symbol:

                                  M <= => 3
                                   I <= => 1
                                  U <= => 0

The correspondence was chosen arbitrarily; the only rhyme or reason is that each symbol
looks a little like the one it is mapped onto. I number is called the Gödel number of the
corresponding letter. Now I sure you can guess what the Gödel number of a multiletter
string will be:

                                  MU <= => 30
                                  MIIU <= => 3110
                                     Etc.


Mumon and Gödel                                                                        269

It is easy. Clearly this mapping between notations is an information preserving
transformation; it is like playing the same melody on two different instruments.

       Let us now take a look at a typical derivation in the MIU-system, written
simultaneously in both notations:

(1)    MI             axiom 31
(2)    MII            rule 2 311
(3)    MIIII          rule 2 31111
(4)    MUI            rule 3 301
(5)    MUIU           rule 1 3010
(6)    MUIUUIU        rule 2 3010010
(7)    MUIIU          rule 4 30110

The left-hand column is obtained by applying our four familiar typographical rules. The
right-hand column, too, could be thought of as having been generated by a similar set of
typographical rules. Yet the right-hand column has a dual nature. Let me explain what
this means.

          Seeing Things Both Typographically and Arithmetically
We could say of the fifth string ('3010') that it was made from the fourth, by appending a
`0' on the right; on the other hand we could equally well view the transition as caused by
an arithmetical operation-multiplication by 10, to be exact. When natural numbers are
written in the decimal system, multiplication by 10 and putting a `0' on the right are
indistinguishable operations. We can take advantage of this to write an arithmetical rule
which corresponds to typographical rule I:

ARITHMETICAL \rule la: A number whose decimal expansion ends on the right in `1'\par

can be multiplied by 10.

We can eliminate the reference to the symbols in the decimal expansion by arithmetically
describing the rightmost digit:

ARITHMETICAL \rule Ib: A number whose remainder when divided by 10 is 1, can\par

be multiplied by 10.

Now we could have stuck with a purely typographical rule, such as the following one:

TYPOGRAPHICAL \rule I: From any theorem whose rightmost symbol is ' 1' a new\par

theorem can be made, by appending `0' to the right of that 1'.
They would have the same effect. This is why the right-hand column has a "dual nature":
it can be viewed either as a series of typographical opera-




Mumon and Gödel                                                                        270

tions changing one pattern of symbols into another, or as a series arithmetical operations
changing one magnitude into another. But the are powerful reasons for being more
interested in the arithmetical version Stepping out of one purely typographical system
into another isomorphic typographical system is not a very exciting thing to do; whereas
stepping clear out of the typographical domain into an isomorphic part of number theory
has some kind of unexplored potential. It is as if somebody h known musical scores all
his life, but purely visually-and then, all o: sudden, someone introduced him to the
mapping between sounds a musical scores. What a rich, new world! Then again, it is as if
somebody h been familiar with string figures all his life, but purely as string figur devoid
of meaning-and then, all of a sudden, someone introduced him the mapping between
stories and strings. What a revelation! The discovery of Gödel-numbering has been
likened to the discovery, by Descartes, of t isomorphism between curves in a plane and
equations in two variables; incredibly simple, once you see it-and opening onto a vast
new world
        Before we jump to conclusions, though, perhaps you would like to a more
complete rendering of this higher level of the isomorphism. It i very good exercise. The
idea is to give an arithmetical rule whose action is indistinguishable from that of each
typographical rule of the MIU-system:

A solution is given below. In the rules, m and k are arbitrary natural numbers, and n is
any natural number which is less than 10m

\rule 1: If we have made 10m + 1, then we can make 10 x (10m + 1)\par

Example: Going from line 4 to line 5. Here, m = 30.

\rule 2: If we have made 3 x 10" + n, then we can make 10' X X (3 x 10"'+n)+n.\par

Example: Going from line 1 to line 2, where both m and n equal 1.

\rule 3: If we have made k x 10 "`+ 111 x 10'+n, then we can make k x 10"+` + n.\par

Example: Going from line 3 to line 4. Here, m and n are 1, and k is 3.

\rule 4: If we have made k x 10rn+z + n, k x 10" +n. then we can make k x 10m + n\par

Example: Going from line 6 to line 7. Here, m = 2, n = 10, and k = 301.

Let us not forget our axiom! Without it we can go nowhere. Therefore, let us postulate
that:

                                     We can make 31.

Now the right-hand column can be seen as a full-fledged arithmetic process, in a new
arithmetical system which we might call the 310-system




Mumon and Gödel                                                                         271

I
(1)      31             given
(2)      311            rule 2 (m=1, n=1)
(3)      31111          rule 2 (m=2, n=11)
(4)      301            rule 3 (m=1, n=1, k=3)
(5)      3010           rule 1 (m=30)
(6)      3010010        rule 2 (m=3, n=10)
(7)      30110          rule 4 (m=2, n=10, k=301)

Notice once again that the lengthening and shortening rules are ever with us in this "310-
system"; they have merely been transposed into the domain of numbers, so that the Godel
numbers go up and down. If you look carefully at what is going on, you will discover that
the rules are based on nothing more profound than the idea that shifting digits to left and
right in decimal representations of integers is related to multiplications and divisions by
powers of 10. This simple observation finds its generalization in the following

      CENTRAL PROPOSITION: If there is a typographical rule which tells how
      certain digits are to be shifted, changed, dropped, or inserted in any number
      represented decimally, then this rule can be represented equally well by an
      arithmetical counterpart which involves arithmetical operations with powers of 10
      as well as additions, subtractions, and so forth.

                                       More briefly:

      Typographical rules for manipulating numerals are actually arithmetical rules for
      operating on numbers.

This simple observation is at the heart of Gödel’s method, and it will have an absolutely
shattering effect. It tells us that once we have a Gödel numbering for any formal system,
we can straightaway form a set of arithmetical rules which complete the Gödel
isomorphism. The upshot is that we can transfer the study of any formal system-in fact
the study of all formal systems-into number theory.

                              MIU-Producible Numbers
Just as any set of typographical rules generates a set of theorems, a corresponding set of
natural numbers will be generated by repeated applications of arithmetical rules. These
producible numbers play the same role inside number theory as theorems do inside any
formal system. Of course, different numbers will be producible, depending on which
rules are adopted. "Producible numbers" are only producible relative to a system of
arithmetical rules. For example, such numbers as 31, 3010010, 3111, and so forth could
be called MIU-producible numbers-an ungainly name, which might be shortened to
MIU-numbers, symbolizing the fact that those numbers are the ones that result when you
transcribe the MIU-system into number theory, via Gödel-numbering. If we were to
Gödel-number the pq-system


Mumon and Gödel                                                                        272

and then "arithmetize" its rules, we could call the producible numbers "pq-numbers"-and
so on.
       Note that the producible numbers (in any given system) are defined by a recursive
method: given numbers which are known to be producible, we have rules telling how to
make more producible numbers. Thus, the class of numbers known to be producible is
constantly extending itself, in much the same way that the list of Fibonacci numbers, or
Q-numbers, does. The set of producible numbers of any system is a recursively
enumerable set. What about its complement-the set of nonproducible numbers? Is that set
always recursively enumerable? Do numbers which are nonproducible share some
common arithmetical feature?
       This is the sort of issue which arises when you transpose the study of formal
systems into number theory. For each system which is arithmetized, one can ask, "Can
we characterize producible numbers in a simple way?" "Can we characterize
nonproducible numbers in a recursively enumerable way?" These are difficult questions
of number theory. Depending on the system which has been arithmetized, such questions
might prove too hard for us to resolve. But if there is any hope for solving such problems,
it would have to reside in the usual kind of step-by-step reasoning as it applies to natural
numbers. And that, of course, was put in its quintessential form in the previous Chapter.
TNT seemed, to all appearances, to have captured all valid mathematical thinking
processes in one single, compact system.

              Answering Questions about Producible Numbers
                          by Consulting TNT

Could it be, therefore, that the means with which to answer any question about any
formal system lies within just a single formal system-TNT? It seems plausible. Take, for
instance, this question:

                          Is MU a theorem of the MIU-system?

Finding the answer is equivalent to determining whether 30 is a MIU number or not.
Because it is a statement of number theory, we should expect that, with some hard work,
we could figure out how to translate the sentence "30 is a MIU-number" into TNT-
notation, in somewhat the same way as we figured out how to translate other number-
theoretical sentences into TNT-notation. I should immediately caution the reader that
such a translation, though it does exist, is immensely complex. If you recall, I pointed out
in Chapter VIII that even such a simple arithmetical predicate as "b is a power of 10" is
very tricky to code into TNT-notation-and the predicate "b is a MIU-number" is a lot
more complicated than that! Still, it can be found; and the numeral
SSSSSSSSSSSSSSSSSSSSSSSSSSSSSSO can be substituted for every b. This will
result in a MONstrous string of TNT, a string of TNT which speaks about the MU-
puzzle. Let us therefore call that string "MUMON". Through MUMON and strings like
it, TNT is now capable of speaking "in code" about the MIU-system.




Mumon and Gödel                                                                         273

                           The Dual Nature of MUMON

In order to gain some benefit from this peculiar transformation of the original question,
we would have to seek the answer to a new question:

                            Is MUMON a theorem of TNT?

All we have done is replace one relatively short string (MU) by another (the monstrous
MUMON), and a simple formal system (the MIU-system) by a complicated one (TNT).
It isn't likely that the answer will be any more forthcoming even though the question has
been reshaped. In fact, TNT has a full complement of both lengthening and shortening
rules, and the reformulation of the question is likely to be far harder than the original.
One might even say that looking at MU via MUMON is an intentionally idiotic way of
doing things. However, MUMON can be looked at on more than one level.
          In fact, this is an intriguing point: MUMON has two different passive meanings.
Firstly, it has the one which was given before:

                                  30 is a MIU-number.

But secondly, we know that this statement is tied (via isomorphism) to the statement

                          MU is a theorem of the MIU-system.

So we can legitimately quote this latter as the second passive meaning of MUMON. It
may seem very strange because, after all, MUMON contains nothing but plus signs,
parentheses, and so forth-symbols of TNT. How can it possibly express any statement
with other than arithmetical content?
        The fact is, it can. Just as a single musical line may serve as both harmony and
melody in a single piece; just as "BACH" may be interpreted as both a name and a
melody; just as a single sentence may be an accurate structural description of a picture by
Escher, of a section of DNA, of a piece by Bach, and of the dialogue in which the
sentence is embedded, so MUMON can be taken in (at least) two entirely different ways.
This state of affairs comes about because of two facts:

   Fact 1. Statements such as "MU is a theorem" can be coded into number theory
        via Gödel’s isomorphism.

   Fact 2. Statements of number theory can be translated into TNT.

It could be said that MUMON is, by Fact 1, a coded message, where the symbols of the
code are, by Fact 2, just symbols of TNT.




Mumon and Gödel                                                                        274

                           Codes and Implicit Meaning
Now it could be objected here that a coded message, unlike an uncod message, does not
express anything on its own-it requires knowledge the code. But in reality there is no
such thing as an uncoded message. There are only messages written in more familiar
codes, and message written in less familiar codes. If the meaning of a message is to be
revealed it must be pulled out of the code by some sort of mechanism, or isomorphism. It
may be difficult to discover the method by which the decoding should be done; but once
that method has been discovered, the message becomes transparent as water. When a
code is familiar enough, it ceases appearing like a code; one forgets that there is a
decoding mechanism. The message is identified with its meaning.
         Here we have a case where the identification of message and meant is so strong
that it is hard for us to conceive of an alternate meaning: residing in the same symbols.
Namely, we are so prejudiced by the symbols of TNT towards seeing number-theoretical
meaning (and only numb( theoretical meaning) in strings of TNT, that to conceive of
certain string of TNT as statements about the MIU-system is quite difficult. But Gödel’s
isomorphism compels us to recognize this second level of meaning certain strings of
TNT.

       Decoded in the more familiar way, MUMON bears the message:

                                  30 is a MIU-number.

This is a statement of number theory, gotten by interpreting each sign the conventional
way.
        But in discovering Gödel-numbering and the whole isomorphism bu upon it, we
have in a sense broken a code in which messages about the MIU-system are written in
strings of TNT. Gödel’s isomorphism is a n( information-revealer, just as the
decipherments of ancient scripts we information-revealers. Decoded by this new and less
familiar mechanism MUMON bears the message

                          MU is a theorem of the MIU-system.

The moral of the story is one we have heard before: that meaning is ; automatic by-
product of our recognition of any isomorphism; therefore there are at least two passive
meanings of MUMON-maybe more!

                   The Boomerang: Gödel-Numbering TNT
Of course things do not stop here. We have only begun realizing the: potential of Gödel’s
isomorphism. The natural trick would be to turn TNT's capability of mirroring other
formal systems back on itself, as the Tortoise turned the Crab's phonographs against
themselves, and as his Goblet G turned against itself, in destroying itself. In order to do
this, we



Mumon and Gödel                                                                        275

will have to Gödel-number TNT itself, just as we did the MIU-system, and then
"arithmetize" its rules of inference. The Gödel-numbering is easy to do. For instance, we
could make the following correspondence:
Symbol...         Codon         Mnemonic Justification
0                 666           Number of the Beast for the Mysterious Zero
S       .....     123           successorship: 1, 2, 3,
=       .....     111           visual resemblance, turned sideways
+       ....      112           1+1=2
.       ....      236           2x3=6
(       ....      362           ends in 2        *
)       ....      323           ends in 3        *
<                 212           ends in 2        * these three pairs
>       ....      213           ends in 3        * form a pattern
[       ....      312           ends in 2        *
]       ....      313           ends in 3        *
α       ....      262           opposite to V (626)
´....             163           163 is prime
∧       .....     161           ´∧´ is a "graph" of the sequence 1-6-1
∨       .....     616           ´∨' is a "graph" of the sequence 6-1-6
⊃       .....     633           ` 6 "implies" 3 and 3, in some sense .
~       .....     223           . 2 + 2 is not 3
ℑ       .....     333           `ℑ' looks like `3'
V       .....     626           opposite to a; also a "graph" of 6-2-6
.:      .....     636           two dots, two sixes
punc. .....       611           special number, as on Bell system (411, 911)
       Each symbol of TNT is matched up with a triplet composed of the digits 1, 2, 3,
and 6, in a manner chosen for mnemonic value. I shall call each such triplet of digits a
Gödel codon, or codon for short. Notice that I have given no codon for b, c, d, or e; we
are using austere TNT. There is a hidden motivation for this, which you will find out
about in Chapter XVI. I will explain the bottom entry, "punctuation", in Chapter XIV.
       Now we can rewrite any string or rule of TNT in the new garb. Here, for instance,
is Axiom 1 in the two notations, the old below the new:
                626,262,636,223,123,262,111,666
                 V a      : ~     S a = 0
Conveniently, the standard convention of putting in a comma every third digit happens to
coincide with our colons, setting them off for "easy" legibility.
        Here is the Rule of Detachment, in the new notation:
RULE: If x and 212x6331213 are both theorems, then 1 is a theorem. Finally, here is an
entire derivation taken from last Chapter, given in austere TNT and also transcribed into
the new notation:

Mumon and Gödel                                                                      276

626,262,636,626262,163,636,362262,112,123,262,163,323,111,123,362,262,112,262,163,323 axiom:
 V α :: V α ' : ( α + 5 α ' ) = S ( α + α ' )
626,262.163,636,362,123,666,112,123,262,163,323,111,123,362,123,666,112,262,163,32 specification
 V α ' : ( S 0 + S α ´ ) = S (                                 S 0 + α ´ )
362,123,666,112,123,666,323,1 11,123,362,123,666,112,666,323                     specification
 (   S 0 + 5          0 ) = S (           S 0      + 0     )
626,262,636,362 262,112,666, 3 23,111,262                                        axiom
 V• α : ( α + 0 ) = α
362,123,666,112,666,323,111,123,666                                              specification
  ( S 0 + 0           ) =     S 0
123,362,123,666,112,666,323,11 1,123,123,666                                     insert '12;
 S    ( S 0 + 0           ) =      S S 0
362,123,666,112,123,666,323,111,123,123,666                                      transitivity
  ( S 0 + 5 0            ) =       S S 0

Notice that I changed the name of the "Add S" rule to "Insert `123' ", since that is the
typographical operation which it now legitimizes.
        This new notation has a pretty strange feel to it. You lose all sense o meaning; but
if you had been brought up on it, you could read strings it this notation as easily as you do
TNT. You would be able to look and, at glance, distinguish well-formed formulas from
ill-formed ones. Naturally since it is so visual, you would think of this as a typographical
operation but at the same time, picking out well-formed formulas in this notation i
picking out a special class of integers, which have an arithmetical definition too.
        Now what about "arithmetizing" all the rules of inference? As matter stand, they
are all still typographical rules. But wait! According to the Central Proposition, a
typographical rule is really equivalent to al arithmetical rule. Inserting and moving digits
in decimally represented numbers is an arithmetical operation, which can be carried out
typographically. Just as appending a 'O' on the end is exactly the same as multiplying b,
10, so each rule is a condensed way of describing a messy arithmetical operation.
Therefore, in a sense, we do not even need to look for equivalent arithmetical rules,
because all of the rules are already arithmetical!

         TNT-Numbers: A Recursively Enumerable Set of Numbers
        Looked at this way, the preceding derivation of the theorem
"362,123,666,112,123,666,323,111,123,123,666" is a sequence of high] convoluted
number-theoretical transformations, each of which acts on one or more input numbers,
and yields an output number, which is, as before, called a producible number, or, to be
more specific, a TNT-number. Some the arithmetical rules take an old TNT-number and
increase it in a particular way, to yield a new TNT-number; some take an old TNT-
number a and decrease it; other rules take two TNT-numbers, operate on each of them
some odd way, and then combine the results into a new TNT-number
and so on and so forth. And instead of starting with just one know: 'TNT-number, we
have five initial TNT-numbers-one for each (austere axiom, of course. Arithmetized TNT
is actually extremely similar to the




Mumon and Gödel                                                                               277

arithmetized MIU-system, only there are more rules and axioms, and to write out
arithmetical equivalents explicitly would be a big bother-and quite unenlightening,
incidentally. If you followed how it was done for the MIU-system, there ought to be no
doubt on your part that it is quite analogous here.
       There is a new number-theoretical predicate brought into being by this
"Godelization" of TNT: the predicate

                                   α is a TNT-number.

For     example,      we     know       from   the     preceding      derivation      that
362,123,666,112,123,666,323,111,123,123,666 is a TNT-number, while on the other
hand, presumably 123,666,111,666 is not a TNT-number.
        Now it occurs to us that this new number-theoretical! predicate is expressible by
some string of TNT with one free variable, say a. We could put a tilde in front, and that
string would express the complementary notion

                                 α is not a TNT-number.

Now if we replaced all the occurrences of a in this second string by the TNT-numeral for
123,666,111,666-a numeral which would contain exactly 123,666,111,666 S's, much too
long to write out-we would have a TNT-string which, just like MUMON, is capable of
being interpreted on two levels. In the first place, that string would say

                         123,666,111,666 is not a TNT-number.

But because of the isomorphism which links TNT-numbers to theorems of TNT, there
would be a second-level meaning of this string, which is:

                             S0=0 is not a theorem of TNT.

                              TNT Tries to Swallow Itself

This unexpected double-entendre demonstrates that TNT contains strings which talk
about other strings of TNT. In other words, the metalanguage in which we, on the
outside, can speak about TNT, is at least partially imitated inside TNT itself. And this is
not an accidental feature of TNT; it happens because the architecture of any formal
system can be mirrored inside N (number theory). It is just as inevitable a feature of TNT
as are the vibrations induced in a record player when it plays a record. It seems as if
vibrations should come from the outside world-for instance, from jumping children or
bouncing balls; but a side effect of producing sounds-and an unavoidable one-is that they
wrap around and shake the very mechanism which produces them. It is no accident; it is a
side effect which cannot be helped. It is in the nature of record players. And it is in the
nature of any formalization of number theory that its metalanguage is embedded within it.




Mumon and Gödel                                                                        278

       We can dignify this observation by calling it the Central Dogma of MIathematical
Logic, and depicting it in a two-step diagram:

                               TNT => N => meta-'TNT

In words: a string of TNT has an interpretation in N; and a statement o may have a
second meaning as a statement about TNT.

                      G: A String Which Talks about Itself in Code

This much is intriguing yet it is only half the story. The rest of the st involves an
intensification of the self-reference. We are now at the st where the Tortoise was when he
realized that a record could be m; which would make the phonograph playing it break-but
now the quest is: "Given a record player, how do you actually figure out what to put the
record?" That is a tricky matter.
        We want to find a string of TNT-which we'll call 'G'-which is ab itself, in the
sense that one of its passive meanings is a sentence about G. particular the passive
meaning will turn out to be

                              "G is not a theorem of TNT."

I should quickly add that G also has a passive meaning which is a statement of number
theory; just like MUMON it is susceptible to being construed in least) two different
ways. The important thing is that each passive mean is valid and useful and doesn't cast
doubt on the other passive meaning in any way. (The fact that a phonograph playing a
record can induce vibrations in itself and in the record does not diminish in any way the
fact t those vibrations are musical sounds!)

                G's Existence Is What Causes TNT's Incompleteness

The ingenious method of creating G, and some important concepts relating to TNT, will
be developed in Chapters XIII and XIV; for now it is interesting to glance ahead, a bit
superficially, at the consequences finding a self-referential piece of TNT. Who knows? It
might blow up! In a sense it does. We focus down on the obvious question:

                             Is G a theorem of TNT, or not?

Let us be sure to form our own opinion on this matter, rather than rely G's opinion about
itself. After all, G may not understand itself any be than a Zen master understands
himself. Like MUMON, G may express a falsity. Like MU, G may be a nontheorem. We
don't need to believe every possible string of TNT-only its theorems. Now let us use our
power of reasoning to clarify the issue as best we can at this point.
        We will make our usual assumption: that TNT incorporates valid




Mumon and Gödel                                                                       279

methods of reasoning, and therefore that TNT never has falsities for theorems. In other
words, anything which is a theorem of TNT expresses a truth. So if G were a theorem, it
would express a truth, namely: "G is not a theorem". The full force of its self-reference
hits us. By being a theorem, G would have to be a falsity. Relying on our assumption that
TNT never has falsities for theorems, we'd be forced to conclude that G is not a theorem.
This is all right; it leaves us, however, with a lesser problem. Knowing that G is not a
theorem, we'd have to concede that G expresses a truth. Here is a situation in which TNT
doesn't live up to our expectations-we have found a string which expresses a true
statement yet the string is not a theorem. And in our amazement, we shouldn't lose track
of the fact that G has an arithmetical interpretation, too-which allows us to summarize
our findings this way:

   A string of TNT has been found; it expresses, unambiguously, a statement about
   certain arithmetical properties of natural numbers; moreover, by reasoning outside
   the system we can determine not only that the statement is a true one, but also that
   the string fails to be a theorem of TNT. And thus, if we ask TNT whether the
   statement is true, TNT says neither yes nor no.

       Is the Tortoise's string in the Mu Offering the analogue of G? Not quite. The
analogue of the Tortoise's string is ~G. Why is this so? Well, let us think a moment about
what -G says. It must say the opposite of what G says. G says, "G is not a theorem of
TNT", so ~G must say "G is a theorem". We could rephrase both G and ~G this way:

       G: "I am not a theorem (of TNT)."
       ~G: "My negation is a theorem (of TNT)."

It is ~G which is parallel to the Tortoise's string, for that string spoke not about itself, but
about the string which the Tortoise first proffered to Achilles -- which had an extra knot
on it (or one too few, however you want to look at it).

                             Mumon Has the Last Word
Mumon penetrated into the Mystery of the Undecidable anyone, in his concise poem on
Joshu's MU:

       Has a dog Buddha-nature?
       This is the most serious question of all.
       If you say yes or no,
       You lose your own Buddha-nature.




Mumon and Gödel                                                                            280

                                     Prelude . .
   Achilles and the Tortoise have come to the residence of their friend the Crab, to
   make the acquaintance of one of his friends, the Anteater. The introductions
   having been made, the four of them settle down to tea.

Tortoise We have brought along a little something for you, Mr. Crab. Crab: That's most
   kind of you. But you shouldn't have.
Tortoise: Just a token of our esteem. Achilles, would you like to give it to Mr. C?
Achilles: Surely. Best wishes, Mr. Crab. I hope you enjoy it.

   (Achilles hands the Crab an elegantly wrapped present, square and very thin. The
   Crab begins unwrapping it.)

Anteater: I wonder what it could be.
Crab: We'll soon find out. (Completes the unwrapping, and pulls out the gif)t Two
   records! How exciting! But there's no label. Uh-oh-is this another of your "specials",
   Mr. T?
Tortoise: If you mean a phonograph-breaker, not this time. But it is in fact a custom-
   recorded item, the only one of its kind in the entire world. In fact, it's never even been
   heard before-except, of course, when Bach played it.
Crab: When Bach played it? What do you mean, exactly?
Achilles: Oh, you are going to be fabulously excited, Mr. Crab, when Mr. T tells you
   what these records in fact are.
Tortoise: Oh, you go ahead and tell him, Achilles.
Achilles: May I? Oh, boy! I'd better consult my notes, then. (Pulls out a small filing card,
   and clears his voice.) Ahem. Would you be interested in hearing about the remarkable
   new result in mathematics, to which your records owe their existence?
Crab: My records derive from some piece of mathematics? How curious Well, now that
   you've provoked my interest, I must hear about it.
Achilles: Very well, then. (Pauses for a moment to sip his tea, then resumes) Have you
   heard of Fermat's infamous "Last Theorem"?
Anteater: I'm not sure ... It sounds strangely familiar, and yet I can't qui place it.
Achilles: It's a very simple idea. Pierre de Fermat, a lawyer by vocation b mathematician
   by avocation, had been reading in his copy of the class text Arithmetica by
   Diophantus, and came across a page containing the equation

                                         a2+b2=c2




Prelude                                                                                  281

          FIGURE 54. Mobius Strip II, by M. C. Escher (woodcut, 1963).




Prelude                                                                  282

He immediately realized that this equation has infinitely many solutions a, b, c, and then
   wrote in the margin the following notorious comment:
   The equation

                                       an +bn=cn

   has solutions in positive integers a, b, c, and n only when n = 2 (an then there are
   infinitely many triplets a, b, c which satisfy the equation); but there are no
   solutions for n > 2. I have discovered a truly marvelous proof of this statement,
   which, unfortunately, this margin is too small to contain.

Ever since that day, some three hundred years ago, mathematicians have been vainly
   trying to do one of two things: either to I Fermat's claim, and thereby vindicate
   Fermat's reputation, whit though very high, has been somewhat tarnished by skeptics
   who he never really found the proof he claimed to have found-or e: refute the claim,
   by finding a counterexample: a set of four integers a, b, c, and n, with n > 2, which
   satisfy the equation. Until recently, every attempt in either direction had met with
   failure. 1 sure, the Theorem has been proven for many specific values of i particular,
   all n up to 125,000.
Anteater: Shouldn't it be called a "Conjecture" rather than a "Theorem it's never been
   given a proper proof?
Achilles: Strictly speaking, you're right, but tradition has kept it this i
Crab: Has someone at last managed to resolve this celebrated quest Achilles: Indeed! In
   fact, Mr. Tortoise has done so, and as usual, by a wizardly stroke. He has not only
   found a PROOF of Fermat's Theorem (thus justifying its name as well as vindicating
   Fermat; also a COUNTEREXAMPLE, thus showing that the skeptics had good
   intuition!
Crab: Oh my gracious! That is a revolutionary discovery.
Anteater: But please don't leave us in suspense. What magical integer they, that satisfy
   Fermat's equation? I'm especially curious about the value of n.
Achilles: Oh, horrors! I'm most embarrassed! Can you believe this? the values at home on
   a truly colossal piece of paper. Unfortunately was too huge to bring along. I wish I
   had them here to show to y( it's of any help to you, I do remember one thing-the value
   of n only positive integer which does not occur anywhere in the continued fraction for
   π
Crab: Oh, what a shame that you don't have them here. But there reason to doubt what
   you have told us.




Prelude                                                                               283

                                                   FIGURE 55. Pierre de Fermat.

Anteater: Anyway, who needs to see n written out decimally? Achilles has just told us
   how to find it. Well, Mr. T, please accept my hearty felicitations, on the occasion of
   your epoch-making discovery!
Tortoise: Thank you. But what I feel is more important than the result itself is the
   practical use to which my result immediately led.
Crab: I am dying to hear about it, since I always thought number theory was the Queen of
   Mathematics -- the purest branch of mathematic -- the one branch of mathematics
   which has No applications!
Tortoise: You're not the only one with that belief, but in fact it is quite impossible to
   make a blanket statement about when or how some branch-or even some individual
   Theorem-of pure mathematics will have important repercussions outside of
   mathematics. It is quite unpredictable-and this case is a perfect example of that
   phenomenon.
Achilles: Mr. Tortoise's double-barreled result has created a breakthrough in the field of
   acoustico-retrieval!
Anteater: What is acoustico-retrieval?
Achilles: The name tells it all: it is the retrieval of acoustic information from extremely
   complex sources. A typical task of acoustico-retrieval is to reconstruct the sound
   which a rock made on plummeting into a lake from the ripples which spread out over
   the lake's surface.
Crab: Why, that sounds next to impossible!
Achilles: Not so. It is actually quite similar to what one's brain does, when it reconstructs
   the sound made in the vocal cords of another person from the vibrations transmitted
   by the eardrum to the fibers in the cochlea.
Crab: I see. But I still don't see where number theory enters the picture, or what this all
   has to do with my new records.




Prelude                                                                                  284

Achilles: Well, in the mathematics of acoustico-retrieval, there arise rr questions which
   have to do with the number of solutions of cer Diophantine equations. Now Mr. T has
   been for years trying to fit way of reconstructing the sounds of Bach playing his
   harpsichord, which took place over two hundred years ago, from calculations in\% ing
   the motions of all the molecules in the atmosphere at the pre time.
Anteater: Surely that is impossible! They are irretrievably gone, g forever!
Achilles: Thus think the nave ... But Mr. T has devoted many year this problem, and
   came to the realization that the whole thing hinged on the number of solutions to the
   equation

                                        an +bn=cn

   in positive integers, with n > 2.
Tortoise: I could explain, of course, just how this equation arises, but I’m sure it would
   bore you.
Achilles: It turned out that acoustico-retrieval theory predicts that Bach sounds can be
   retrieved from the motion of all the molecule the atmosphere, provided that EITHER
   there exists at least one solution to the equation
Crab: Amazing!
Anteater: Fantastic!
Tortoise: Who would have thought!
Achilles: I was about to say, "provided that there exists EITHER such a solution OR a
   proof that there are tic) solutions!" And therefore, Mr. T, in careful fashion, set about
   working at both ends of the problem, simultaneously. As it turns out, the discovery of
   the counterexample was the key ingredient to finding the proof, so the one led directly
   to the other.
Crab: How could that be?
Tortoise: Well, you see, I had shown that the structural layout of any pr of Fermat's Last
   Theorem-if one existed-could be described by elegant formula, which, it so happened,
   depended on the values ( solution to a certain equation. When I found this second
   equation my surprise it turned out to be the Fermat equation. An amusing accidental
   relationship between form and content. So when I found the counterexample, all I
   needed to do was to use those numbers blueprint for constructing my proof that there
   were no solutions to equation. Remarkably simple, when you think about it. I can't
   imagine why no one had ever found the result before.
Achilles: As a result of this unanticipatedly rich mathematical success, Mr. T was able to
   carry out the acoustico-retrieval which he had long dreamed of. And Mr. Crab's
   present here represents a palpable realization of all this abstract work.




Prelude                                                                                 285

Crab: Don't tell me it's a recording of Bach playing his own works for harpsichord!
Achilles: I'm sorry, but I have to, for that is indeed just what it is! This is a set of two
   records of Johann Sebastian Bach playing all of his Well Tempered Clavier. Each
   record contains one of the two volumes of the Well-Tempered Clavier; that is to say,
   each record contains 24 preludes and fugues-one in each major and minor key.
Crab: Well, we must absolutely put one of these priceless records on, immediately! And
   how can I ever thank the two of you?
Tortoise: You have already thanked us plentifully, with this delicious tea which you have
   prepared.

   (The Crab slides one of the records out of its jacket, and puts it on. The sound of
   an incredibly masterful harpsichordist fills the room, in the highest imaginable
   fidelity. One even hears-or is it one's imagination?-the soft sounds of Bach singing
   to himself as he plays ...)

Crab: Would any of you like to follow along in the score? I happen to have a unique
   edition of the Well-Tempered Clavier, specially illuminated by a teacher of mine who
   happens also to be an unusually fine calligrapher. Tortoise: I would very much enjoy
   that.

   (The Crab goes to his elegant glass-enclosed wooden bookcase, opens the doors, and
   draws out two large volumes.)

Crab: Here you are, Mr. Tortoise. I've never really gotten to know all the beautiful
   illustrations in this edition. Perhaps your gift will provide the needed impetus for me
   to do so.
Tortoise: I do hope so.
Anteater: Have you ever noticed how in these pieces the prelude always sets the mood
   perfectly for the following fugue?
Crab: Yes. Although it may be hard to put it into words, there is always some subtle
   relation between the two. Even if the prelude and fugue do not have a common
   melodic subject, there is nevertheless always some intangible abstract quality which
   underlies both of them, binding them together very strongly.
Tortoise: And there is something very dramatic about the few moments of silent suspense
   hanging between prelude and fugue-that moment where the the theme of the fugue is
   about to ring out, in single tones, and then to join with itself in ever-increasingly
   complex levels of weird, exquisite harmony.
Achilles: I know just what you mean. There are so many preludes and fugues which I
   haven't yet gotten to know, and for me that fleeting interlude of silence is very
   exciting; it's a time when I try to second-guess old Bach. For example, I always
   wonder what the fugue's tempo will be: allegro, or adagio? Will it be in 6/8, or 4/4?
   Will it have three voices, or five-or four? And then, the first voice starts ... Such an
   exquisite moment.




Prelude                                                                                 286

Crab: Ah, yes, well do I remember those long-gone days of my youth, days when I
   thrilled to each new prelude and fugue, filled with excitement of their novelty and
   beauty and the many unexpected' surprises which they conceal.
Achilles: And now? Is that thrill all gone?
Crab: It's been supplanted by familiarity, as thrills always will be. But that familiarity
   there is also a kind of depth, which has its own compensations. For instance, I find
   that there are always new surprises whit hadn't noticed before.
Achilles: Occurrences of the theme which you had overlooked?
Crab: Perhaps-especially when it is inverted and hidden among several other voices, or
   where it seems to come rushing up from the dept out of nowhere. But there are also
   amazing modulations which ii marvelous to listen to over and over again, and wonder
   how old B2 dreamt them up.
Achilles: I am very glad to hear that there is something to look forward after I have been
   through the first flush of infatuation with the Well Tempered Clavier-although it also
   makes me sad that this stage cot not last forever and ever.
Crab: Oh, you needn't fear that your infatuation will totally die. One the nice things about
   that sort of youthful thrill is that it can always resuscitated, just when you thought it
   was finally dead. It just takes the right kind of triggering from the outside.
Achilles: Oh, really? Such as what?
Crab: Such as hearing it through the ears, so to speak, of someone whom it is a totally
   new experience-someone such as you, Achilles. Somehow the excitement transmits
   itself, and I can feel thrilled again.
Achilles: That is intriguing. The thrill has remained dormant somewhere inside you, but
   by yourself, you aren't able to fish it up out of your subconscious.
Crab: Exactly. The potential of reliving the thrill is "coded", in sot unknown way, in the
   structure of my brain, but I don't have the power to summon it up at will; I have to
   wait for chance circumstance trigger it.
Achilles: I have a question about fugues which I feel a little embarrass about asking, but
   as I am just a novice at fugue-listening, I was wondering if perhaps one of you
   seasoned fugue-listeners might help me learning .. .
Tortoise: I'd certainly like to offer my own meager knowledge, if it might prove of' some
   assistance.
Achilles: Oh, thank you. Let me come at the question from an angle. Are you familiar
   with the print called Cube with Magic Ribbons, by M. Escher?
Tortoise: In which there are circular bands having bubble-like distortions which, as soon
   as you've decided that they are bumps, seem to turn it dents-and vice versa?




Prelude                                                                                 287

      FIGURE 56. Cube with Magic Ribbons, by M. C. Escher (lithograph, 1957).

Achilles: Exactly.
Crab: I remember that picture. Those little bubbles always seem to flip back and forth
   between being concave and convex, depending on the direction that you approach
   them from. There's no way to see them simultaneously as concave AND convex-
   somehow one's brain doesn't allow that. There are two mutually exclusive "modes" in
   which one can perceive the bubbles.
Achilles: Just so. Well, I seem to have discovered two somewhat analogous modes in
   which I can listen to a fugue. The modes are these: either to follow one individual
   voice at a time, or to listen to the total effect of all of them together, without trying to
   disentangle one from another. I have tried out both of these modes, and, much to my
   frustration, each one of them shuts out the other. It's simply not in my power to follow
   the paths of individual voices and at the same time to hear the whole effect. I find that
   I flip back and forth between one mode and the other, more or less spontaneously and
   involuntarily.




Prelude                                                                                    288

Anteater: Just as when you look at the magic bands, eh?
Achilles: Yes. I was just wondering ... does my description of they modes of fugue-
   listening brand me unmistakably as a naive, inexperienced listener, who couldn't even
   begin to grasp the deeper mo, perception which exist beyond his ken?
Tortoise: No, not at all, Achilles. I can only speak for myself, but I to myself shifting
   back and forth from one mode to the other without exerting any conscious control
   over which mode should he dominant don't know if our other companions here have
   also experience( thing similar.
Crab: Most definitely. It's quite a tantalizing phenomenon, since you feel that the essence
   of the fugue is flitting about you, and you can't grasp all of it, because you can't quite
   make yourself function ways at once.
Anteater: Fugues have that interesting property, that each of their voices is a piece of
   music in itself; and thus a fugue might be thought o collection of several distinct
   pieces of music, all based on one theme, and all played simultaneously. And it is up to
   the listener subconscious) to decide whether it should be perceived as a unit, c
   collection of independent parts, all of which harmonize.
Achilles: You say that the parts are "independent", yet that can't be literally true. There
   has to be some coordination between them, otherwise when they were put together
   one would just have an unsystematic clashing of tones-and that is as far from the truth
   as could b,
Anteater: A better way to state it might be this: if you listened to each on its own, you
   would find that it seemed to make sense all by its could stand alone, and that is the
   sense in which I meant that it is independent. But you are quite right in pointing out
   that each of individually meaningful lines fuses with the others in a highly nonrandom
   way, to make a graceful totality. The art of writing a beautiful fugue lies precisely in
   this ability, to manufacture several diff lines, each one of which gives the illusion of
   having been written I own beauty, and yet which when taken together form a whole, ,
   does not feel forced in any way. Now, this dichotomy between he a fugue as a whole,
   and hearing its component voices, is a part: example of a very general dichotomy,
   which applies to many kit structures built up from lower levels.
Achilles: Oh, really? You mean that my two "modes" may have some general type of
   applicability, in situations other than fugue-listening?
Anteater: Absolutely.
Achilles: I wonder how that could be. I guess it has to do with alternating between
   perceiving something as a whole, and perceiving it as a collection of parts. But the
   only place I have ever run into that dichotomy is in listening to fugues.
Tortoise: Oh, my, look at this! I just turned the page while following the music, and came
   across this magnificent illustration facing the page of the fugue.




Prelude                                                                                  289

Crab: I have never seen that illustration before. Why don't you pass it 'round?

   (The Tortoise passes the book around. Each of the foursome looks at it in a
   characteristic way-this one from afar, that one from close up, everyone tipping his
   head this way and that in puzzlement. Finally it has made the rounds, and returns
   to the Tortoise, who peers at it rather intently.)

Achilles: Well, I guess the prelude is just about over. I wonder if, as I listen to this fugue,
  I will gain any more insight into the question, "What is the right way to listen to a
  fugue: as a whole, or as the sum of its parts?"
TTortoise: Listen carefully, and you will!

(The prelude ends. There is a moment of silence; and ...

                                                                                [ATTACCA]




Prelude                                                                                    290


                          Levels of Description,
                         and Computer Systems

                                Levels of Description

GODEL'S STRING G, and a Bach fugue: they both have the property that they can be
understood on different levels. We are all familiar with this kind of thing; and yet in some
cases it confuses us, while in others w handle it without any difficulty at all. For example,
we all know that w human beings are composed of an enormous number of cells (around
twenty-five trillion), and therefore that everything we do could in principle be described
in terms of cells. Or it could even be described on the level c molecules. Most of us
accept this in a rather matter-of-fact way; we go t the doctor, who looks at us on lower
levels than we think of ourselves. W read about DNA and "genetic engineering" and sip
our coffee. We seem t have reconciled these two inconceivably different pictures of
ourselves simply by disconnecting them from each other. We have almost no way t relate
a microscopic description of ourselves to that which we feel ourselves to be, and hence it
is possible to store separate representations of ourselves in quite separate "compartments"
of our minds. Seldom do we have to fir back and forth between these two concepts of
ourselves, wondering "How can these two totally different things be the same me?"
        Or take a sequence of images on a television screen which show Shirley
MacLaine laughing. When we watch that sequence, we know that we are actually looking
not at a woman, but at sets of flickering dots on a flat surface. We know it, but it is the
furthest thing from our mind. We have these two wildly opposing representations of what
is on the screen, but that does not confuse us. We can just shut one out, and pay attention
to th other-which is what all of us do. Which one is "more real"? It depends o; whether
you're a human, a dog, a computer, or a television set.

                             Chunking and Chess Skill
One of the major problems of Artificial Intelligence research is to figure out how to
bridge the gap between these two descriptions; how to construe a system which can
accept one level of description, and produce the other One way in which this gap enters
Artificial Intelligence is well illustrated b the progress in knowledge about how to
program a computer to play goof chess. It used to be thought in the 1950's and on into the
1960's-that the




Levels of Description, and Computer Systems                                              291

trick to making a machine play well was to make the machine look further ahead into the
branching network of possible sequences of play than any chess master can. However, as
this goal gradually became attained, the level of computer chess did not have any sudden
spurt, and surpass human experts. In fact, a human expert can quite soundly and
confidently trounce the best chess programs of this day.
         The reason for this had actually been in print for many years. In the 1940's, the
Dutch psychologist Adriaan de Groot made studies of how chess novices and chess
masters perceive a chess situation. Put in their starkest terms, his results imply that chess
masters perceive the distribution of pieces in chunks. There is a higher-level description
of the board than the straightforward "white pawn on K5, black rook on Q6" type of
description, and the master somehow produces such a mental image of the board. This
was proven by the high speed with which a master could reproduce an actual position
taken from a game, compared with the novice's plodding reconstruction of the position,
after both of them had had five-second glances at the board. Highly revealing was the fact
that masters' mistakes involved placing whole groups of pieces in the wrong place, which
left the game strategically almost the same, but to a novice's eyes, not at all the same. The
clincher was to do the same experiment but with pieces randomly assigned to the squares
on the board, instead of copied from actual games. The masters were found to be simply
no better than the novices in reconstructing such random boards.
         The conclusion is that in normal chess play, certain types of situation recur-
certain patterns-and it is to those high-level patterns that the master is sensitive. He thinks
on a different level from the novice; his set of concepts is different. Nearly everyone is
surprised to find out that in actual play, a master rarely looks ahead any further than a
novice does-and moreover, a master usually examines only a handful of possible moves!
The trick is that his mode of perceiving the board is like a filter: he literally does not see
bad moves when he looks at a chess situation-no more than chess amateurs see illegal
moves when they look at a chess situation. Anyone who has played even a little chess has
organized his perception so that diagonal rook-moves, forward captures by pawns, and so
forth, are never brought to mind. Similarly, master-level players have built up higher
levels of organization in the way they see the board; consequently, to them, bad moves
are as unlikely to come to mind as illegal moves are, to most people. This might be called
implicit pruning of the giant branching tree of possibilities. By contrast, explicit pruning
would involve thinking of a move, and after superficial examination, deciding not to
pursue examining it any further.
         The distinction can apply just as well to other intellectual activities -- for instance,
doing mathematics. A gifted mathematician doesn't usually think up and try out all sorts
of false pathways to the desired theorem, as less gifted people might do; rather, he just
"smells" the promising paths, and takes them immediately.
         Computer chess programs which rely on looking ahead have not been taught to
think on a higher level; the strategy has just been to use brute




Levels of Description, and Computer Systems                                                 292

force look-ahead, hoping to crush all types of opposition. But it h worked. Perhaps
someday, a look-ahead program with enough brute ,gill indeed overcome the best human
players-but that will be a intellectual gain, compared to the revelation that intelligence de
crucially on the ability to create high-level descriptions of complex such as chess boards,
television screens, printed pages, or painting

                                     Similar Levels
usually, we are not required to hold more than one level of understanding of a situation in
our minds at once. Moreover, the different descriptions a single system are usually so
conceptually distant from each other tl was mentioned earlier, there is no problem in
maintaining them both are just maintained in separate mental compartments. What is
confusing though, is when a single system admits of two or more descriptions different
levels which nevertheless resemble each other in some way. we find it hard to avoid
mixing levels when we think about the system can easily get totally lost.
        Undoubtedly this happens when we think about our psychology-for instance,
when we try to understand people's motivations: for various actions. There are many
levels in the human m structure-certainly it is a system which we do not understand very
we But there are hundreds of rival theories which tell why people act the way they do,
each theory based on some underlying assumptions about he down in this set of levels
various kinds of psychological "forces" are f( Since at this time we use pretty much the
same kind of language f mental levels, this makes for much level-mixing and most
certain] hundreds of wrong theories. For instance, we talk of "drives"-for se power, for
fame, for love, etc., etc.-without knowing where these drives come from in the human
mental structure. Without belaboring the pc simply wish to say that our confusion about
who we are is certainly r( to the fact that we consist of a large set of levels, and we use
overlapping language to describe ourselves on all of those levels.

                                  Computer Systems
There is another place where many levels of description coexist for a system, and where
all the levels are conceptually quite close to one an( I am referring to computer systems.
When a computer program is ping, it can be viewed on a number of levels. On each level,
the description is given in the language of computer science, which makes all the de
descriptions similar in some ways to each other-yet there are extremely imp( differences
between the views one gets on the different levels. At the 1 level, the description can be
so complicated that it is like the dot-description of a television picture. For some
purposes, however, this is by far the important view. At the highest level, the description
is greatly chunked and




Levels of Description, and Computer Systems                                              293

takes on a completely different feel, despite the fact that many of the same concepts
appear on the lowest and highest levels. The chunks on the high-level description are like
the chess expert's chunks, and like the chunked description of the image on the screen:
they summarize in capsule form a number of things which on lower levels are seen as
separate. (See Fig. 57.) Now before things become too abstract, let us pass on to the




FIGURE 57. The idea of "chunking": a group of items is reperceived as a single "chunk".
The chunk's boundary is a little like a cell membrane or a national border: it establishes
a separate identity for the cluster within. According to context, one may wish to ignore
the chunk's internal structure or to take it into account.

concrete facts about computers, beginning with a very quick skim of what a computer
system is like on the lowest level. The lowest level? Well, not really, for I am not going
to talk about elementary particles-but it is the lowest level which we wish to think about.
        At the conceptual rock-bottom of a computer, we find a memory, a central
processing unit (CPU), and some input-output (I/O) devices. Let us first describe the
memory. It is divided up into distinct physical pieces, called words. For the sake of
concreteness, let us say there are 65,536 words of memory (a typical number, being 2 to
the 16th power). A word is further divided into what we shall consider the atoms of
computer science-bits. The number of bits in a typical word might be around thirty-six.
Physically, a bit is just a magnetic "switch" that can be in either of two positions.




--- a word of 36 bits --




Levels of Description, and Computer Systems                                            294

you could call the two positions "up" and "down", or "x" and "o", o and "0" ... The third
is the usual convention. It is perfectly fine, but i the possibly misleading effect of making
people think that a comp deep down, is storing numbers. This is not true. A set of thirty-
six bits not have to be thought of as a number any more than two bits has i thought of as
the price of an ice cream cone. Just as money can do va things depending on how you use
it, so a word in memory can serve r functions. Sometimes, to be sure, those thirty-six bits
will indeed repn a number in binary notation. Other times, they may represent thin dots
on a television screen. And other times, they may represent a letters of text. How a word
in memory is to be thought of depends eni on the role that this word plays in the program
which uses it. It ma course, play more than one role-like a note in a canon.

                                Instructions and Data
There is one interpretation of a word which I haven't yet mentioned, that is as an
instruction. The words of memory contain not only data t acted on, but also the program
to act on the data. There exists a lin repertoire of operations which can be carried out by
the central proce5 unit-the CPU-and part of a word, usually its first several bits-is it
pretable as the name of the instruction-type which is to be carried What do the rest of the
bits in a word-interpreted-as-instruction stand Most often, they tell which other words in
memory are to be acted upoi other words, the remaining bits constitute a pointer to some
other wor( words) in memory. Every word in memory has a distinct location, li house on
a street; and its location is called its address. Memory may have "street", or many
"streets"-they are called "pages". So a given wo addressed by its page number (if memory
is paged) together wit position within the page. Hence the "pointer" part of an instruction
i numerical address of some word(s) in memory. There are no restric on the pointer, so an
instruction may even "point" at itself, so that whet executed, it causes a change in itself to
be made.
How does the computer know what instruction to execute at any € time? This is kept
track of in the CPU. The CPU has a special pointer w points at (i.e., stores the address of)
the next word which is to be inter ed as an instruction. The CPU fetches that word from
memory, and c it electronically into a special word belonging to the CPU itself. (Wor the
CPU are usually not called "words", but rather, registers.) Then the executes that
instruction. Now the instruction may call for any of a number of types of operations to be
carried out. Typical ones include:

       ADD the word pointed to in the instruction, to a register.
             (In this case, the word pointed to is obviously interpreted as number.)




Levels of Description, and Computer Systems                                               295

               PRINT the word pointed to in the instruction, as letters.
               (In this case, the word is obviously interpreted not as a number, but as a
       string of letters.)
               JUMP to the word pointed to in the instruction.
               (In this case, the CPU is being told to interpret that particular word as its
       next instruction.)

               Unless the instruction explicitly dictates otherwise, the CPU will pick up
       the very next word and interpret it as an instruction. In other words, the CPU
       assumes that it should move down the "street" sequentially, like a mailman,
       interpreting word after word as an instruction. But this sequential order can be
       broken by such instructions as the JUMP instruction, and others.

                              Machine Language vs. Assembly language

       This is a very brief sketch of machine language. In this language, the types of
       operations which exist constitute a finite repertoire which cannot be extended.
       Thus all programs, no matter how large and complex, must be made out of
       compounds of these types of instructions. Looking at a program written in
       machine language is vaguely comparable to looking at a DNA molecule atom by
       atom. If you glance back to Fig. 41, showing the nucleotide sequence of a DNA
       molecule--and then if you consider that each nucleotide contains two dozen atoms
       or so-and if you imagine trying to write the DNA, atom by atom, for a small virus
       (not to mention a human being!)-then you will get a feeling for what it is like to
       write a complex program in machine language, and what it is like to try to grasp
       what is going on in a program if you have access only to its machine language
       description. ,
               It must be mentioned, however, that computer programming was
       originally done on an even lower level, if possible, than that of machine language-
       -namely, connecting wires to each other, so that the proper operations were "hard-
       wired" in. This is so amazingly primitive by modern standards that it is painful
       even to' imagine. Yet undoubtedly the people who first did it experienced as much
       exhilaration as the pioneers of modern computers ever do .. .
               We now wish to move to a higher level of the hierarchy of levels of
       description of programs. This is the assembly language level. There is not a
       gigantic spread between assembly language and machine language; indeed, the
       step is rather gentle. In essence, there is a one-to-one correspondence between
       assembly language instructions and machine language instructions. The idea of
       assembly language is to "chunk" the individual machine language instructions, so
       that instead of writing the sequence of bits "010111000" when you want an
       instruction which adds one number to another, you simply write ADD, and then
       instead of giving the address in binary representation, you can refer to the word in
       memory by a name.




Levels of Description, and Computer Systems                                             296

       Therefore, a program in assembly language is very much like a machine language
       program made legible to humans. You might compare the machine language
       version of a program to a TNT-derivation done in the obscure Gödel-numbered
       notation, and the assembly language version to the isomorphic TNT-derivation,
       done in the original TNT-notation, which is much easier to understand. Or, going
       back to the DNA image, we can liken the difference between machine language
       and assembly language to the difference between painfully specifying each
       nucleotide, atom by atom, and specifying a nucleotide by simply giving its name
       (i.e., 'A', 'G', 'C', or 'T'). There is a tremendous saving of labor in this very
       simple "chunking" operation, although conceptually not much has been changed.

                              Programs That Translate Programs

       Perhaps the central point about assembly language is not its differences from
       machine language, which are not that enormous, but just the key idea that
       programs could be written on a different level at all! Just think about it: the
       hardware is built to "understand" machine language programs-sequences of bits-
       but not letters and decimal numbers. What happens when hardware is fed a
       program in assembly language\% It is as if you tried to get a cell to accept a piece
       of paper with the nucleotide sequence written out in letters of the alphabet, instead
       of in chemicals. What can a cell do with a piece of paper? What can a computer
       do with an assembly language program?
               And here is the vital point: someone can write, in machine language, a
       translation program. This program, called an assembler, accepts mnemonic
       instruction names, decimal numbers, and other convenient abbreviations which a
       programmer can remember easily, and carries out the conversion into the
       monotonous but critical bit-sequences. After the assembly language program has
       been assembled (i.e., translated), it is run-or rather, its machine language
       equivalent is run. But this is a matter of terminology. Which level program is
       running? You can never go wrong if you say that the machine language program
       is running, for hardware is always involved when any program runs-but it is also
       quite reasonable to think of the running program in terms of assembly language.
       For instance, you might very well say, "Right now, the CPU is executing a JUMP
       instruction", instead of saying, "Right now, the CPU is executing a ' 1 11010000'
       instruction". A pianist who plays the notes G-E-B E-G-B is also playing an
       arpeggio in the chord of E minor. There is no reason to be reluctant about
       describing things from a higher-level point of view. So one can think of the
       assembly language program running concurrently with the machine language
       program. We have two modes of describing what the CPU is doing.




Levels of Description, and Computer Systems                                             297

                        Higher-Level Languages, Compilers, and Interpreters

       The next level of the hierarchy carries much further the extremely powerful idea
       of using the computer itself to translate programs from a high level into lower
       levels. After people had programmed in assembly language for a number of years,
       in the early 1950's, they realized that there were a number of characteristic
       structures which kept reappearing in program after program. There seemed to be,
       just as in chess, certain fundamental patterns which cropped up naturally when
       human beings tried to formulate algorithms--exact descriptions of processes they
       wanted carried out. In other words, algorithms seemed to have certain higher-
       level components, in terms of which they could be much more easily and
       esthetically specified than in the very restricted machine language, or assembly
       language. Typically, a high-level algorithm component consists not of one or two
       machine language instructions, but of a whole collection of them, not necessarily
       all contiguous in memory. Such a component could be represented in a higher-
       level language by a single item-a chunk.
               Aside from standard chunks-the newly discovered components out of
       which all algorithms can be built-people realized that almost all programs contain
       even larger chunks-superchunks, so to speak. These superchunks differ from
       program to program, depending on the kinds of high-level tasks the j program is
       supposed to carry out. We discussed superchunks in Chapter V, calling them by
       their usual names: "subroutines" and "procedures". It was clear that a most
       powerful addition to any programming language would be the ability to define
       new higher-level entities in terms of previously known ones, and then to call them
       by name. This would build the chunking operation right into the language. Instead
       of there being a determinate repertoire of instructions out of which all programs
       had to be explicitly assembled, the programmer could construct his own modules,
       each with its own name, each usable anywhere inside the program, just as if it had
       been a built-in feature of the language. Of course, there is no getting away from
       the fact that down below, on a machine language level, everything would still be
       composed of the same old machine language instructions, but that would not be
       explicitly visible to the highlevel programmer; it would be implicit.
               The new languages based on these ideas were called compiler languages.
       One of the earliest and most elegant was called "Algol", for "Algorithmic
       Language". Unlike the case with assembly language, there is no straightforward
       one-to-one correspondence between statements in Algol and machine language
       instructions. To be sure, there is still a type of mapping from Algol into machine
       language, but it is far more "scrambled" than that between assembly language and
       machine language. Roughly speaking, an Algol program is to its machine
       language translation as a word problem in an elementary algebra text is to the
       equation it translates into. (Actually, getting from a word problem to an equation
       is far more complex, but it gives some inkling of the types of "unscrambling" that
       have to be carried out in translating from a high-level language to a lower-level
       Ian




Levels of Description, and Computer Systems                                          298

       guage.) In the mid-1950's, successful programs called compilers were written
       whose function was to carry out the translation from compiler languages to
       machine language.
                Also, interpreters were invented. Like compilers, interpreters translate
       from high-level languages into machine language, but instead of translating all the
       statements first and then executing the machine code, they read one line and'
       execute it immediately. This has the advantage that a user need not have written a
       complete program to use an interpreter. He may invent his program line by line,
       and test it out as he goes along. Thus, an interpreter is to a compiler as a
       simultaneous interpreter is to a translator of a written speech. One of the most
       important and fascinating of all computer languages is LISP (standing for "List
       Processing"), which was invented by John McCarthy around the time Algol was
       invented. Subsequently, LISP has enjoyed great popularity with workers in
       Artificial Intelligence.
                There is one interesting difference between the way interpreters work and
       compilers work. A compiler takes input (a finished Algol program, for instance)
       and produces output (a long sequence of machine language instructions). At this
       point, the compiler has done its duty. The output is then given to the computer to
       run. By contrast, the interpreter is constantly running while the programmer types
       in one LISP statement after another, and each one gets executed then' and there.
       But this doesn't mean that each statement gets first translated, then executed, for
       then an interpreter would be nothing but a line-by-line compiler. Instead, in an
       interpreter, the operations of reading a new line, "understanding" it, and executing
       it are intertwined: they occur simultaneously.
                Here is the idea, expanded a little more. Each time a new line of LISP is
       typed in, the interpreter tries to process it. This means that the interpreter jolts into
       action, and certain (machine language) instructions inside it get executed.
       Precisely which ones get executed depends on the LISP statement itself, of
       course. There are many JUMP instructions inside the interpreter, so that the new
       line of LISP may cause control to move around in a complex way-forwards,
       backwards, then forwards again, etc.. Thus, each LISP statement gets converted
       into a "pathway" inside the interpreter, and the act of following that pathway
       achieves the desired effect.
                Sometimes it is helpful to think of the LISP statements as mere pieces of
       data which are fed sequentially to a constantly running machine language
       program (the LISP interpreter). When you think of things this way, you get a
       different image of the relation between a program written in a higher-level
       language and the machine which is executing it.

                                               Bootstrapping

       Of course a compiler, being itself a program, has to be written in some language.
       The first compilers were written in assembly language, rather than machine
       language, thus taking full advantage of the already ac-




Levels of Description, and Computer Systems                                                299

       omplished first step up from machine language. A summary of these rather tricky
       concepts is presented in Figure 58.

                                                     FIGURE 58. Assemblers and
                                                     compilers are both translators into
                                                     machine language. This is indicated
                                                     by the solid lines. Moreover, since
                                                     they are themselves programs, they
                                                     are originally written in a language
                                                     also. The wavy lines indicate that aa
                                                     compiler can be written in assembly
                                                     language, and an assembler in
                                                     machine language.




       Now as sophistication increased, people realized that a partially written compiler
       could be used to compile extensions of itself. In other words, once i certain
       minimal core of a compiler had been written, then that minimal compiler could
       translate bigger compilers into machine language-which n turn could translate yet
       bigger compilers, until the final, full-blown :compiler had been compiled. This
       process is affectionately known as `bootstrapping"-for obvious reasons (at least if
       your native language is English it is obvious). It is not so different from the
       attainment by a child of a critical level of fluency in his native language, from
       which point on his 'vocabulary and fluency can grow by leaps and bounds, since
       he can use language to acquire new language.

                       Levels on Which to Describe Running Programs

       Compiler languages typically do not reflect the structure of the machines which
       will run programs written in them. This is one of their chief advantages over the
       highly specialized assembly and machine languages. Of course, when a compiler
       language program is translated into machine language, the resulting program is
       machine-dependent. Therefore one can describe a program which is being
       executed in a machine-independent way or a machine-dependent way. It is like
       referring to a paragraph in a book by its subject matter (publisher-independent), or
       its page number and position on the page (publisher-dependent).
               As long as a program is running correctly, it hardly matters how you
       describe it or think of its functioning. It is when something goes wrong that




Levels of Description, and Computer Systems                                            300

       it is important to be able to think on different levels. If, for instance, the machine
       is instructed to divide by zero at some stage, it will come to a halt and let the user
       know of this problem, by telling where in the program the questionable event
       occurred. However, the specification is often given on a lower level than that in
       which the programmer wrote the program. Here are three parallel descriptions of
       a program grinding to a halt:

       Machine Language Level:
             "Execution of the program stopped in location 1110010101110111"

       Assembly Language Level•:
             "Execution of the program stopped when the DIV (divide) instruction was
             hit"

       Compiler Language Level:
             "Execution of the program stopped during evaluation of the algebraic
             expression `(A + B)/Z'

       One of the greatest problems for systems programmers (the people who write
       compilers, interpreters, assemblers, and other programs to be used by many
       people) is to figure out how to write error-detecting routines in such a way that
       the messages which they feed to the user whose program has a "bug" provide
       high-level, rather than low-level, descriptions of the problem. It is an interesting
       reversal that when something goes wrong in a genetic "program" (e.g., a
       mutation), the "bug" is manifest only to people on a high level-namely on the
       phenotype level, not the genotype level. Actually, modern biology uses mutations
       as one of its principal windows onto genetic processes, because of their multilevel
       traceability.

                         Microprogramming and Operating Systems

       In modern computer systems, there are several other levels of the hierarchy. For
       instance, some systems-often the so-called "microcomputers" come with machine
       language instructions which are even more rudimentary than the instruction to add
       a number in memory to a number in a register. It is up to the user to decide what
       kinds of ordinary machine-level instructions he would like to be able to program
       in; he "microprograms" these instructions in terms of the "micro-instructions"
       which are available. Then the "higher-level machine language" instructions which
       he has designed may be burned into the circuitry and become hard-wired,
       although they need not be. Thus microprogramming allows the user to step a little
       below the conventional machine language level. One of the consequences is that a
       computer of one manufacturer can be hard-wired (via microprogramming) so as
       to have the same machine language instruction set as a computer of the same, or
       even another, manufacturer. The microprogrammed computer is said to be
       "emulating" the other computer. Then there is the level of the operating system,
       which fits between the

Levels of Description, and Computer Systems                                              301

       machine language program and whatever higher level the user is programming in.
       The operating system is itself a program which has the functions of shielding the
       bare machine from access by users (thus protecting the system), and also of
       insulating the programmer from the many extremely intricate and messy problems
       of reading the program, calling a translator, running the translated program,
       directing the output to the proper channels at the proper time, and passing control
       to the next user. If there are several users "talking" to the same CPU at once, then
       the operating system is the program that shifts attention from one to the other in
       some orderly fashion. The complexities of operating systems are formidable
       indeed, and I shall only hint at them by the following analogy.
               Consider the first telephone system. Alexander Graham Bell could phone
       his assistant in the next room: electronic transmission of a voice! Now that is like
       a bare computer minus operating system: electronic computation! Consider now a
       modern telephone system. You have a choice of other telephones to connect to.
       Not only that, but many different calls can be handled simultaneously. You can
       add a prefix and dial into different areas. You can call direct, through the
       operator, collect, by credit card, person-to-person, on a conference call. You can
       have a call rerouted or traced. You can get a busy signal. You can get a siren-like
       signal that says that the number you dialed isn't "well-formed", or that you have
       taken too in long in dialing. You can install a local switchboard so that a group of
       phones are all locally connected--etc., etc. The list is amazing, when you think of
       how much flexibility there is, particularly in comparison to the erstwhile miracle
       of a "bare" telephone. Now sophisticated operating systems carry out similar
       traffic-handling and level-switching operations with respect to users and their
       programs. It is virtually certain that there are somewhat parallel things which take
       place in the brain: handling of many stimuli at the same time; decisions of what
       should have priority over what and for how long; instantaneous "interrupts"
       caused by emergencies or other unexpected occurrences; and so on.

                       Cushioning the User and Protecting the System

       The many levels in a complex computer system have the combined effect of
       "cushioning" the user, preventing him from having to think about the many lower-
       level goings-on which are most likely totally irrelevant to him anyway. A
       passenger in an airplane does not usually want to be aware of the levels of fuel in
       the tanks, or the wind speeds, or how many chicken dinners are to be served, or
       the status of the rest of the air traffic around the destination-this is all left to
       employees on different levels of the airlines hierarchy, and the passenger simply
       gets from one place to another. Here again, it is when something goes wrong-such
       as his baggage not arriving that the passenger is made aware of the confusing
       system of levels underneath him.




Levels of Description, and Computer Systems                                            302

                       Are Computers Super-Flexible or Super-Rigid?

       One of the major goals of the drive to higher levels has always been to make as
       natural as possible the task of communicating to the computer what you want it to
       do. Certainly, the high-level constructs in compiler languages are closer to the
       concepts which humans naturally think in, than are lower-level constructs such as
       those in machine language. But in this drive towards ease of communication, one
       aspect of "naturalness" has been quite neglected. That is the fact that interhuman
       communication is far less rigidly constrained than human-machine
       communication. For instance, we often produce meaningless sentence fragments
       as we search for the best way to express something, we cough in the middle of
       sentences, we interrupt each other, we use ambiguous descriptions and "improper"
       syntax, we coin phrases and distort meanings-but our message still gets through,
       mostly. With programming languages, it has generally been the rule that there is a
       very strict syntax which has to be obeyed one hundred per cent of the time; there
       are no ambiguous words or constructions. Interestingly, the printed equivalent of
       coughing (i.e., a nonessential or irrelevant comment) is allowed, but only
       provided it is signaled in advance by a key word (e.g., COMMENT), and then
       terminated by another key word (e.g., a semicolon). This small gesture towards
       flexibility has its own little pitfall, ironically: if a semicolon (or whatever key
       word is used for terminating a comment) is used inside a comment, the translating
       program will interpret that semicolon as signaling the end of the comment, and
       havoc will ensue.
       If a procedure named INSIGHT has been defined and then called seventeen times
       in the program, and the eighteenth time it is misspelled as INSIHGT, woe to the
       programmer. The compiler will balk and print a rigidly unsympathetic error
       message, saying that it has never heard of INSIHGT. Often, when such an error
       is detected by a compiler, the compiler tries to continue, but because of its lack of
       insihgt, it has not understood what the programmer meant. In fact, it may very
       well suppose that something entirely different was meant, and proceed under that
       erroneous assumption. Then a long series of error messages will pepper the rest of
       the program, because the compiler-not the programmer-got confused. Imagine the
       chaos that would result if a simultaneous English-Russian interpreter, upon
       hearing one phrase of French in the English, began trying to interpret all the
       remaining English as French. Compilers often get lost in such pathetic ways. C'est
       la vie.
       Perhaps this sounds condemnatory of computers, but it is not meant to be. In some
       sense, things had to be that way. When you stop to think what most people use
       computers for, you realize that it is to carry out very definite and precise tasks,
       which are too complex for people to do. If the computer is to be reliable, then it is
       necessary that it should understand, without the slightest chance of ambiguity,
       what it is supposed to do. It is also necessary that it should do neither more nor
       less than it is explicitly instructed to do. If there is, in the cushion underneath the
       programmer, a program whose purpose is to "guess" what the programmer wants
       or

Levels of Description, and Computer Systems                                               303

       means, then it is quite conceivable that the programmer could try to communicate
       his task and be totally misunderstood. So it is important that the high-level
       program, while comfortable for the human, still should be unambiguous and
       precise.

                               Second-Guessing the Programmer

       Now it is possible to devise a programming language-and a program which
       translates it into the lower levels-which allows some sorts of imprecision. One
       way of putting it would be to say that a translator for such a programming
       language tries to make sense of things which are done "outside of the rules of the
       language". But if a language allows certain "transgressions", then transgressions
       of that type are no longer true transgressions, because they have been included
       inside the rules' If a programmer is aware that he may make certain types of
       misspelling, then he may use this feature of the language deliberately, knowing
       that he is actually operating within the rigid rules of the language, despite
       appearances. In other words, if the user is aware of all the flexibilities
       programmed into the translator for his convenience, then he knows the bounds
       which he cannot overstep, and therefore, to him, the translator still appears rigid
       and inflexible, although it may allow him much more freedom than early versions
       of the language, which did not incorporate "automatic compensation for human
       error".
                With "rubbery" languages of that type, there would seem to be two
       alternatives: (1) the user is aware of the built-in flexibilities of the language and
       its translator; (2) the user is unaware of them. In the first case, the language is still
       usable for communicating programs precisely, because the programmer can
       predict how the computer will interpret the programs he writes in the language. In
       the second case, the "cushion" has hidden features which may do things that are
       unpredictable (from the vantage point of a user who doesn't know the inner
       workings of the translator). This may result in gross misinterpretations of
       programs, so such a language is unsuitable for purposes where computers are used
       mainly for their speed and reliability.
                Now there is actually a third alternative: (3) the user is aware of the built-
       in flexibilities of the language and its translator, but there are so many of them
       and they interact with each other in such a complex way that he cannot tell how
       his programs will be interpreted. This may well apply to the person who wrote the
       translating program; he certainly knows its insides as well as anyone could-but he
       still may not be able to anticipate how it will react to a given type of unusual
       construction.
                One of the major areas of research in Artificial Intelligence today is called
       automatic programming, which is concerned with the development of yet higher-
       level languages-languages whose translators are sophisticated. in that they can do
       at least some of the following impressive things: generalize from examples,
       correct some misprints or grammatical errors,




Levels of Description, and Computer Systems                                                304

       try to make sense of ambiguous descriptions, try to second-guess the user by
       having a primitive user model, ask questions when things are unclear, use English
       itself, etc. The hope is that one can walk the tightrope between reliability and
       flexibility.

                             AI Advances Are Language Advances

       It is striking how tight the connection is between progress in computer science
       (particularly Artificial Intelligence) and the development of new languages. A
       clear trend has emerged in the last decade: the trend to consolidate new types of
       discoveries in new languages. One key for the understanding and creation of
       intelligence lies in the constant development and refinement of the languages in
       terms of which processes for symbol manipulation are describable. Today, there
       are probably three or four dozen experimental languages which have been
       developed exclusively for Artificial Intelligence research. It is important to realize
       that any program which can be written in one of these languages is in principle
       programmable in lower-level languages, but it would require a supreme effort for
       a human; and the resulting program would be so long that it would exceed the
       grasp of humans. It is not that each higher level extends the potential of the
       computer; the full potential of the computer already exists in its machine language
       instruction set. It is that the new concepts in a high-level language suggest
       directions and perspectives by their very nature.
                The "space" of all possible programs is so huge that no one can have a
       sense of what is possible. Each higher-level language is naturally suited for
       exploring certain regions of "program space"; thus the programmer, by using that
       language, is channeled into those areas of program space. He is not forced by the
       language into writing programs of any particular type, but the language makes it
       easy for him to do certain kinds of things. Proximity to a concept, and a gentle
       shove, are often all that is needed for a major discovery-and that is the reason for
       the drive towards languages of ever higher levels.
                Programming in different 'languages is like composing pieces in different
       keys, particularly if you work at the keyboard. If you have learned or written
       pieces in many keys, each key will have its own special emotional aura. Also,
       certain kinds of figurations "lie in the hand" in one key but are awkward in
       another. So you are channeled by your choice of key. In some ways, even
       enharmonic keys, such as C-sharp and D-flat, are quite distinct in feeling. This
       shows how a notational system can play a significant role in shaping the final
       product.
                A "stratified" picture of Al is shown in Figure 59, -with machine
       components such as transistors on the bottom, and "intelligent programs" on the
       top. The picture is taken from the book Artificial Intelligence by Patrick Henry
       Winston, and it represents a vision of Al shared by nearly all Al workers.
       Although I agree with the idea that Al must be stratified in some such way, I do
       not think that, with so few layers, intelligent programs



Levels of Description, and Computer Systems                                              305

                                                FIGURE 59. To create intelligent
                                                programs, one needs to build up a series
                                                of levels of hardware and software, so
                                                that one is spared the agonT of seeing
                                                everything only on the lowest level.
                                                Descriptions of a single process on
                                                different levels will sound verb different
                                                from each other, only the top one being
                                                sufficiently    chunked    that    it    is
                                                comprehensible to us. [Adapted from P.
                                                H. Winston, Artificial Intelligence
                                                (Reading, Mass.: Addison-ifele'', 1977)]


       can he reached. Between the machine language level and the level where rue
       intelligence will be reached, I am convinced there will lie perhaps mother dozen
       (or even several dozen!) layers, each new layer building on and extending the
       flexibilities of the layer below. What they will be like we can hardly dream of
       now ...


                               The Paranoid and the Operating System

       The similarity of all levels in a computer system can lead to some strange level-
       mixing experiences. I once watched a couple of friends-both computer novices-
       playing with the program "PARRY" on a terminal. PARRY s a. rather infamous
       program which simulates a paranoid in an extremely rudimentary way, by spitting
       out canned phrases in English chosen from a vide repertoire; its plausibility is due
       to its ability to tell which of its stock phrases might sound reasonable in response
       to English sentences typed to t by a human.
                At one point, the response time got very sluggish-PARRY was taking very
       long to reply-and I explained to my friends that this was probably because of the
       heavy load on the time-sharing system. I told them they could find out how many
       users were logged on, by typing a special "control" character which would go
       directly to the operating system, and would )e unseen by PARRY. One of my
       friends pushed the control character. In a lash, some internal data about the
       operating system's status overwrote some of PARRY's words on the screen.
       PARRY knew nothing of this: it is a program with "knowledge" only of horse
       racing and bookies-not operating systems and terminals and special control
       characters. But to my friends, both PARRY and the operating system were just
       "the computer"-a mysterious, remote, amorphous entity that responded to them
       when they typed. And so it made perfect sense when one of them blithely typed,
       in 3nglish, "Why are you overtyping what's on the screen?" The idea that PARRY
       could know' nothing about the operating system it was running



Levels of Description, and Computer Systems                                            306

       under was not clear to my friends. The idea that "you" know all about "yourself"
       is so familiar from interaction with people that it was natural to extend it to the
       computer-after all, it was intelligent enough that it could "talk" to them in
       English! Their question was not unlike asking a person, "Why are you making so
       few red blood cells today?" People do not know about that level-the "operating
       system level"-of their bodies.
                The main cause of this level-confusion was that communication with all
       levels of the computer system was taking place on a single screen, on a single
       terminal. Although my friends' naiveté might seem rather extreme, even
       experienced computer people often make similar errors when several levels of a
       complex system are all present at once on the same screen. They forget "who"
       they are talking to, and type something which makes no sense at that level,
       although it would have made perfect sense on another level. It might seem
       desirable, therefore, to have the system itself sort out the levels-to interpret
       commands according to what "makes sense". Unfortunately, such interpretation
       would require the system to have a lot of common sense, as well as perfect
       knowledge of the programmer's overall intent-both of which would require more
       artificial intelligence than exists at the present time.

                            The Border between Software and Hardware

       One can also be confused by the flexibility of some levels and the rigidity of
       others. For instance, on some computers there are marvelous text-editing systems
       which allow pieces of text to be "poured" from one format into another,
       practically as liquids can be poured from one vessel into another. A thin page can
       turn into a wide page, or vice versa. With such power, you might expect that it
       would be equally trivial to change from one font to another-say from roman to
       italics. Yet there may be only a single font available on the screen, so that such
       changes are impossible. Or it may be feasible on the screen but not printable by
       the printer-or the other way around. After dealing with computers for a long time,
       one gets spoiled, and thinks that everything should be programmable: no printer
       should be so rigid as to have only one character set, or even a finite repertoire of
       them-typefaces should be user-specifiable! But once that degree of flexibility has
       been attained, then one may be annoyed that the printer cannot print in different
       colors of ink, or that it cannot accept paper of all shapes and sizes, or that it does
       not fix itself when it breaks ...
                The trouble is that somewhere, all this flexibility has to "bottom out", to
       use the phrase from Chapter V. There must be a hardware level which underlies it
       all, and which is inflexible. It may lie deeply hidden, and there may be so much
       flexibility on levels above it that few users feel the hardware limitations-but it is
       inevitably there.
                What is this proverbial distinction between software and hardware? It is
       the distinction between programs and machines-between long complicated
       sequences of instructions, and the physical machines which carry




Levels of Description, and Computer Systems                                              307

       them out. I like to think of software as "anything which you could send over he
       telephone lines", and hardware as "anything else". A piano is hardware, gut
       printed music is software. A telephone set is hardware, but a telephone lumber is
       software. `The distinction is a useful one, but not always so clear-cut.
               We humans also have "software" and "hardware" aspects, and the
       difference is second nature to us. We are used to the rigidity of our physiology:
       the fact that we cannot, at will, cure ourselves of diseases, or ;row hair of any
       color-to mention just a couple of simple examples. We an, however, "reprogram"
       our minds so that we operate in new conceptual frameworks. The amazing
       flexibility of our minds seems nearly irreconcilable with the notion that our brains
       must be made out of fixed-rule hardware, which cannot be reprogrammed. We
       cannot make our neurons ire faster or slower, we cannot rewire our brains, we
       cannot redesign the interior of a neuron, we cannot make anti choices about the
       hardware-and 'et, we can control how we think.
               But there are clearly aspects of thought which are beyond our control. We
       cannot make ourselves smarter by an act of will; we cannot learn a new language
       as fast as we want; we cannot make ourselves think faster than we lo; we cannot
       make ourselves think about several things at once; and so on. This is a kind of
       primordial self-knowledge which is so obvious that it is lard to see it at all; it is
       like being conscious that the air is there. We never really bother to think about
       what might cause these "defects" of our minds: lamely, the organization of our
       brains. To suggest ways of reconciling the software of mind with the hardware of
       brain is a main goal of this book.

                                 Intermediate Levels and the Weather

       We have seen that in computer systems, there are a number of rather sharply
       defined strata, in terms of any one of which the operation of a running program
       can be described. Thus there is not merely a single low bevel and a single high
       level-there are all degrees of lowness and highness. s the existence of intermediate
       levels a general feature of systems which lave low and high levels? Consider, for
       example, the system whose 'hardware" is the earth's atmosphere (not very hard,
       but no matter), and whose "software" is the weather. Keeping track of the motions
       of all of the molecules simultaneously would be a very low-level way of
       "understanding" he weather, rather like looking at a huge, complicated program
       on the machine language level. Obviously it is way beyond human
       comprehension. 3ut we still have our own peculiarly human ways of looking at,
       and describing, weather phenomena. Our chunked view of the weather is based >n
       very high-level phenomena, such as: rain, fog, snow, hurricanes, cold fronts,
       seasons, pressures, trade winds, the jet stream, cumulo-nimbus clouds,
       thunderstorms, inversion layers, and so on. All of these phenomena involve
       astronomical numbers of molecules, somehow behaving in concert o that large-
       scale trends emerge. This is a little like looking at the weather n a compiler
       language.




Levels of Description, and Computer Systems                                             308

               Is there something analogous to looking at the weather in an intermediate-
       level language, such as assembly language? For instance, are there very small
       local "mini-storms", something like the small whirlwinds which one occasionally
       sees, whipping up some dust in a swirling column a few feet wide, at most? Is a
       local gust of wind an intermediate-level chunk which plays a role in creating
       higher-level weather phenomena? Or is there just no practical way of combining
       knowledge of such kinds of phenomena to create a more comprehensive
       explanation of the weather?
               Two other questions come to my mind. The first is: "Could it be that the
       weather phenomena which we perceive on our scale-a tornado, a drought-are just
       intermediate-level phenomena: parts of vaster, slower phenomena?" If so, then
       true high-level weather phenomena would be global, and their time scale would
       be geological. The Ice Age would be a high-level weather event. The second
       question is: "Are there intermediate level weather phenomena which have so far
       escaped human perception, but which, if perceived, could give greater insight into
       why the weather is as it is?"

                                       From Tornados to Quarks

               This last suggestion may sound fanciful, but it is not all that far-fetched.
       We need only look to the hardest of the hard sciences-physics-to find peculiar
       examples of systems which are explained in terms of interacting "parts" which are
       themselves invisible. In physics, as in any other discipline, a system is a group of
       interacting parts. In most systems that we know, the parts retain their identities
       during the interaction, so that we still see the parts inside the system. For
       example, when a team of football players assembles, the individual players retain
       their separateness-they do not melt into some composite entity, in which their
       individuality is lost. Still-and this is important-some processes are going on in
       their brains which are evoked by the team-context, and which would not go on
       otherwise, so that in a minor way, the players change identity when they become
       part of the larger system, the team. This kind of system is called a nearly
       decomposable system (the term comes from H. A. Simon's article "The
       Architecture of Complexity"; see the Bibliography). Such a system consists of
       weakly interacting modules, each of which maintains its own private identity
       throughout the interaction but by becoming slightly different from how it is when
       outside of the system,, contributes to the cohesive behavior of the whole system.
       The systems studied in physics are usually of this type. For instance, an atom is
       seen as made of 'a nucleus whose positive charge captures a number of electrons
       in "orbits", or bound states. The bound electrons are very much like free electrons,
       despite their being internal to a composite object.
               Some systems studied in physics offer a contrast to the relatively
       straightforward atom. Such systems involve extremely strong interactions, as a
       result of which the parts are swallowed up into the larger system, and lose some
       or all of their individuality. An example of this is the nucleus of an atom, which is
       usually described as being "a collection of protons and



Levels of Description, and Computer Systems                                             309

       neutrons". But the forces which pull the component particles together are strong
       that the component particles do not survive to anything like their “free" form (the
       form they have when outside a nucleus). And in fact a nucleus acts in many ways
       as a single particle, rather than as a collection of interacting particles. When a
       nucleus is split. protons and neutrons are ten released. but also other particles.
       such as pi-mesons and gamma rays, are commonly produced. Are all those
       different particles physically present side a nucleus before it is split, or are then
       just "sparks" which fly off ten the nucleus is split- It is perhaps not meaningful to
       try to give an answer to such a question. On the level of particle physics, the
       difference between storing the potential to make "sparks" and storing actual sub
       particles is not so clear.
               A nucleus is thus one systems whose "parts!, even though they are not
       visible while on the inside, can be pulled out and made risible. However, ere are
       more pathological cases. such as the proton and neutron seen as stems themselves.
       Each of them has been hypothesized to be constituted from a trio of "quarks"-
       hypothetical particles which can be combined in twos or threes to make many
       known fundamental particles. However, the interaction between quarks is so
       strong that not only can they not he seen [side the proton and neutron, but they
       cannot even be pulled out at all'. bus, although quarks help to give a theoretical
       understanding of certain properties of protons and neutrons, their own existence
       may perhaps ever be independently established. Here see have the antithesis of a
       nearly decomposable system"-it is a system which, if anything, is "nearly
       indecomposable", Yet what is curious is that a quark-based theory of rotors and
       neutrons (and other particles) has considerable explanatory power. in that many
       experimental results concerning the particles which narks supposedly compose
       can be accounted for quite well, quantitatively. by using the "quark model".

                   Superconductivity: A "Paradox" of Renormalization
       In Chapter V we discussed how renormalized particles emerge from their bare
       cores, by recursively compounded interactions with virtual particles. A
       renormalized particle can be seen either as this complex mathematical construct,
       or as the single lump which it is, physically. One of the strangest rid most
       dramatic consequences of this way of describing particles is the explanation it
       provides for the famous phenomenon of superconductivity resistance-free flow of
       electrons in certain solids, at extremely low temperatures.
               It turns out that electrons in solids are renormalized by their interactions
       with strange quanta of vibration called phonons (themselves renormalized as
       well!). These renormalized electrons are called polarons. Calculation shows that
       at very low temperatures, two oppositely spinning polarons sill begin to attract
       each other, and can actually become bound together in i certain way. Under the
       proper conditions. all the current-carrying polar




Levels of Description, and Computer Systems                                             310

       ons will get paired up, forming Cooper pains. Ironically, this pairing comes about
       precisely because electrons-the hare cores of the paired polarons--repel each other
       electrically. In contrast to the electrons, each Cooper pair feels neither attracted to
       nor repelled by an other Cooper pair, and consequently it can slip freely through a
       metal as if the metal were a vacuum. If you convert the mathematical description
       of such a metal from one whose primitive units are polarons into one whose
       primitive units are Cooper pairs. you get a considerable- simplified set of
       equations. This mathematical simplicity is the physicist's way of knowing that
       "chunking" into Cooper pairs is the natural way to look at superconductivity.
               Here we have several levels of particle: the Cooper pair itself: the two
       oppositely-spinning polarons which compose it: the electrons and phonons which
       make up the polarons: and then, within the electrons, the virtual photons and
       positrons, etc. etc. We can look at each level and perceive phenomena there,
       which are explained by an understanding of the levels below.

                                              "Sealing-off"
       Similarly, and fortunately. one does not have to know all about quarks to
       understand many things about the particles which they may compose. Thus, a
       nuclear physicist can proceed with theories of nuclei that are based on protons and
       neutrons, and ignore quark theories and their rivals. The nuclear physicist has a
       chunked picture of protons and neutrons-a description derived from lower-level
       theories buf which does not require understanding the lower-level theories.
       Likewise, an atomic physicist has a chunked picture of an atomic nucleus derived
       from nuclear theory. Then a chemist has a chunked picture of the electrons and
       their orbits, and builds theories of small molecules, theories which can be taken
       over in a chunked way by the molecular biologist, who has an intuition for how
       small molecules hang together, but whose technical expertise is in the field of
       extremely large molecules and how they interact. Then the cell biologist has a
       chunked picture of the units which the molecular biologist pores over, and tries to
       use them to account f'or the ways that cells interact. The point is clear. Each level
       is, in some sense, "sealed off'' from the levels below it. This is another of Simon's
       vivid terms, recalling the way in which a submarine is built in compartments, so
       that if one part is damaged, and water begins pouring in, the trouble can be
       prevented from spreading, by closing the doors, thereby sealing off the damaged
       compartment from neighboring compartments.
                Although there is always some "leakage" between the hierarchical levels
       of science, so that a chemist cannot afford to ignore lower-level physics totally, or
       a biologist to ignore chemistry totally, there is almost no leakage from one level
       to a distant level. That is why people earl, have intuitive understandings of other
       people without necessarily understanding the quark model, the structure of nuclei,
       the nature of electron orbits,




Levels of Description, and Computer Systems                                               311

       the chemical bond, the structure of proteins, the organelles in a cell, the methods
       of intercellular communication, the physiology 'of the various organs of the
       human body, or the complex interactions among organs. All at a person needs is a
       chunked model of how the highest level acts; and as all know, such models are
       very realistic and successful.

                     The Trade-off between Chunking and Determinism

       There is, however, perhaps one significant negative feature of a chunked model: it
       usually does not have exact predictive power. That is, we save ourselves from the
       impossible task of seeing people as collections of quarks (or whatever is at the
       lowest level) by using chunked models: but of course such models only give us
       probabilistic estimates of how other people feel, wil1 react to what we say or do,
       and so on. In short, in using chunked high-level models, we sacrifice determinism
       for simplicity. Despite not being sure how people will react to a joke, we tell it
       with the expectation at they will do something such as laugh, or not laugh-rather
       than, say, climb the nearest flagpole. (Zen masters might well do the latter!) A
       chunked model defines a "space" within which behavior is expected to fall, and
       specifies probabilities of its falling in different parts of that space.

                   "Computers Can Only Do What You Tell Them to Do"

               Now these ideas can be applied as well to computer programs as to
       compose physical systems. There is an old saw which says, "Computers can only
       what you tell them to do." This is right in one sense, but it misses the hint: you
       don't know in advance the consequences of what you tell a computer to do;
       therefore its behavior can be as baffling and surprising id unpredictable to you as
       that of a person. You generally know in advance the space in which the output
       will fall, but you don't know details of here it will fall. For instance, you might
       write a program to calculate the first million digits of 7r. Your program will spew
       forth digits of 7r much faster than you can-but there is no paradox in the fact that
       the computer outracing its programmer. You know in advance the space in which
       the output will lie-namely the space of digits between 0 and 9-which is to say, )u
       have a chunked model of the program's behavior; but if you'd known ie rest, you
       wouldn't have written the program.
               There is another sense in which this old saw is rusty. This involves the ct
       that as you program in ever higher-level languages, you know less and ss
       precisely what you've told the computer to do! Layers and layers of translation
       may separate the "front end" of a complex program from the actual machine
       language instructions. At the level you think and program, your statements may
       resemble declaratives and suggestions more than they resemble imperatives or
       commands. And all the internal rumbling provoked by the input of a high-level
       statement is invisible to you, generally, just as when you eat a sandwich, you are
       spared conscious awareness of the digestive processes it triggers




Levels of Description, and Computer Systems                                            312

               In any case, this notion that "computers can only do what they are told to
       do," first propounded by Lady Lovelace in her famous memoir, is so prevalent
       and so connected with the notion that "computers cannot think" that we shall
       return to it in later Chapters when our level of sophistication is greater.

                                      Two Types of System

       There is an important division between two types of system built up from many
       parts. There are those systems in which the behavior of some parts tends to cancel
       out the behavior of other parts, with the result that it does not matter too much
       what happens on the low level, because most anything will yield similar high-
       level behavior. An example of this kind of system is a container of gas, where all
       the molecules bump and bang against each other in very complex microscopic
       ways; but the total outcome, from a macroscopic point of view, is a very calm,
       stable system with a certain temperature, pressure, and volume. Then there are
       systems where the effect of a single low-level event may get magnified into an
       enormous high-level consequence. Such a system is a pinball machine, where the
       exact angle with which a ball strikes each post is crucial in determining the rest of
       its descending pathway.
                A computer is an elaborate combination of these two types of system. It
       contains subunits such as wires, which behave in a highly predictable fashion:
       they conduct electricity according to Ohm's law, a very precise, chunked law
       which resembles the laws governing gases in containers, since it depends on
       statistical effects in which billions of random effects cancel each other out,
       yielding a predictable overall behavior. A computer also contains macroscopic
       subunits, such as a printer, whose behavior is completely determined by delicate
       patterns of currents. What the printer prints is not by any means created by a
       myriad canceling microscopic effects. In fact, in the case of most computer
       programs, the value of every single bit in the program plays a critical role in the
       output that gets printed. If any bit were changed, the output would also change
       drastically.
                Systems which are made up of "reliable" subsystems only-that is,
       subsystems whose behavior can be reliably predicted from chunked descriptions-
       play inestimably important roles in our daily lives, because they are pillars of
       stability. We can rely on walls not to fall down, on sidewalks to go where they
       went yesterday, on the sun to shine, on clocks to tell the time correctly, and so on.
       Chunked models of such systems are virtually entirely deterministic. Of course,
       the other kind of system which plays a very large role in our lives is a system that
       has variable behavior which depends on some internal microscopic parameters-
       often a very large number of them, moreover-which we cannot directly observe.
       Our chunked model of such a system is necessarily in terms of the "space" of
       operation, and involves probabilistic estimates of landing in different regions of
       that space.
                A container of gas, which, as I already pointed out, is a reliable system



Levels of Description, and Computer Systems                                             313

       Because of many canceling effects, obeys precise, deterministic laws of physics.
       Such laws are chunked laws, in that they deal with the gas as a whole, nd ignore
       its constituents. Furthermore, the microscopic and macroscopic descriptions of a
       gas use entirely different terms. The former requires the pacification of the
       position and velocity of every single component molecule; the latter requires only
       the specification of three new quantities: temperature, pressure, and volume, the
       first two of which do not even have microscopic counterparts. The simple
       mathematical relationship which elates these three parameters- pV = cT, where c
       is a constant-is a law which depends on, yet is independent of, the lower-level
       phenomena. Less paradoxically, this law can be derived from the laws governing
       the molecular level; in that sense it depends on the lower level. On the other hand,
       it is law which allows you to ignore the lower level completely, if you wish: in hat
       sense it is independent of the lower level.
                It is important to realize that the high-level law cannot be stated in the
       vocabulary of the low-level description. "Pressure" and "temperature" are new
       terms which experience with the low level alone cannot convey. We humans
       perceive temperature and pressure directly; that is how we are guilt, so that it is
       not amazing that we should have found this law. But creatures which knew gases
       only as theoretical mathematical constructs would have to have an ability to
       synthesize new concepts, if they were to discover this law.

                                         Epiphenomena

       In drawing this Chapter to a close, I would like to relate a story about a complex
       system. I was talking one day with two systems programmers for he computer I
       was using. They mentioned that the operating system seemed to be able to handle
       up to about thirty-five users with great comfort, but at about thirty-five users or
       so, the response time all of a sudden hot up, getting so slow that you might as well
       log off and go home and wait until later. Jokingly I said, "Well, that's simple to fix
       just find the place in he operating system where the number `35' is stored, and
       change it to 60'!" Everyone laughed. The point is, of course, that there is no such
       place. where, then, does the critical number-35 users-come from? The answer is:
       It is a visible consequence of the overall system organization-an
       "epiphenometon,,.
               Similarly, you might ask about a sprinter, "Where is the `9.3' stored, hat
       makes him be able to run 100 yards in 9.3 seconds?" Obviously, it is not stored
       anywhere. His time is a result of how he is built, what his reaction time is, a
       million factors all interacting when he runs. The time is quite 'reproducible, but it
       is not stored in his body anywhere. It is spread around among all the cells of his
       body and only manifests itself in the act of the print itself.
               Epiphenomena abound. In the game of "Go", there is the feature that “two
       eyes live”. It is not built into the rules, but it is a consequence of the




Levels of Description, and Computer Systems                                              314

       rules. In the human brain, there is gullibility. How gullible are you? Is your
       gullibility located in some "gullibility center" in your brain? Could a
       neurosurgeon reach in and perform some delicate operation to lower your
       gullibility, otherwise leaving you alone? If you believe this, you are pretty
       gullible, and should perhaps consider such an operation.

                                         Mind vs. Brain

       In coming Chapters, where we discuss the brain, we shall examine whether the
       brain's top level-the mind-can be understood without understanding the lower
       levels on which it both depends and does not depend. Are there laws of thinking
       which are "sealed off" from the lower laws that govern the microscopic activity in
       the cells of the brain? Can mind be "skimmed" off of brain and transplanted into
       other systems? Or is it impossible to unravel thinking processes into neat and
       modular subsystems? Is the brain more like an atom, a renormalized electron, a
       nucleus, a neutron, or a quark? Is consciousness an epiphenomenon? To
       understand the mind, must one go all the way down to the level of nerve cells?




Levels of Description, and Computer Systems                                          315

FIGURE 60. [Drawing by the author.)

                                    ... Ant Fugue
        . . . . then, one by one, the four voices of the fugue chime in.)
Achilles: I know the rest of you won't believe this, but the answer to the question is
   staring us all in the face, hidden in the picture. It is simply one word-but what an
   important one: "MU"!
CCrab: I know the rest of you won't believe this, but the answer to the question is staring
   us all in the face, hidden in the picture. It is simply one word-but what an important
   one: "HOLISM"!
Achilles: Now hold on a minute. You must be seeing things. It's plain as day that the
   message of this picture is "MU", not "HOLISM"!
Crab: I beg your pardon, but my eyesight is extremely good. Please look again, and then
   tell me if the the picture doesn't say what I said it says!
Anteater: I know the rest of you won't believe this, but the answer to the question is
   staring us all in the face, hidden in the picture. It is simply one word-but what an
   important one: "REDUCTIONISM"!
Crab: Now hold on a minute. You must be seeing things. It's plain as day that the
   message of this picture is "HOLISM", not "REDUCTIONISM"!
Achilles: Another deluded one! Not "HOLISM", not "REDUCTIONISM", but "MU" is
   the message of this picture, and that much is certain.
Anteater: I beg your pardon, but my eyesight is extremely clear. Please look again, and
   then see if the picture doesn't say what I said it says.
Achilles: Don't you see that the picture is composed of two pieces, and that each of them
   is a single letter?
Crab: You are right about the two pieces, but you are wrong in your identification of
   what they are. The piece on the left is entirely composed of three copies of one word:
   "HOLISM"; and the piece on the right is composed of many copies, in smaller letters,
   of the same word. Why the letters are of different sizes in the two parts, I don't know,
   but I know what I see, and what I see is "HOLISM", plain as day. How you see
   anything else is beyond me.
Anteater: You are right about the two pieces, but you are wrong in your identification of
   what they are. The piece on the left is entirely composed of many copies of one
   word: "REDUCTIONISM"; and the piece on the right is composed of one single
   copy, in larger letters, of the same word. Why the letters are of different sizes in the
   two parts, I don't know, but I know what I see, and what I see is
   "REDUCTIONISM", plain as day. How you see anything else is beyond me.
Achilles: I know what is going on here. Each of you has seen letters which compose, or
   are composed of, other letters. In the left-hand piece,




Ant Fugue                                                                              317

    there are indeed three "HOLISM"'s, but each one of them is composed out of smaller
    copies of the word "REDUCTIONISM". And in complementary fashion, in the right-
    hand piece, there is indeed one "REDUCTIONISM", but it is composed out of
    smaller copies of the word "HOLISM". Now this is all fine and good, but in your
    silly squabble, the two of you have actually missed the forest for the trees. You see,
    what good is it to argue about whether "HOLISM" or "REDUCTIONISM" is right,
    when the proper way to understand the matter is to transcend the question, by
    answering "Mu",
Crab: I now see the picture as you have described it, Achilles, but I have no idea of what
    you mean by the strange expression "transcending the question".
Anteater: I now see the picture as you have described it, Achilles, but I have no idea of
    what you mean by the strange expression "Mu". .illes: I will be glad to indulge both
    of you, if you will first oblige me, by telling me the meaning of these strange
    expressions, "HOLISM" and "REDUCTIONISM".
Crab: HOLISM is the most natural thing in the world to grasp. It's simply the belief that
    "the whole is greater than the sum of its parts". No one in his right mind could reject
    holism.
Anteater: REDUCTIONISM is the most natural thing in the world to grasp. It's simply
    the belief that "a whole can be understood completely if you understand its parts, and
    the nature of their 'sum'". No one in her left brain could reject reductionism.
Crab: I reject reductionism. I challenge you to tell me, for instance, how to understand a
    brain reductionistically. Any reductionistic explanation of a brain will inevitably fall
    far short of explaining where the consciousness experienced by a brain arises from.
Anteater: I reject holism. I challenge you to tell me, for instance, how a holistic
    description of an ant colony sheds any more light on it than is shed by a description
    of the ants inside it, and their roles, and their interrelationships. Any holistic
    explanation of an ant colony will inevitably fall far short of explaining where the
    consciousness experienced by an ant colony arises from.
Antilles: Oh, no! The last thing which I wanted to do was to provoke another argument.
    Anyway, now that I understand the controversy, I believe that my explanation of
    "Mu" will help greatly. You see, "Mu" is an ancient Zen answer which, when given
    to a question, UNASKS the question. Here, the question seems to be, "Should the
    world be understood via holism, or via reductionism?" And the answer of "Mu" here
    rejects the premises of the question, which are that one or the other must be chosen.
    By unasking the question, it reveals a wider truth: that there is a larger context into
    which both holistic and reductionistic explanations fit.
Anteater: Absurd! Your "Mu" is as silly as a cow's moo. I'll have none of this Zen washy-
    wishiness.




Ant Fugue                                                                               318

Crab: Ridiculous! Your "ML" is as silly as a kitten's mew. I'll have none of this Zen
    washy-wishiness.
Achilles: Oh, dear! We're getting nowhere fast. Why have you stayed so strangely silent,
    Mr. Tortoise? It makes me very uneasy. Surely you must somehow be capable of
    helping straighten out this mess?
Tortoise: I know the rest of you won't believe this, but the answer to the question is
    staring us all in the face, hidden in the picture. It is simply one word-but what an
    important one: "Mu"!
Gust as he says this, the fourth voice in the fugue being played enters, exactly one octave
    below the first entry.)
Achilles: Oh, Mr. T, for once you have let me down. I was sure that you, who always see
    the most deeply into things, would be able to resolve this dilemma-but apparently,
    you have seen no further than I myself saw. Oh, well, I guess I should feel pleased to
    have seen as far as Mr. Tortoise, for once.
Tortoise: I beg your pardon, but my eyesight is extremely fine. Please look again, and
    then tell me if the picture doesn't say what I said it says.
Achilles: But of course it does! You have merely repeated my own original observation.
Tortoise: Perhaps "Mu" exists in this picture on a deeper level than you imagine,
    Achilles-an octave lower (figuratively speaking). But for now I doubt that we can
    settle the dispute in the abstract. I would like to see both the holistic and
    reductionistic points of view laid out more explicitly; then there may be more of a
    basis for a decision. I would very much like to hear a reductionistic description of an
    ant colony, for instance.
Crab: Perhaps Dr. Anteater will tell you something of his experiences in that regard.
    After all, he is by profession something of an expert on that subject.
Tortoise: I am sure that we have much to learn from you, Dr. Anteater. Could you tell us
    more about ant colonies, from a reductionistic point of view?
Anteater: Gladly. As Mr. Crab mentioned to you, my profession has led me quite a long
    way into the understanding of ant colonies.
Achilles: I can imagine! The profession of anteater would seem to be synonymous with
    being an expert on ant colonies!
Anteater: I beg your pardon. "Anteater" is not my profession; it is my species. By
    profession, I am a colony surgeon. I specialize in correcting nervous disorders of the
    colony by the technique of surgical removal.
Achilles: Oh, I see. But what do you mean by "nervous disorders" of an ant colony?
Anteater: Most of my clients suffer from some sort of speech impairment. You know,
    colonies which have to grope for words in everyday situations. It can be quite tragic.
    I attempt to remedy the situation by, uhh—removing the defective part of the colony.
    These operations




Ant Fugue                                                                              319

    are sometimes quite involved, and of course years of study are required before one
    can perform them.
Achilles: But-isn't it true that, before one can suffer from speech impairment, one must
    have the faculty of speech?
Anteater: Right.
Achilles: Since ant colonies don't have that faculty, I am a little confused. Crab: It's too
    bad, Achilles, that you weren't here last week, when Dr.
Anteater and Aunt Hillary were my house guests. I should have thought of having you
    over then.
Achilles: Is Aunt Hillary your aunt, Mr. Crab? Crab: Oh, no, she's not really anybody's
    aunt.
Anteater: But the poor dear insists that everybody should call her that, even strangers. It's
    just one of her many endearing quirks.
Crab: Yes, Aunt Hillary is quite eccentric, but such a merry old soul. It's a shame I didn't
    have you over to meet her last week.
Anteater: She's certainly one of the best-educated ant colonies I have ever had the good
    fortune to know. The two of us have spent many a long evening in conversation on
    the widest range of topics.
Achilles: I thought anteaters were devourers of ants, not patrons of antintellectualism!
Anteater: Well, of course the two are not mutually inconsistent. I am on the best of terms
    with ant colonies. It's just ANTS that I eat, not colonies-and that is good for both
    parties: me, and the colony.
Achilles: How is it possible that--
Tortoise: How is it possible that--
Achilles: -having its ants eaten can do an ant colony any good? Crab: How is it possible
    that
Tortoise: -having a forest fire can do a forest any good? Anteater: How is it possible that
Crab: -having its branches pruned can do a tree any good? Anteater: -having a haircut can
    do Achilles any good?
Tortoise: Probably the rest of you were too engrossed in the discussion to notice the
    lovely stretto which just occurred in this Bach fugue.
Achilles: What is a stretto?
Tortoise: Oh, I'm sorry; I thought you knew the term. It is where one theme repeatedly
    enters in one voice after another, with very little delay between entries.
Achilles: If I listen to enough fugues, soon I'll know all of these things and will be able to
    pick them out myself, without their having to be pointed out.
Tortoise: Pardon me, my friends. I am sorry to have interrupted. Dr. Anteater was trying
    to explain how eating ants is perfectly consistent with being a friend of an ant colony.
Achilles: Well, I can vaguely see how it might be possible for a limited and regulated
    amount of ant consumption to improve the overall health of




Ant Fugue                                                                                 320

    a colony-but what is far more perplexing is all this talk about having conversations
    with ant colonies. That's impossible. An ant colony is simply a bunch of individual
    ants running around at random looking for food and making a nest.
Anteater: You could put it that way if you want to insist on seeing the trees but missing
    the forest, Achilles. In fact, ant colonies, seen as wholes, are quite well-defined units,
    with their own qualities, at times including the mastery of language.
Achilles: I find it hard to imagine myself shouting something out loud in the middle of
    the forest, and hearing an ant colony answer back.
Anteater: Silly fellow! That's not the way it happens. Ant colonies don't converse out
    loud, but in writing. You know how ants form trails leading them hither and thither?
Achilles: Oh, yes-usually straight through the kitchen sink and into my peach jam.
Anteater: Actually, some trails contain information in coded form. If you know the
    system, you can read what they're saying just like a book. Achilles: Remarkable. And
    can you communicate back to them? Anteater: Without any trouble at all. That's how
    Aunt Hillary and I have conversations for hours. I take a stick and draw trails in the
    moist ground, and watch the ants follow my trails. Presently, a new trail starts getting
    formed somewhere. I greatly enjoy watching trails develop. As they are forming, I
    anticipate how they will continue (and more often I am wrong than right). When the
    trail is completed, I know what Aunt Hillary is thinking, and I in turn make my reply.
Achilles: There must be some amazingly smart ants in that colony, I'll say that.
Anteater: I think you are still having some difficulty realizing the difference in levels
    here. Just as you would never confuse an individual tree with a forest, so here you
    must not take an ant for the colony. You see, all the ants in Aunt Hillary are as dumb
    as can be. They couldn't converse to save their little thoraxes!
Achilles: Well then, where does the ability to converse come from? It must reside
    somewhere inside the colony! I don't understand how the ants can all be unintelligent,
    if Aunt Hillary can entertain you for hours with witty banter.
Tortoise: It seems to me that the situation is not unlike the composition of a human brain
    out of neurons. Certainly no one would insist that individual brain cells have to be
    intelligent beings on their own, in order to explain the fact that a person can have an
    intelligent conversation.
Achilles: Oh, no, clearly not. With brain cells, I see your point completely. Only ... ants
    are a horse of another color. I mean, ants just roam about at will, completely
    randomly, chancing now and then upon a morsel of food ... They are free to do what
    they want to do, and with that freedom, I don’t see at all how their behaviour, seen as
    a whole, can




Ant Fugue                                                                                 321

   amount to anything coherent-especially something so coherent as the brain behavior
   necessary for conversing.
Crab: It seems to me that the ants are free only within certain constraints. For example,
   they are free to wander, to brush against each other, to pick up small items, to work
   on trails, and so on. But they never step out of that small world, that ant-system,
   which they are in. It would never occur to them, for they don't have the mentality to
   imagine anything of the kind. Thus the ants are very reliable components, in the
   sense that you can depend on them to perform certain kinds of tasks in certain ways.
Achilles: But even so, within those limits they are still free, and they just act at random,
   running about incoherently without any regard for the thought mechanisms of a
   higher-level being which Dr. Anteater asserts they are merely components of.
Anteater: Ah, but you fail to recognize one thing. Achilles-the regularity of statistics.
Achilles: How is that?
Anteater: For example, even though ants as individuals wander about in what seems a
   random way, there are nevertheless overall trends, involving large numbers of ants,
   which can emerge from that chaos.
Achilles: Oh, I know what you mean. In fact, ant trails are a perfect example of such a
   phenomenon. There, you have really quite unpredictable motion on the part of any
   single ant-and yet, the trail itself seems to remain well-defined and stable. Certainly
   that must mean that the individual ants are not just running about totally at random.
Anteater: Exactly, Achilles. There is some degree of communication among the ants, just
   enough to keep them from wandering off completely at random. By this minimal
   communication they can remind each other that they are not alone but are
   cooperating with teammates. It takes a large number of ants, all reinforcing each
   other this way, to sustain any activity-such as trail-building-for any length of time.
   Now my very hazy understanding of the operation of brains leads me to believe that
   something similar pertains to the firing of neurons. Isn't it true, Mr. Crab, that it takes
   a group of neurons firing in order to make another neuron fire?
Crab: Definitely. Take the neurons in Achilles' brain, for example. Each neuron receives
   signals from neurons attached to its input lines, and if the sum total of inputs at any
   moment exceeds a critical threshold. then that neuron will fire and send its own
   output pulse rushing off to other neurons, which may in turn fire-and on down the
   line it goes. The neural flash swoops relentlessly in its Achillean path, in shapes
   stranger then the dash of a gnat-hungry swallow; every twist, every turn foreordained
   by the neural structure in Achilles' brain, until sensory input messages interfere.
Achilles: Normally, I think that I'M in control of what I think-but the way you put it turns
   it all inside out, so that it sounds as though "I" am just




Ant Fugue                                                                                 322

    what comes out of all this neural structure, and natural law. It makes what I consider
    my SELF sound at best like a by-product of an organism governed by natural law,
    and at worst, an artificial notion produced by my distorted perspective. In other
    words, you make me feel like I don't know who or what-I am, if anything.
Tortoise: You'll come to understand much better as we go along. But Dr.
Anteater-what do you make of this similarity?
Anteater: I knew there was something parallel going on in the two very different systems.
    Now I understand it much better. It seems that group phenomena which have
    coherence-trail-building, for example-will take place only when a certain threshold
    number of ants get involved. If an effort is initiated, perhaps at random, by a few ants
    in some locale, one of two things can happen: either it will fizzle out after a brief
    sputtering start
Achilles: When there aren't enough ants to keep the thing rolling?
Anteater: Exactly. The other thing that can happen is that a critical mass of ants is
    present, and the thing will snowball, bringing more and more ants into the picture. In
    the latter case, a whole "team" is brought into being which works on a single project.
    That project might be trailmaking, or food-gathering, or it might involve nest-
    keeping. Despite the extreme simplicity of this scheme on a small scale, it can give
    rise to very complex consequences on a larger scale.
Achilles: I can grasp the general idea of order emerging from chaos, as you sketch it, but
    that still is a long way from the ability to converse. After all, order also emerges from
    chaos when molecules of a gas bounce against each other randomly-yet all that
    results there is an amorphous mass with but three parameters to characterize it:
    volume, pressure, and temperature. Now that's a far cry from the ability to understand
    the world, or to talk about it!
Anteater: That highlights a very interesting difference between the explanation of the
    behavior of an ant colony and the explanation of the behavior of gas inside a
    container. One can explain the behavior of the gas simply by calculating the
    statistical properties of the motions of its molecules. There is no need to discuss any
    higher elements of structure than molecules, except the full gas itself. On the other
    hand, in an ant colony, you can't even begin to understand the activities of the colony
    unless you go through several layers of structure.
Achilles: I see what you mean. In a gas, one jump takes you from the lowest level-
    molecules-to the highest level-the full gas. There are no intermediate levels of
    organization. Now how do intermediate levels of organized activity arise in an ant
    colony?
Anteater: It has to do with the existence of several different varieties of ants inside any
    colony.
Achilles: Oh, yes. I think I have heard about that. They are called "castes", aren't they?
Anteater: That's correct. Aside from the queen, there are males, who do




Ant Fugue                                                                                323

    practically nothing towards, the upkeep of the nest, and then—
Achilles: And of course there are soldiers-Glorious Fighters Against Communism!
Crab: Hmm ... I hardly think that could be right, Achilles. An ant colony is quite
    communistic internally, so why would its soldiers fight against communism? Or am I
    right, Dr. Anteater? .
Anteater: Yes, about colonies you are right, Mr. Crab; they are indeed based on
    somewhat communistic principles. But about soldiers Achilles is somewhat naive. In
    fact, the so-called "soldiers" are hardly adept at fighting at all. They are slow,
    ungainly ants with giant heads, who can snap with their strong jaws, but are hardly to
    be glorified. As in a true communistic state, it is rather the workers who are to be
    glorified. It is they who do most of the chores, such as food-gathering, hunting, and
    nursing of the young. It is even they who do most of the fighting.
Achilles: Bah. That is an absurd state of affairs. Soldiers who won't fight!
Anteater: Well, as I just said, they really aren't soldiers at all. It's the workers who are
    soldiers; the soldiers are just lazy fatheads.
Achilles: Oh, how disgraceful! Why, if I were an ant, I'd put some discipline in their
    ranks! I'd knock some sense into those fatheads!
Tortoise: If you were an ant? How could you be an ant? There is no way to map your
    brain onto an ant brain, so it seems to me to be a pretty fruitless question to worry
    over. More reasonable would be the proposition of mapping your brain onto an ant
    colony ... But let us not get sidetracked. Let Dr. Anteater continue with his most
    illuminating description of castes and their role in the higher levels of organization.
Anteater: Very well. There are all sorts of tasks which must be accomplished in a colony,
    and individual ants develop specializations. Usually an ant's specialization changes as
    the ant ages. And of course it is also dependent on the ant's caste. At any one
    moment, in any small area of a colony, there are ants of all types present. Of course,
    one caste may be be very sparse in some places and very dense in others.
Crab: Is the density of a given caste, or specialization, just a random thing? Or is there a
    reason why ants of one type might be more heavily concentrated in certain areas, and
    less heavily in others?
Anteater: I'm glad you brought that up, since it is of crucial importance in understanding
    how a colony thinks. In fact, there evolves, over a long period of time, a very delicate
    distribution of castes inside a colony. And it is this distribution which allows the
    colony to have the complexity which underlies the ability to converse with me.
Achilles: It would seem to me that the constant motion of ants to and fro would
    completely prevent the possibility of a very delicate distribution.Any delicate
    distribution would be quickly destroyed by all the random motions of ants, just as
    any delicate pattern among molecules in a gas would not survive for an instant, due
    to the random bombardment from all sides.
Anteater: In an ant colony. the situation is quite the contrary. In fact, it is just exactly the
    to-ing and fro-ing of ants inside the colony




Ant Fugue                                                                                  324

    which adapts the caste distribution to varying situations, and thereby preserves the
    delicate caste distribution. You see, the caste distribution cannot remain as one single
    rigid pattern; rather, it must constantly be changing so as to reflect, in some manner,
    the real-world situation with which the colony is dealing, and it is precisely the
    motion inside the colony which updates the caste distribution, so as to keep it in line
    with the present circumstances facing the colony.
Tortoise: Could you give an example?
Anteater: Gladly. When I, an anteater, arrive to pay a visit to Aunt Hillary, all the foolish
    ants, upon sniffing my odor, go into a panic-which means, of course, that they begin
    running around completely differently from the way they were before I arrived.
Achilles: But that's understandable, since you're a dreaded enemy of the colony.
Anteater: Oh, no. I must reiterate that, far from being an enemy of the colony, I am Aunt
    Hillary's favorite companion. And Aunt Hillary is my favorite aunt. I grant you, I'm
    quite feared by all the individual ants in the colony-but that's another matter entirely.
    In any case, you see that the ants' action in response to my arrival completely changes
    the internal distribution of ants.
Achilles: That's clear.
Anteater: And that sort of thing is the updating which I spoke of. The new distribution
    reflects my presence. One can describe the change from old state to new as having
    added a "piece of knowledge" to the colony.
Achilles: How can you refer to the distribution of different types of ants inside a colony
    as a "piece of knowledge"?
Anteater: Now there's a vital point. It requires some elaboration. You see, what it comes
    down to is how you choose to describe the caste distribution. If you continue to think
    in terms of the lower levels-individual ants-then you miss the forest for the trees.
    That's just too microscopic a level, and when you think microscopically, you're
    bound to miss some large-scale features. You've got to find the proper high-level
    framework in which to describe the caste distribution-only then will it make sense
    how the caste distribution can encode many pieces of knowledge.
Achilles: Well, how DO you find the proper-sized units in which to describe the present
    state of the colony, then?
Anteater: All right. Let's begin at the bottom. When ants need to get something done,
    they form little "teams", which stick together to perform a chore. As I mentioned
    earlier, small groups of ants are constantly forming and unforming. Those which
    actually exist for a while are the teams, and the reason they don't fall apart is that
    there really is something for them to do.
Achilles: Earlier you said that a group will stick together if its size exceeds a certain
    threshold. Now you're saying that a group will stick together if there is something for
    it to do.
Anteater: They are equivalent statements. For instance, in food-gathering,




Ant Fugue                                                                                325

    if there is an inconsequential amount of food somewhere which gets discovered by
    some wandering Ant who then attempts to communicate its enthusiasm to other ants,
    the number of ants who respond will be proportional to the size of the food sample-
    and an inconsequential amount will not attract enough ants to surpass the threshold.
    Which is exactly what I meant by saying there is nothing to do-too little food ought
    to be ignored.
Achilles: I see. I assume that these "teams" are one of the levels of structure falling
    somewhere in between the single-ant level and the colony level.
Anteater: Precisely. There exists a special kind of team, which I call a "signal"-and all the
    higher levels of structure are based on signals. In fact, all the higher entities are
    collections of signals acting in concert. There are teams on higher levels whose
    members are not ants, but teams on lower levels. Eventually you reach the lowest-
    level teams which is to say, signals-and below them, ants.
Achilles: Why do signals deserve their suggestive name?
Anteater: It comes from their function. The effect of signals is to transport ants of various
    specializations to appropriate parts of the colony. So the typical story of a signal is
    thus: it comes into existence by exceeding the threshold needed for survival, then it
    migrates for some distance through the colony, and at some point it more or less
    disintegrates into its individual members, leaving them on their own.
Achilles: It sounds like a wave, carrying sand dollars and seaweed from afar, and leaving
    them strewn, high and dry, on the shore.
4nteater: In a way that's analogous, since the team does indeed deposit something which
    it has carried from a distance, but whereas the water in the wave rolls back to the sea,
    there is no analogous carrier substance in the case of a signal, since the ants
    themselves compose it.
Tortoise: And I suppose that a signal loses its coherency just at some spot in the colony
    where ants of that type were needed in the first place.
Anteater: Naturally.
Achilles: Naturally? It's not so obvious to ME that a signal should always go just where it
    is needed. And even if it goes in the right direction, how does it figure out where to
    decompose? How does it know it has arrived?
Anteater: Those are extremely important matters, since they involve explaining the
    existence of purposeful behavior-or what seems to be purposeful behavior-on the part
    of signals. From the description, one would be inclined to characterize the signals'
    behavior as being oriented towards filling a need, and to call it "purposeful". But you
    can look at it otherwise.
Achilles: Oh, wait. Either the behavior is purposeful, or it is NOT. I don't see how you
    can have it both ways.
Anteater: Let me explain my way of seeing things, and then see if you agree. Once a
    signal is formed, there is no awareness on its part that it




Ant Fugue                                                                                326

   should head off in any particular direction. But here, the delicate caste distribution
   plays a crucial role. It is what determines the motion of signals through the colony,
   and also how long a signal will remain stable, and where it will "dissolve".
Achilles: So everything depends on the caste distribution, eh?
Anteater: Right. Let's say a signal is moving along. As it goes, the ants which compose it
   interact, either by direct contact or by exchange of scents, with ants of the local
   neighborhoods which it passes through. The contacts and scents provide information
   about local matters of urgency, such as nest-building, or nursing, or whatever. The
   signal will remain glued together as long as the local needs are different from what it
   can supply; but if it CAN contribute, it disintegrates, spilling a fresh team of usable
   ants onto the scene. Do you see now how the caste distribution acts as an overall
   guide of the teams inside the colony?
Achilles: I do see that.
Anteater: And do you see how this way of looking at things requires attributing no sense
   of purpose to the signal?
Achilles: I think so. Actually, I'm beginning to see things from two different vantage
   points. From an ant's-eye point of view, a signal has NO purpose. The typical ant in a
   signal is just meandering around the colony, in search of nothing in particular, until it
   finds that it feels like stopping. Its teammates usually agree, and at that moment the
   team unloads itself by crumbling apart, leaving just its members but none of its
   coherency. No planning is required, no looking ahead; nor is any search required, to
   determine the proper direction. But from the COLONY'S point-of view, the team has
   just responded to a message which was written in the language of the caste
   distribution. Now from this perspective, it looks very much like purposeful activity.
Crab: What would happen if the caste distribution were entirely random? Would signals
   still band and disband?
Anteater: Certainly. But the colony would not last long, due to the meaninglessness of the
   caste distribution.
Crab: -Precisely the point I wanted to make. Colonies survive because their caste
   distribution has meaning, and that meaning is a holistic aspect, invisible on lower
   levels. You lose explanatory power unless you take that higher level into account.
Anteater: I see your side; but I believe you see things too narrowly.
Crab: How so?
Anteater: Ant colonies have been subjected to the rigors of evolution for billions of years.
   A few mechanisms were selected for, and most were selected against. The end result
   was a set of mechanisms which make ant colonies work as we have been describing.
   If you could watch the whole process in a movie-running a billion or so times faster
   than life, of course-the emergence of various mechanisms would be seen as natural
   responses to external pressures, just as bubbles in boiling water are natural responses
   to an external heat source. I don't suppose you




Ant Fugue                                                                               327

    see "meaning" and "purpose", in the bubbles in boiling water-or do you?
Crab: No, but
Anteater: Now that's MY point. No matter how big a bubble is, it owes its existence to
    processes on the molecular level, and you can forget about any "higher-level laws".
    The same goes for ant colonies and their teams. By looking at things from the vast
    perspective of evolution, you can drain the whole colony of meaning and purpose.
    They become superfluous notions.
Achilles: Why, then, Dr. Anteater, did you tell me that you talked with Aunt Hillary? It
    now seems that you would deny that she can talk or think at all.
Anteater: I am not being inconsistent, Achilles. You see, I have as much difficulty as
    anyone else in seeing things on such a grandiose time scale, so I find it much easier
    to change points of view. When I do so, forgetting about evolution and seeing things
    in the here and now, the vocabulary of teleology comes back: the MEANING of the
    caste distribution and the PURPOSEFULNESS of signals. This not only happens
    when I think of ant colonies, but also when I think about my own brain and other
    brains. However, with some effort I can always remember the other point of view if
    necessary, and drain all these systems of meaning, too.
Crab: Evolution certainly works some miracles. You never know the next trick it will pull
    out of its sleeve. For instance, it wouldn't surprise me one bit if it were theoretically
    possible for two or more "signals" to pass through each other, each one unaware that
    the other one is also a signal; each one treating the other as if it were just part of the
    background population.
Anteater: It is better than theoretically possible; in fact it happens routinely!
Achilles: Hmm ... What a strange image that conjures up in my mind. I can just imagine
    ants moving in four different directions, some black, some white, criss-crossing,
    together forming an orderly pattern, almost like-like
Tortoise: A fugue, perhaps?
Achilles: Yes-that's it! An ant fugue!
Crab: An interesting image, Achilles. By the way, all that talk of boiling water made me
    think of tea. Who would like some more? Achilles: I could do with another cup, Mr.
    C.
Crab: Very good.
Achilles: Do you suppose one could separate out the different visual "voices" of such an
    "ant fugue"? I know how hard it is for me
Tortoise: Not for me, thank you.
Achilles: -to track a single voice
Anteater: I'd like some, too, Mr. Crab
Achilles: -- in a musical fugue--




Ant Fugue                                                                                 328

              FIGURE 61. "Ant Fugue", by M. C. Escher (woodcut, 1953).

Anteater: -if it isn't too much trouble
Achilles: . -when all of them
Crab: Not at all. Four cups of tea
Tortoise: Three?
Achilles: -are going at once.
Crab: -coming right up!
Anteater: That's an interesting thought, Achilles. But its unlikely that anyone could draw
    such a picture in a convincing way
Achilles: That's too bad.
Tortoise: Perhaps you could answer this, Dr. Anteater. Does a signal, from its creation
    until its dissolution, always consist of the same set of ants?
Anteater: As a matter of fact, the individuals in a signal sometimes break off and get
    replaced by others of the same caste, if there are a few in the area. Most often, signals
    arrive at their disintegration points with nary an ant in common with their starting
    lineup.




Ant Fugue                                                                                329

Crab: I can see that the signals are constantly affecting the caste distribution throughout
    the colony, and are doing so in response to the internal needs of the colony-which in
    turn reflect the external situation which the colony is faced with. Therefore the caste
    distribution, as you said, Dr. Anteater, gets continually updated in a way which
    ultimately reflects the outer world.
Achilles: But what about those intermediate levels of structure? You were saying that the
    caste distribution should best be pictured not in terms of ants or signals, but in terms
    of teams whose members were other teams, whose members were other teams, and
    so on until you come down to the ant level. And you said that that was the key to
    understanding how it was possible to describe the caste distribution as encoding
    pieces of information about the world.
Anteater: Yes, we are coming to all that. I prefer to give teams of a sufficiently high level
    the name of "symbols". Mind you, this sense of the word has some significant
    differences from the usual sense. My "symbols" are ACTIVE SUBSYSTEMS of a
    complex system, and they are composed of lower-level active subsystems ... They are
    therefore quite different from PASSIVE symbols, external to the system, such as
    letters of the alphabet or musical notes, which sit there immobile, waiting for an
    active system to process them.
Achilles: Oh, this is rather complicated, isn't it? I just had no idea that ant colonies had
    such an abstract structure.
Anteater: Yes, it's quite remarkable. But all these layers of structure are necessary for the
    storage of the kinds of knowledge which enable an organism to be "intelligent" in
    any reasonable sense of the word. Any system which has a mastery of language has
    essentially the same underlying sets of levels.
Achilles: Now just a cotton-picking minute. Are you insinuating that my brain consists
    of, at bottom, just a bunch of ants running around?
Anteater: Oh, hardly. You took me a little too literally. The lowest level may be utterly
    different. Indeed, the brains of anteaters, for instance, are not composed of ants. But
    when you go up a level or two in a brain, you reach a level whose elements have
    exact counterparts in other systems of equal intellectual strength-such as ant colonies.
Tortoise: That is why it would be reasonable to think of mapping your brain, Achilles,
    onto an ant colony, but not onto the brain of a mere ant.
Achilles: I appreciate the compliment. But how would such a mapping be carried out?
    For instance, what in my brain corresponds to the low level teams which you call
    signals?
Anteater: Oh, I but dabble in brains, and therefore couldn't set up the map in its glorious
    detail. But-and correct me if I'm wrong, Mr. Crab-I would surmise that the brain
    counterpart to an ant colony's signal is the firing of a neuron; or perhaps it is a larger-
    scale event, such as a pattern of neural firings.




Ant Fugue                                                                                 330

Crab: I would tend to agree. But don't you think that, for the purposes of our discussion,
   delineating the exact counterpart is not in itself crucial, desirable though it might be?
   It seems to me that the main idea is that such a correspondence does exist, even if we
   don't know exactly how to define it right now. I would only question one point, Dr.
   Anteater, which you raised, and that concerns the level at which one can have faith
   that the correspondence begins. You seemed to think that a SIGNAL might have a
   direct counterpart in a brain; whereas I feel that it is only at the level of your
   ACTIVE SYMBOLS and above that it is likely that a correspondence must exist.
Anteater: Your interpretation may very well be more accurate than mine, Mr. Crab.
   Thank you for bringing out that subtle point.
Achilles: What does a symbol do that a signal couldn't do?
Anteater: It is something like the difference between words and letters. Words, which are
   meaning-carrying entities, are composed of letters, which in themselves carry no
   meaning. This gives a good idea of the difference between symbols and signals. In
   fact it is a useful analogy, as long as you keep in mind the fact that words and letters
   are PASSIVE, symbols and signals are ACTIVE.
Achilles: I'll do so, but I'm not sure I understand why it is so vital to stress the difference
   between active and passive entities.
Anteater: The reason is that the meaning which you attribute to any passive symbol, such
   as a word on a page, actually derives from the meaning which is carried by
   corresponding active symbols in your brain. So that the meaning of passive symbols
   can only be properly understood when it is related to the meaning of active symbols.
Achilles: All right. But what is it that endows a SYMBOL-an active one, to be sure-with
   meaning, when you say that a SIGNAL, which is a perfectly good entity in its own
   right, has none? Anteater: It all has to do with the way that symbols can cause other
   symbols to be triggered. When one symbol becomes active, it does not do so in
   isolation. It is floating about, indeed, in a medium, which is characterized by its caste
   distribution.
Crab: Of course, in a brain there is no such thing as a caste distribution, but the
   counterpart is the "brain state". There, you describe the states of all the neurons, and
   all the interconnections, and the threshold for firing of each neuron.
Anteater: Very well; let's lump "caste distribution" and "brain state" under a common
   heading, and call them just the "state". Now the state can be described on a low level
   or on a high level. A low-level description of the state of an ant colony would involve
   painfully specifying the location of each ant, its age and caste, and other similar
   items. A very detailed description, yielding practically no global insight as to WHY it
   is in that state. On the other hand, a description on a high level would involve
   specifying which symbols could be triggered by which combinations of other
   symbols, under what conditions, and so forth.




Ant Fugue                                                                                  331

Achilles: What about a description on the level of signals, or teams?
Anteater: A description on that level would fall somewhere in between the low-level and
   symbol-level descriptions. It would contain a great deal of information about what is
   actually going on in specific locations throughout the colony, although certainly less
   than an ant-by-ant description, since teams consist of clumps of ants. A team-by-team
   description is like a summary of an ant-by-ant description. However, you have to add
   extra things which were not present in the ant-by-ant description-such as the
   relationships between teams, and the supply of various castes here and there. This
   extra complication is the price you pay for the right to summarize.
Achilles: It is interesting to me to compare the merits of the descriptions at various levels.
   The highest-level description seems to carry the most explanatory power, in that it
   gives you the most intuitive picture of the ant colony, although strangely enough, it
   leaves out seemingly- the most important feature-the ants.
Anteater: But you see, despite appearances, the ants are not the most important feature.
   Admittedly, were it not for them, the colony Wouldn't exist: but something
   equivalent-a brain-can exist, ant-free. So, at least from a high-level point of view, the
   ants are dispensable. .Achilles: I'm sure no ant would embrace your theory with
   eagerness.
Anteater: Well, I never met an ant with a high-level point of view.
Crab: What a counterintuitive picture you paint, Dr. Anteater. It seems that, if what you
   say is true, in order to grasp the whole structure, you have to describe it omitting any
   mention of its fundamental building blocks.
Anteater: Perhaps I can make it a little clearer by an analogy. Imagine you have before
   you a Charles Dickens novel.
Achilles: The Pickwick Papers-will that do?
Anteater: Excellently! And now imagine trying the following game: you must find a way
   of mapping letters onto ideas, so that the entire Pickwick Papers makes sense when
   you read it letter by letter.
Achilles: Hmm ... You mean that every time I hit a word such as "the", I have to think of
   three definite concepts, one after another, with no room for variation?
Anteater: Exactly. They are the `t'-concept, the `h'-concept, and the `e'-concept-and every
   time, those concepts are as they were the preceding time.
Achilles: Well, it sounds like that would turn the experience of "reading" The Pickwick
   Papers into an indescribably boring nightmare. It would be an exercise in
   meaninglessness, no matter what concept I associated with each letter.
Anteater: Exactly. There is no natural mapping from the individual letters into the real
   world. The natural mapping occurs on a higher level between words, and parts of the
   real world. If you wanted to describe the book, therefore, you would make no
   mention of the letter level.




Ant Fugue                                                                                 332

Achilles: Of course not! I'd describe the plot and the characters, and so forth.
Anteater: So there you are. You would omit all mention of the building blocks, even
    though the book exists thanks to them. They are the medium, but not the message.
Achilles: All right-but what about ant colonies?
Anteater: Here, there are active signals instead of passive letters, and active symbols
    instead of passive words-hut the idea carries over.
Achilles: Do you mean I couldn't establish a mapping between signals and things in the
    real world?
Anteater: You would find that you could not do it in such a way that the triggering of new
    signals would make am sense. Nor could you succeed on any lower level-for example
    the ant level. Only on the symbol level do the triggering patterns make sense.
    Imagine, for instance, that one day you were watching Aunt Hillary when I arrived to
    pay a call. You could watch as carefully as you wanted, and yet you would probably
    perceive nothing more than a rearrangement of ants.
Achilles: I'm sure that's accurate.
Anteater: And yet, as I watched, reading the higher level instead of the lower level, I
    would see several dormant symbols being awakened, those which translate into the
    thought, "Oh, here's that charming Dr. Anteater again-how pleasant!"-or words to
    that effect.
Achilles: That sounds like what happened when the four of us all found different levels to
    read in the MU-picture--or at least THREE of us did .. .
Tortoise: What an astonishing coincidence that there should be such a resemblance
    between that strange picture which I chanced upon in the Well-Tempered Clavier,
    and the trend of our conversation.
Achilles: Do you think it's just coincidence?
Tortoise: Of course.
Anteater: Well, I hope you can grasp now how the thoughts in Aunt Hillary emerge from
    the manipulation of symbols composed of signals composed of teams composed of
    lower-level teams, all the way down to ants.
Achilles: Why do you call it "symbol manipulation"? Who does the manipulating, if the
    symbols are themselves active? Who is the agent?
Anteater: This gets back to the question which you earlier raised about purpose. You're
    right that symbols themselves are active, but the activities which they follow are
    nevertheless not absolutely free. The activities of all symbols are strictly determined
    by the state of the full system in which they reside. Therefore, the full system is
    responsible for how its symbols trigger each other, and so it is quite reasonable to
    speak of the full system as the "agent". As the symbols operate, the state of the
    system gets slowly transformed, or updated. But there are many features which
    remain over time. It is this partially constant, partially varying system which is the
    agent. One can give a name to the




Ant Fugue                                                                              333

    full system. For example, Aunt Hillary is the "who" who can be said to manipulate
    her symbols; and you are similar, Achilles.
Achilles: That's quite a strange characterization of the notion of who I am. I'm not sure I
    can fully understand it, but I will give it some thought.
Tortoise: It would be quite interesting to follow the symbols in your brain as you do that
    thinking about the symbols in your brain.
Achilles: That's too complicated for me. I have trouble enough just trying to picture how
    it is possible to look at an ant colony and read it on the symbol level. I can certainly
    imagine perceiving it at the ant level; and with a little trouble, I can imagine what it
    must be like to perceive it at the signal level; but what in the world can it be like to
    perceive an ant colony at the symbol level?
Anteater: One only learns through long practice. But when one is at my stage, one reads
    the top level of an ant colony as easily as you yourself read the "MU" in the MU-
    picture.
Achilles: Really? That must be an amazing experience.
Anteater: In a way-but it is also one which is quite familiar to you, Achilles.
Achilles: Familiar to me? What do you mean? I have never looked at an ant colony on
    anything but the ant level.
Anteater: Maybe not; but ant colonies are no different from brains in many respects.
Achilles: I have never seen nor read any brain either, however.
Anteater: What about your OWN brain? Aren't you aware of your own thoughts? Isn't
    that the essence of consciousness? What else are you doing but reading your own
    brain directly at the symbol level?
Achilles: I never thought of it that way. You mean that I bypass all the lower levels, and
    only see the topmost level?
Anteater: That's the way it is, with conscious systems. They perceive themselves on the
    symbol level only, and have no awareness of the lower levels, such as the signal
    levels.
Achilles: Does it follow that in a brain, there are active symbols which are constantly
    updating themselves so that they reflect the overall state of the brain itself, always on
    the symbol level?
Anteater: Certainly. In any conscious system there are symbols which represent the brain
    state, and they are themselves part of the very brain state which they symbolize. For
    consciousness requires a large degree of self-consciousness.
Achilles: That is a weird notion. It means that although there is frantic activity occurring
    in my brain at all times, I am only capable of registering that activity in one way-on
    the symbol level; and I am completely insensitive to the lower levels. It is like being
    able to read a Dickens novel by direct visual perception, without ever having learned
    the letters of the alphabet. I can't imagine anything as weird as that really happening.
Crab: But precisely that sort of thing can happen when you read “MU”,




Ant Fugue                                                                                334

    without perceiving the lower levels "HOLISM" and "REDUCTIONISM".
Achilles: You're right-I bypassed the lower levels, and saw only the top. I wonder if I'm
    missing all sorts of meaning on lower levels of my brain as well, by reading only the
    symbol level. It's too bad that the top level doesn't contain all the information about
    the bottom level, so that by reading the top, one also learns what the bottom level
    says. But I guess it would be naive to hope that the top level encodes anything from
    the bottom level-it probably doesn't percolate up. The MU-picture is the most striking
    possible example of that: there, the topmost level says only "ML which bears no
    relation whatever to the lower levels!
Crab: That's absolutely true. (Picks up the MU-picture, to inspect it more closely.) Hmm
    ... There's something strange about the smallest letters in this picture; they're very
    wiggly ... '
Anteater: Let me take a look. (Peers closely at the MU-picture.) I think there's yet
    another level, which all of us missed!
Tortoise: Speak for yourself, Dr. Anteater.
Achilles: Oh, no-that can't be! Let me see. (Looks very carefully.) I know the rest of you
    won't believe this, but the message of this picture is staring us all in the face, hidden
    in its depths. It is simply one word, repeated over and over again, like a mantra-but
    what an important one: "Mu"! What do you know! It is the same as the top level!
    And none of us suspected it in the least.
Crab: We would never have noticed it if it hadn't been for you, Achilles. Anteater: I
    wonder if the coincidence of the highest and lowest levels happened by chance? Or
    was it a purposeful act carried out by some creator?
Crab: How could one ever decide that?
Tortoise: I don't see any way to do so, since we have no idea why that particular picture is
    in the Crab's edition of the Well-Tempered Clavier. Anteater: Although we have been
    having a lively discussion, I have still managed to listen with a good fraction of an
    ear to this very long and complex four-voice fugue. It is extraordinarily beautiful.
Tortoise: It certainly is. And now, in just a moment, comes an organ point.
Achilles: Isn't an organ point what happens when a piece of music slows down slightly,
    settles for a moment or two on a single note or chord, and then resumes at normal
    speed after a short silence?
Tortoise: No, you're thinking of a "fermata"-a sort of musical semicolon. Did you notice
    there was one of those in the prelude?
Achilles: I guess I must have missed it.
Tortoise: Well, you have another chance coming up to hear a fermata-in fact, there are a
    couple of them coming up, towards the end of this fugue.
Achilles: Oh, good. You'll point them out in advance, won't you? Tortoise: If you like.
Achilles: But do tell me, what is an organ point?
Tortoise: An organ point is the sustaining of a single note by one of the




Ant Fugue                                                                                335

    voices in a polyphonic piece (often the lowest voice), while the other voices continue
    their own independent lines, This organ point is on the note of G. Listen carefully,
    and you'll hear it.
Anteater:. There occurred an incident one day when I visited with Aunt Hillary which
    reminds me of your suggestion of observing the symbols in Achilles' brain as they
    create thoughts which are about themselves.
Crab: Do tell us about it.
Anteater: Aunt Hillary had been feeling very lonely, and was very happy to have
    someone to talk to that day. So she gratefully told me to help myself to the juiciest
    ants I could find. (She's always been most generous with her ants.)
Achilles: Gee!
Anteater: It just happened that I had been watching the symbols which were carrying out
    her thoughts, because in them were some particularly juicy-looking ants.
Achilles: Gee!
Anteater: So I helped myself to a few of the fattest ants which had been parts of the
    higher-level symbols which I had been reading. Specifically, the symbols which they
    were part of were the ones which had expressed the thought, "Help yourself to any of
    the ants which look appetizing."
Achilles: Gee!
Anteater: Unfortunately for them, but fortunately for me, the little bugs didn't have the
    slightest inkling of what they were collectively telling me, on the symbol level.
Achilles: Gee! That is an amazing wraparound. They were completely unconscious of
    what they were participating in. Their acts could be seen as part of a pattern on a
    higher level, but of course they were completely unaware of that. Ah, what a pity-a
    supreme irony, in fact-that they missed it.
Crab: You are right, Mr. T-that was a lovely organ point.
Anteater: I had never heard one before, but that one was so conspicuous that no one could
    miss it. Very effective.
Achilles: What? Has the organ point already occurred? How can I not have noticed it, if it
    was so blatant?
Tortoise: Perhaps you were so wrapped up in what you were saying that you were
    completely unaware of it. Ah, what a pity-a supreme irony, in fact-that you missed it.
Crab: Tell me, does Aunt Hillary live in an anthill?
Anteater: Well, she owns a rather large piece of property. It used to belong to someone
    else, but that is rather a sad story. In any case, her estate is quite expansive. She lives
    rather sumptuously, compared to many other colonies.
!chilies: How does that jibe with the communistic nature of ant colonies which you
    earlier described to us? It sounds quite inconsistent, to me, to preach communism and
    to live in a fancy estate.




Ant Fugue                                                                                 336

Anteater: The communism is on the ant level. In an ant colony all ants work for the
   common good, even to their own individual detriment at times. Now this is simply a
   built-in aspect of Aunt Hillary's structure, but for all I know, she may not even be
   aware of this internal communism. Most human beings are not aware of anything
   about their neurons; in fact they probably are quite content not to know anything
   about their brains, being somewhat squeamish creatures. Aunt Hillary is also
   somewhat squeamish; she gets rather antsy whenever she starts to think about ants at
   all. So she avoids thinking about them whenever possible. I truly doubt that she
   knows anything about the communistic society which is built into her very structure.
   She herself is a staunch believer in libertarianism-you know, laissez-faire and all that.
   So it makes perfect sense, to me at least, that she should live in a rather sumptuous
   manor.




Tortoise: As I turned the page just now, while following along in this lovely edition of
    the Well-Tempered Clavier, I noticed that the first of the two fermatas is coming up
    soon-so you might listen for it, Achilles. Achilles: I will, I will.
Tortoise: Also, there's a most curious picture facing this page. Crab: Another one? What
    next?
Tortoise: See for yourself. (Passes the score over to the Crab.)
Crab: Aha! It's just a few bunches of letters. Let's see-there are various numbers of the
    letters `J', 'S', `B', `m', `a', and 't'. It's strange, how the first three letters grow, and then
    the last three letters shrink again. Anteater: May I see it?
Crab: Why, certainly.
Anteater: Oh, by concentrating on details, you have utterly missed the big picture. In
    reality, this group of letters is `f', `e', `r', 'A', `C', 'H', without any repetitions. First
    they get smaller, then they get bigger. Here, Achilles-what do you make of it?
Achilles: Let me see. Hmm. Well, I see it as a set of upper-case letters which grow as you
    move to the right.
Tortoise: Do they spell anything?




Ant Fugue                                                                                        337

Achilles: Ah ... "J. S. BACH". Oh! I understand now. It's Bach's name!
Tortoise: Strange that you should see it that way. I see it as a set of lower-case letters,
    shrinking as they move to the right, and ... spelling out ... the name of ... (Slows down
    slightly, especialh drawing out the last few words. Then there is a brief silence.
    Suddenly he resumes as if nothing unusual had happened.) -"fermat".
Achilles: Oh, you've got Fermat on the brain, I do believe. You see Fermat's Last
    Theorem everywhere.
Anteater: You were right, Mr. Tortoise-I just heard a charming little fermata in the fugue.
Crab: So did I.
Achilles: Do you mean everybody heard it but me? I'm beginning to feel stupid.
Tortoise: There, there, Achilles-don't feel bad. I'm sure you won't miss Fugue's Last
    Fermata (which is coming up quite soon). But, to return to our previous topic, Dr.
    Anteater, what is the very sad story which you alluded to, concerning the former
    owner of Aunt Hillary's property
Anteater: The former owner was an extraordinary individual, one of the most creative ant
    colonies who ever lived. His name was Johant Sebastiant Fermant, and he was a
    mathematiciant by vocation, but a musiciant by avocation.
Achilles: How very versantile of him!
Anteater: At the height of his creative powers, he met with a most untimely demise. One
    day, a very hot summer day, he was out soaking up the warmth, when a freak
    thundershower-the kind that hits only once every hundred years or so-appeared from
    out of the blue, and thoroughly drenched J. S F. Since the storm came utterly without
    warning, the ants got completely disoriented and confused. The intricate organization
    which had been so finely built up over decades, all went down the drain in a matter of
    minutes. It was tragic.
Achilles: Do you mean that all the ants drowned, which obviously would spell the end of
    poor J. S. F.
Anteater: Actually, no. The ants managed to survive, every last one of them, by crawling
    onto various sticks and logs which floated above the raging torrents. But when the
    waters receded and left the ants back on their home grounds, there was no
    organization left. The caste distribution was utterly destroyed, and the ants
    themselves had no ability to reconstruct what had once before been such a finely
    tuned organization. They were as helpless as the pieces of Humpty Dumpty in putting
    themselves back together again. I myself tried, like all the king's horses and all the
    king's men, to put poor Fermant together again. I faithfully put out sugar and cheese,
    hoping against hope that somehow Fermant would reappear ... (Pulls out a
    handkerchief and wipes his eyes.)
Achilles: How valiant of you! I never knew Anteaters had such big hearts.
Anteater: But it was all to no avail. He was Bone, beyond reconstitution.




Ant Fugue                                                                                338

    However, something very strange then began to take place: over the next few
    months, the ants which had been components of J. S. F. slowly regrouped, and built
    up a new organization. And thus was Aunt Hillary born.
Crab: Remarkable! Aunt Hillary is composed of the very same ants as Fermant was
Anteater: Well, originally she was, yes. By now, some of the older ants have died, and
    been replaced. But there are still many holdovers from the J. S. F.-days.
Crab: And can't you recognize some of J. S. F.'s old traits coming to the fore, from time
    to time, in Aunt Hillary\%
Anteater: Not a one. They have nothing in common. And there is no reason they should,
    as I see it. There are, after all, often several distinct ways to rearrange a group of
    parts to form a "sum". And Aunt Hillary was just a new "sum" of the old parts. Not
    MORE than the sum, mind you just that particular KIND of sum.
Tortoise: Speaking of sums, I am reminded of number theory, where occasionally one
    will be able to take apart a theorem into its component symbols, rearrange them in a
    new order, and come up with a new theorem.
Anteater: I've never heard of such a phenomenon, although I confess to being a total
    ignoramus in the field.
Achilles: Nor have I heard of it-and I am rather well versed in the field, If I don't say so
    myself. I suspect Mr. T is just setting up one of his elaborate spoofs. I know him
    pretty well by now. Anteater: Speaking of number theory, I am reminded of J. S. F.
    again, for number theory is one of the domains in which he excelled. In fact, he made
    some rather rema, ..able contributions to number theory. Aunt Hillary, on the other
    hand, is remarkably dull-witted in anything that has even the remotest connection
    with mathematics. Also, she has only a rather banal taste in music, whereas
    Sebastiant was extremely gifted in music.
Achilles: I am very fond of number theory. Could you possibly relate to us something of
    the nature of Sebastiant's contributions,
Anteater: Very well, then. (Pauses for a moment to sip his tea, then resumes.)
    Have you heard of Fourmi's infamous "Well-Tested Conjecture”?
Achilles. I'm not sure ... It sounds strangely familiar, and yet I can't quite place it.
Anteater: It's a very simple idea. Lierre de Fourmi, a mathematiciant by vocation but
    lawyer by avocation, had been reading in his copy-of the classic text Arithmetica by
    Di of Antus, and came across a page containing the equation
                                        2a+2b=2c
    He immediately realized that this equation has infinitely many solutions a. b, c, and
    then wrote in the margin the following notorious comment:




Ant Fugue                                                                               339

   FIGURE 63. During emigrations arm' ants sometimes create living bridges of their
   own bodies. In this photograph of such a bridge (de Fourmi Lierre), the workers of
   an Eciton burchelli colony can be seen linking their legs and, along the top of the
   bridge, hooking their tarsal claws together to form irregular systems of chains. .A
   symbiotic silverfish, Trichatelura manni, is seen crossing the bridge in the center.
   [From E. O. Wilson, The Insect Societies 'Cambridge, Mass.: Harvard University
   Press, 1971), p. 62)

   The equation
                                        na+nb=nc

   has solutions in positive integers a, b, c, and n only when n = 2 (and then there are
   infinitely many triplets a, b, c which satisfy the equation); but there are no solutions
   for n > 2. I have discovered a truly marvelous proof of this statement, which,
   unfortunately. is so small that it would be well-nigh invisible if written in the margin.
   Ever since that year, some three hundred days ago, mathematiciants have been vainly
   trying to do one of two things: either to prove Fourmi's claim, and thereby vindicate
   Fourmi's reputation, which, although very high, has been somewhat tarnished by
   skeptics who think he never really found the proof he claimed to have found-or else
   to refute the claim, by finding a counterexample: a set of four integers a, b, c, and n,
   with n > 2, which satisfy the equation. Until very recently, every attempt in either
   direction had met with failure. To be sure, the Conjecture has been verified for many
   specific values of n-in particular, all n up to 125,000. But no one had succeeded in
   proving it for ALL n-no one, that is, until Johant Sebastiant Fermant came upon the
   scene. It was he who found the proof that cleared Fourmi's name.



Ant Fugue                                                                               340

        It now goes under the name "Johant Sebastiant's Well-Tested Conjecture".
   Achilles: Shouldn't it be called a "Theorem" rather than a "Conjecture", if it's finally been
        given a proper proof;
   Anteater: Strictly speaking, you're right, but tradition has kept it this way.
   Tortoise: What sort of music did Sebastiant do?
   Anteater: He had great gifts for composition. Unfortunately, his greatest work is shrouded
        in mystery, for he never reached the point of publishing it. Some believe that he had
        it all in his mind; others are more unkind, saying that he probably never worked it out
        at all, but merely blustered about it.
   Achilles: What was the nature of this magnum opus?
   Anteater: It was to be a giant prelude and fugue; the fugue was to have
        twenty-four voices, and to involve twenty-four distinct subjects, one in
        each of the major and minor keys.
   Achilles: It would certainly be hard to listen to a twenty-four-voice fugue
        as a whole!
   Crab: Not to mention composing one!
Anteater: But all that we know of it is Sebastiant's description of it, which he wrote in the
   margin of his copy of Buxtehude's Preludes and Fugues for Organ. The last words which
   he wrote before his tragic demise were:

                I have composed a truly marvelous fugue. In it, I have added
                together the power of 24 keys, and the power of 24 themes; I
                came up with a fugue with the power of 24 voices. Unfortunately,
                this margin is too narrow to contain it.

        And the unrealized masterpiece simply goes by the name, "Fermant's Last Fugue".
    Achilles: Oh, that is unbearably tragic.
    Tortoise: Speaking of fugues, this fugue which we have been listening to is nearly over.
        Towards the end, there occurs a strange new twist on its theme. (Flips the page in the
        Well-Tempered Clavier.) Well, what have we here? A new illustration-how
        appealing! (Shows it to the Crab.)




    Ant Fugue                                                                               341

Crab: Well, what have we here? Oh, I see: It´s HOLISMIONSIM”, written in large letters
    that first shrink and then grow back to their original size. But that doesn't make any
    sense, because it's not a word. Oh me, oh my! (Passes it to the Anteater.)
Anteater: Well, what have we here? Oh, I see: it's "REDUCTHOLISM", written in small
    letters that first grow and then shrink back to their original size. But that doesn't make
    any sense, because it's not a word. Oh my, oh me! (Passes it to Achilles.)
Achilles: I know the rest of you won't believe this, but in fact this picture consists of the
    word "HOLISM" written twice, with the letters continually shrinking as they proceed
    from left to right. (Returns it to the Tortoise.)
Tortoise: I know the rest of you won't believe this, but in fact this picture consists of the
    word "REDUCTIONISM" written once, with the letters continually growing as they
    proceed from left to right.
Achilles: At last-I heard the new twist on the theme this time! I am so glad that you
    pointed it out to me, Mr. Tortoise. Finally, I think I am beginning to grasp the art of
    listening to fugues.




Ant Fugue                                                                                 342

                            Brains and Thoughts
                          New Perspectives on Thought
IT WAS ONLY with the advent of computers that people actually tried to create
"thinking" machines, and witnessed bizarre variations on the theme, of thought. Programs
were devised whose "thinking" was to human thinking as a slinky flipping end over end
down a staircase is to human locomotion. All of a sudden the idiosyncrasies, the
weaknesses and powers, the vagaries and vicissitudes of human thought were hinted at by
the newfound ability to experiment with alien, yet hand-tailored forms of thought-or
approximations of thought. As a result, we have acquired, in the last twenty years or so, a
new kind of perspective on what thought is, and what it is not. Meanwhile, brain
researchers have found out much about the small-scale and large-scale hardware of the
brain. This approach has not yet been able to shed much light on how the brain
manipulates concepts, but it gives us some ideas about the biological mechanisms on
which thought manipulation rests.
In the coming two Chapters, then, we will try to unite some insights gleaned from
attempts at computer intelligence with some of the facts learned from ingenious
experiments on living animal brains, as well as with results from research on human
thought processes done by cognitive psychologists. The stage has been set by the
Prelude, Ant Fugue; now we develop the ideas more deeply.

                        Intensionality and Extensionality
Thought must depend on representing reality in the hardware of the brain. In the
preceding Chapters, we have developed formal systems which represent domains of
mathematical reality in their symbolisms. To what extent is it reasonable to use such
formal systems as models for how the brain might manipulate ideas?
We saw, in the pq-system and then in other more complicated systems, how meaning, in
a limited sense of the term, arose as a result of an isomorphism which maps typographical
symbols onto numbers, operations, and relations; and strings of typographical symbols
onto statements. Now in the brain we don't have typographical symbols, but we have
something even better: active elements which can store information and transmit it and
receive it from other active elements. Thus we have active symbols, rather than passive
typographical symbols. In the brain, the rules




Brains and Thoughts                                                                    343

re mixed right in with the symbols themselves, whereas on paper, the symbols are static
entities, and the rules are in our heads.
It is important not to get the idea, from the rather strict nature of all ie formal systems we
have seen, that the isomorphism between symbols and real things is a rigid, one-to-one
mapping, like the strings which link a marionette and the hand guiding it. In TNT, the
notion "fifty" can be expressed in different symbolic ways; for example,

                           ((SSSSSSSO.SSSSSSSO)+(SO-SO))
                        ((SSSSSO•SSSSSO)+(SSSSSO.SSSSSO))

'hat these both represent the same number is not a priori clear. You can manipulate each
expression independently, and at some point stumble cross a theorem which makes you
exclaim, "Oh-it's that number!"
In your mind, you can also have different mental descriptions for a single person; for
example,

      The person whose book I sent to a friend in Poland a while back.

      The stranger who started talking with me and my friends tonight in this coffee
      house.

:'hat they both represent the same person is not a priori clear. Both descriptions may sit in
your mind, unconnected. At some point during the evening you may stumble across a
topic of conversation which leads to the revelation that they designate the same person,
making you exclaim, Oh-you're that person!"
Not all descriptions of a person need be attached to some central symbol for that person,
which stores the person's name. Descriptions can be manufactured and manipulated in
themselves. We can invent nonexistent people by making descriptions of them; we can
merge two descriptions 'hen we find they represent a single entity; we can split one
description into two when we find it represents two things, not one-and so on. This
calculus of descriptions" is at the heart of thinking. It is said to be intentional and not
extensional, which means that descriptions can "float" without Being anchored down to
specific, known objects. The intensionality of thought is connected to its flexibility; it
gives us the ability to imagine hypothetical worlds, to amalgamate different descriptions
or chop one description into separate pieces, and so on.
Suppose a friend who has borrowed your car telephones you to say hat your car skidded
off a wet mountain road, careened against a bank, .nd overturned, and she narrowly
escaped death. You conjure up a series \& images in your mind, which get progressively
more vivid as she adds details, and in the end you "see it all in your mind's eye". Then
she tells you hat it's all been an April Fool's joke, and both she and the car are fine! In
many ways that is irrelevant. The story and the images lose nothing of their vividness,
and the memory will stay with you for a long, long time. Later, you may even think of
her as an unsafe driver because of the strength of




Brains and Thoughts                                                                       344

the first impression, which should have been wiped out when you learned it was all
untrue. Fantasy and fact intermingle very closely in our minds, and this is because
thinking involves the manufacture and manipulation of complex descriptions, which need
in no way be tied down to real events or things.
A flexible, intensional representation of the world is what thinking is all about. Now how
can a physiological system such as the brain support such a system?

                                 The Brain's "Ants"
The most important cells in the brain are nerve cells, or neurons (see Fig. 65), of which
there are about ten billion. (Curiously, outnumbering the neurons by about ten to one are
the glial cells, or glia. Glia are believed to play more of a supporting role to the neurons'
starring role, and therefore we will not discuss them.) Each neuron possesses a number of
synapses ("entry ports") and one axon ("output channel"). The input and output are
electrochemical flows: that is, moving ions. In between the entry ports of a neuron and its
output channel is its cell body, where "decisions" are made.




FIGURE 65. Schematic drawing of a neuron. [Adapted From D. Wooldridge, The
Machinery of the Brain (New York:"- McGraw-Hill, 1963), p. 6.


Brains and Thoughts                                                                      345

The type of decision which a neuron faces-and this can take place up to a thousand times
per second-is this: whether or not to fire-that is, to ease ions down its axon, which -
eventually will cross over into the entry its of one or more other neurons, thus causing
them to make the same sort of decision. The decision is made in a very simple manner: if
the sum all inputs exceeds a certain threshold, yes; otherwise, no. Some of the inputs can
be negative inputs, which cancel out positive inputs coming from somewhere else. In any
case, it is simple addition which rules the lowest 'el of the mind. To paraphrase Descartes'
famous remark, "I think, therefore I sum" (from the Latin Cogito, ergo am).
         Now although the manner of making the decision sounds very simple, here is one
fact which complicates the issue: there may be as many as 200,000 separate entry ports to
a neuron, which means that up to 200,000 Karate summands may be involved in
determining the neuron's next ion. Once the decision has been made, a pulse of ions
streaks down the on towards its terminal end. Before the ions reach the end, however, ey
may encounter a bifurcation-or several. In such cases, the single output pulse splits up as
it moves down the bifurcating axon, and by the tine it has reached the end, "it" has
become "they"-and they may reach their destinations at separate times, since the axon
branches along which they travel may be of different lengths and have different
resistivities. The important thing, though, is that they all began as one single pulse,
moving 'ay from the cell body. After a neuron fires, it needs a short recovery time fore
firing again; characteristically this is measured in milliseconds, so at a neuron may fire up
to about a thousand times per second.

                          Larger Structures in the Brain
Now we have described the brain's "ants". What about "teams", or "signals"? What about
"symbols"? We make the following observation: despite e complexity of its input, a
single neuron can respond only in a very primitive way-by firing, or not firing. This is a
very small amount of Formation. Certainly for large amounts of information to be carried
or processed, many neurons must be involved. And therefore one might guess at larger
structures, composed from many neurons, would exist, which handle concepts on a
higher level. This is undoubtedly true, but the most naive assumption-that there is a fixed
group of neurons for each different concept-is almost certainly false.
         There are many anatomical portions of the brain which can be distinguished from
each other, such as the cerebrum, the cerebellum, the hypothalamus (see Fig. 66). The
cerebrum is the largest part of the human am, and is divided into a left hemisphere and a
right hemisphere. The outer few millimeters of each cerebral hemisphere are coated with
a layered "bark", or cerebral cortex. The amount of cerebral cortex is the major
distinguishing feature, in terms of anatomy, between human brains and brains of less
intelligent species. We will not describe any of the brain's suborgans in detail because, as
it turns out, only the roughest mapping can




Brains and Thoughts                                                                      346

FIGURE 66. The human brain, seen from the left side. It is strange that the visual area is
in the back of the head. [From Steven Rose, The Conscious Brain, updated ed. (New
York: Vintage, 1966), p. 50. ]

at this time be made between such large-scale suborgans and the activities, mental or
physical, which they are responsible for. For instance, it is known that language is
primarily handled in one of the two cerebral hemispheres-in fact, usually the left
hemisphere. Also, the cerebellum is the place where trains of impulses are sent off to
muscles to control motor activity. But how these areas carry out their functions is still
largely a mystery.

                            Mappings between Brains
Now an extremely important question comes up here. If thinking does take place in the
brain, then how are two brains different from each other? How is my brain different from
yours? Certainly you do not think exactly as I do, nor as anyone else does. But we all
have the same anatomical divisions in our brains. How far does this identity of brains
extend? Does it go to the neural level? Yes, if you look at animals on a low enough level
of the thinking-hierarchy-the lowly earthworm, for instance. The following quote is from
the neurophysiologist, David Hubel, speaking at a conference on communication with
extraterrestrial intelligence:

     The number of nerve cells in an animal like a worm would be measured, I
     suppose, in the thousands. One very interesting thing is that we may point to a
     particular individual cell in a particular earthworm, and then identify the same
     cell, the corresponding cell in another earthworm of the same species.'




Brains and Thoughts                                                                     347

Earthworms have isomorphic brains! One could say, "There is only one earthworm."
  But such one-to-one mappability between individuals' brains disappears very soon as
you ascend in the thinking-hierarchy and the number of neurons increases-confirming
one's suspicions that there is not just one pan! Yet considerable physical similarity can be
detected between different human brains when they are compared on a scale larger than a
;le neuron but smaller than the major suborgans of the brain. What s this imply about how
individual mental differences are represented in physical brain? If we looked at my
neurons' interconnections, could we l various structures that could be identified as coding
for specific things -tow, specific beliefs I have, specific hopes, fears, likes and dislikes I
harbor? If mental experiences can be attributed to the brain, can knowledge and other
aspects of mental life likewise be traced to specific locations de the brain, or to specific
physical subsystems of the brain? This will be a central question to which we will often
return in this Chapter and the next.

                       Localization of Brain Processes: An Enigma

In an attempt to answer this question, the neurologist Karl Lashley, in a series of
experiments beginning around 1920 and running for many ,s, tried to discover where in
its brain a rat stores its knowledge about :e running. In his book The Conscious Brain,
Steven Rose describes Lashley's trials and tribulations this way:

   Lashley was attempting to identify the locus of memory within the cortex, and, to do so,
   first trained rats to run mazes, and then removed various cortical regions. He allowed the
   animals to recover and tested the retention of the maze-running skills. To his surprise it
   was not possible to find a particular region corresponding to the ability to remember the
   way through a maze. instead all the rats which had had cortex regions removed suffered
   some kind f impairment, and the extent of the impairment was roughly proportional to the
   amount of cortex taken off. Removing cortex damaged the motor and sensory capacities
   of the animals, and they would limp, hop, roll, or stagger, but somehow they always
   managed to traverse the maze. So far as memory 'as concerned, the cortex appeared to be
   equipotential, that is, with all regions of equal possible utility. Indeed, Lashley concluded
   rather gloomily in is last paper "In Search of the Engram", which appeared in 1950, that
   the only conclusion was that memory was not possible at all.'

  Curiously, evidence for the opposite point of view was being developed :in Canada at
roughly the same time that Lashley was doing his last work, in late 1940's. The
neurosurgeon Wilder Penfield was examining the reactions of patients whose brains had
been operated on, by inserting electrodes into various parts of their exposed brains, and
then using small electrical pulses to stimulate the neuron or neurons to which the
electrodes been attached. These pulses were similar to the pulses which come other
neurons. What Penfield found was that stimulation of certain




Brains and Thoughts                                                                            348

neurons would reliably create specific images or sensations in the patient. These
artificially provoked impressions ranged from strange but indefinable fears to buzzes and
colors, and, most impressively of all, to entire successions of events recalled from some
earlier time of life, such as a childhood birthday party. The set of locations which could
trigger such specific events was extremely small-basically centered upon a single neuron.
Now these results of Penfield dramatically oppose the conclusions of Lashley, since they
seem to imply that local areas are responsible for specific memories, after all.
  What can one make of this? One possible explanation could be that memories are coded
locally, but over and over again in different areas of the cortex-a strategy perhaps
developed in evolution as security against possible loss of cortex in fights, or in
experiments conducted by neurophysiologists. Another explanation would be that
memories can be reconstructed from dynamic processes spread over the whole brain, but
can be triggered from local spots. This theory is based on the notion of modern telephone
networks, where the routing of a long-distance call is not predictable in advance, for it is
selected at the time the call is placed, and depends on the situation all over the whole
country. Destroying any local part of the network would not block calls; it would just
cause them to be routed around the damaged area. In this sense any call is potentially
nonlocalizable. Yet any call just connects up two specific points; in this sense any call is
localizable.

                             Specificity in Visual Processing

Some of the most interesting and significant work on localization of brain processes has
been done in the last fifteen years by David Hubel and Torsten Wiesel, at Harvard. They
have mapped out visual pathways in the brains of cats, starting with the neurons in the
retina, following their connections towards the rear of the head, passing through the
"relay station" of the lateral geniculate, and ending up in the visual cortex, at the very
back of the brain. First of all, it is remarkable that there exist well defined neural
pathways, in light of Lashley's results. But more remarkable are the properties of the
neurons located at different stages along the pathway.
  It turns out that retinal neurons are primarily contrast sensors. More specifically, the
way they act is this. Each retinal neuron is normally firing at a "cruising speed". When its
portion of the retina is struck by light, it may either fire faster or slow down and even
stop firing. However, it will do so only provided that the surrounding part of the retina is
less illuminated. So this means that there are two types of neuron: ."on-center", and "off-
center". The on-center neurons are those whose firing rate increases whenever, in the
small circular retinal area to which they are sensitive, the center is bright but the outskirts
are dark; the off-center neurons are those which fire faster when there is darkness in the
center and brightness in the




Brains and Thoughts                                                                        349

outer ring. If an on-center pattern is shown to an off-center neuron, the neuron will slow
down in firing (and vice versa). Uniform illumination will .leave both types of retinal
neuron unaffected; they will continue to fire at cruising speed.
  From the retina, signals from these neurons proceed via the optic nerve to the lateral
geniculate, located somewhere towards the middle of the brain. There, one can find a
direct mapping of the retinal surface in the .use that there are lateral-geniculate neurons
which are triggered only by specific stimuli falling on specific areas of the retina. In that
sense, the general geniculate is disappointing; it seems to be only a "relay station", and
not a further processor (although to give it its due, the contrast sensitivity ,ms to be
enhanced in the lateral geniculate). The retinal image is coded a straightforward way in
the firing patterns of the neurons in the lateral geniculate, despite the fact that the neurons
there are not arranged on a o-dimensional surface in the form of the retina, but in a three-
dimensional block. So two dimensions get mapped onto three, yet the formation is
preserved: an isomorphism. There is probably some deep meaning to the change in the
dimensionality of the representation, which is not yet fully appreciated. In any case, there
are so many further unexplained stages of vision that we should not be disappointed but
pleased the fact that-to some extent-we have figured out this one stage!
  From the lateral geniculate, the signals proceed back to the visual cortex. Here, some
new types of processing occur. The cells of the visual cortex are divided into three
categories: simple, complex, and hyper complex. Simple cells act very much like retinal
cells or lateral geniculate [Is: they respond to point-like light or dark spots with
contrasting surrounds, in particular regions of the retina. Complex cells, by contrast,
usually receive input from a hundred or more other cells, and they detect light dark bars
oriented at specific angles on the retina (see Fig. 67). Hyper complex cells respond to
corners, bars, or even "tongues" moving in specific directions (again see Fig. 67). These
latter cells are so highly specialized at they are sometimes called "higher-order hyper
complex cells".

                                A "Grandmother Cell"?
  Because of the discovery of cells in the visual cortex which can be triggered stimuli of
ever-increasing complexity, some people have wondered if things are not leading in the
direction of "one cell, one concept"-for ample, you would have a "grandmother cell"
which would fire if, and only if, your grandmother came into view. This somewhat
humorous ample of a "superhypercomplex cell" is not taken very seriously. Rower, it is
not obvious what alternative theory seems reasonable. One possibility is that larger neural
networks are excited collectively by sufficiently complex visual stimuli. Of course, the
triggering of these larger multineuron units would somehow have to come from
integration of signals emanating from the many hyper complex cells. How this might be
done nobody knows> Just when we seem to be approaching the threshold where




Brains and Thoughts                                                                        350

FIGURE       67.    Responses      to    patterns    by    certain    sample     neurons.
             (a) This edge-detecting neuron looks for vertical edges with light on the left
             and dark on the right. The first column shows how the orientation of an
             edge is relevant to this neuron. The second column shows how the position
             of the edge within the field is irrelevant, for this particular neuron. (b)
             Showing how a hyper complex cell responds more selectively: here, only
             when the descending tongue is in the middle of the field. (c) The responses
             of a hypothetical "grandmother cell" to various random stimuli; the reader
             may enjoy pondering how an "octopus cell" would respond to the same
             stimuli.

  "symbol" might emerge from "signal", the trail gets lost-a tantalizingly unfinished story.
We will return to this story shortly, however, and try to fill in some of it.
  Earlier I mentioned the coarse-grained isomorphism between all human brains which
exists on a large anatomical scale, and the very fine-grained, neural-level isomorphism
which exists between earthworm brains. It is quite interesting that there is also an
isomorphism between the visual processing apparatus of cat, monkey, and human, the
"grain" of which is somewhere between coarse and fine. Here is how that isomorphism
works. First of all, all three species have "dedicated" areas of cortex at the back of their
brains where visual processing is done: the visual cortex. Secondly, in




Brains and Thoughts                                                                     351

each of them, the visual cortex breaks up into three subregions, called areas 18, and 19 of
the cortex. These areas are still universal, in the sense that y can be located in the brain of
any normal individual in any of the three ties. Within each area you can go still further,
reaching the "columnar" organization of the visual cortex. Perpendicular to the surface of
the (ex, moving radially inwards towards the inner brain, visual neurons are inged in
"columns"-that is, almost all connections move along the ial, columnar direction, and not
between columns. And each column ps onto a small, specific retinal region. The number
of columns is not same in each individual, so that one can't find "the same column". ally,
within a column, there are layers in which simple neurons tend to found, and other layers
in which complex neurons tend to be found. to hypercomplex neurons tend to be found in
areas 18 and 19 predominately, while the simple and complex ones are found mostly in
area 17.) appears that we run out of isomorphisms at this level of detail. From here down
to the individual neuron level, each individual cat, monkey, or man has a completely
unique pattern-somewhat like a fingerprint or a signature.
  One minor but perhaps telling difference between visual processing in ;'brains and
monkeys' brains has to do with the stage at which informal from the two eyes is
integrated to yield a single combined higher-level 1al. It turns out that it takes place
slightly later in the monkey than in the cat, which gives each separate eye's signal a
slightly longer time to get processed by itself. This is not too surprising, since one would
expect that higher a species lies in the intelligence hierarchy, the more complex will the
problems which its visual system will be called upon to handle; and before signals ought
to pass through more and more early processing ore receiving a final "label". This is quite
dramatically confirmed by observations of the visual abilities of a newborn calf, which
seems to be born with as much power of visual discrimination as it will ever have. It will
shy away from people or dogs, but not from other cattle. Probably its entire visual system
is "hard-wired" before birth, and involves relatively little optical processing. On the other
hand, a human's visual system, so deeply ant on the cortex, takes several years to reach
maturity.

                              Funneling into Neural Modules

  A puzzling thing about the discoveries so far made about the organization the brain is
that few direct correspondences have been found between large-scale hardware and high-
level software. The visual cortex, for instance, is a large-scale piece of hardware, which is
entirely dedicated to a it software purpose-the processing of visual information-yet all of
processing so far discovered is still quite low-level. Nothing approaching recognition of
objects has been localized in the visual cortex. This means that no one knows where or
how the output from complex and hypercomplex cells gets transformed into conscious
recognition of shapes,




Brains and Thoughts                                                                        352

rooms, pictures, faces, and so on. People have looked for evidence of the "funneling" of
many low-level neural responses into fewer and fewer higher-level ones, culminating in
something such as the proverbial grandmother cell, or some kind of multineuron network,
as mentioned above. It is evident that this will not be found in some gross anatomical
division of the brain, but rather in a more microscopic analysis.
  One possible alternative to the the grandmother cell might be a fixed set of neurons, say
a few dozen, at the thin end of the "funnel", all of which fire when Granny comes into
view. And for each different recognizable object, there would be a unique network and a
funneling process that would focus down onto that network. There are more complicated
alternatives along similar lines, involving networks which can be excited in different
manners, instead of in a fixed manner. Such networks would be the "symbols" in our
brains.
  But is such funneling necessary? Perhaps an object being looked at is implicitly
identified by its "signature" in the visual cortex-that is, the collected responses of simple,
complex, and hypercomplex cells. Perhaps the brain does not need any further recognizer
for a particular form. This theory, however, poses the following problem. Suppose you
are looking at a scene. It registers its signature on your visual cortex; but then how do you
get from that signature to a verbal description of the scene? For instance, the paintings of
Edouard Vuillard, a French post-impressionist, often take a few seconds of scrutiny, and
then suddenly a human figure will jump out at you. Presumably the signature gets
imprinted on the visual cortex in the first fraction of a second-but the picture is only
understood after a few seconds. This is but one example of what is actually a common
phenomenon-a sensation of something "crystallizing" in your mind at the moment of
recognition, which takes place not when the light rays hit your retina, but sometime later,
after some part of your intelligence has had a chance to act on the retinal signals.
  The crystallization metaphor yields a pretty image derived from statistical mechanics,
of a myriad microscopic and uncorrelated activities in a medium, slowly producing local
regions of coherence which spread and enlarge; in the end, the myriad small events will
have performed a complete structural revamping of their medium from the bottom up,
changing' it from a chaotic assembly of independent elements into one large, coherent,
fully linked structure. If one thinks of the early neural activities as independent, and of
the end result of their many independent firings as the triggering of a well-defined large
"module" of neurons, then the word "crystallization" seems quite apt.
  Another argument for funneling is based on the fact that there are a myriad distinct
scenes which can cause you to feel you have perceived the same object-for example, your
grandmother, who may be smiling or frowning, wearing a hat or not, in a bright garden or
a dark train station, seen from near or far, from side or front, and so on. All these scenes
produce extremely different signatures on the visual cortex; yet all of them could prompt
you to say "Hello, Granny." So a funneling process must take




Brains and Thoughts                                                                       353

place at some point after the reception of the visual signature and before e words are
uttered. One could claim that this funneling is not part of the perception of Granny, but
just part of verbalization. But it seems quite unnatural to partition the process that way,
for you could internally use the formation that it is Granny without verbalizing it. It
would be very it unwieldy to handle all of the information in the entire visual cortex,
when much of it could be thrown away, since you don't care about where shadows fall or
how many buttons there are on her blouse, etc.
  Another difficulty with a non-funneling theory is to explain how there in be different
interpretations for a single signature-for example, the Escher picture Convex a4 Concave
(Fig. 23). Just as it seems obvious to us tat we do not merely perceive dots on a television
screen, but chunks, likewise it seems ridiculous to postulate that perception has taken
place hen a giant dot-like "signature" has been created on the visual cortex. here must be
some funneling, whose end result is to trigger some specific modules of neurons, each of
which is associated with the concepts-the funks-in the scene.

                       Modules Which Mediate Thought Processes

  Thus we are led to the conclusion that for each concept there is a fairly ell-defined
module which can be triggered-a module that consists of a nail group of neurons-a
"neural complex" of the type suggested earlier. problem with this theory-at least if it is
taken naively-is that it would suggest that one should be able to locate such modules
somewhere within to brain. This has not yet been done, and some evidence, such as the
experiments by Lashley, points against localization. However, it is still too early to tell.
There may be many copies of each module spread around, or modules may overlap
physically; both of these effects would tend to obscure any division of neurons into
"packets". Perhaps the complexes are like very thin pancakes packed in layers which
occasionally pass through each other; perhaps they are like long snakes which curl
around each other, here and there flattening out, like cobras' heads; perhaps they are like
spiderwebs; or perhaps they are circuits in which signals travel round id round in shapes
stranger than the dash of a gnat-hungry swallow. here is no telling. It is even possible that
these modules are software, ether than hardware, phenomena-but this is something which
we will discuss later
  There are many questions that come to mind concerning these hypothesized neural
complexes. For instance:

         Do they extend into the lower regions of the brain, such as the
            midbrain, the hypothalamus, etc.?
         Can a single neuron belong to more than one such complex?
         To how many such complexes can a single neuron belong?
         By how many neurons can such complexes overlap?




Brains and Thoughts                                                                      354

         Are these complexes pretty much the same for everybody?
         Are corresponding ones found in corresponding places in different
          people's brains?
         Do they overlap in the same way in everybody's brain?

  Philosophically, the most important question of all is this: "hat would the existence of
modules-for instance, a grandmother module-tell us? Would this give us any insight into
the phenomenon of our own consciousness? Or would it still leave us as much in the dark
about what consciousness is, as does knowledge that a brain is built out of neurons and
glia? As you might guess from reading the Ant Fugue, my feeling is that it would go a
long way towards giving us an understanding of the phenomenon of consciousness. The
crucial step that needs to be taken is from a low-level-neuron-by-neuron-description of
the state of a brain, to a high-level-module-by-module-description of the same state of the
same brain. Or, to revert to the suggestive terminology of the Ant Fugue, we want to shift
the description of the brain state from the signal level to the symbol, level.

                                      Active Symbols

  Let us from now on refer to these hypothetical neural complexes, neural modules,
neural packets, neural networks, multineuron units-call them what you will, whether they
come in the form of pancakes, garden rakes, rattlesnakes, snowflakes, or even ripples on
lakes-as symbols. A description of a brain state in terms of symbols was alluded to in the
Dialogue. What would such a description be like? What kinds of concepts is it reasonable
to think actually might be "symbolized"? What kinds of interrelations would symbols
have? And what insights would this whole picture provide into consciousness?
  The first thing to emphasize is that symbols can be either dormant, or awake (activated).
An active symbol is one which has been triggered-that is, one in which a threshold
number of neurons have been caused to fire by stimuli coming from outside. Since a
symbol can be triggered in many different ways, it can act in many different ways when
awakened. This suggests that we should think of a symbol not as a fixed entity, but as a
variable entity. Therefore it would not suffice to describe a brain state by saying
"Symbols A, B, ..., N are all active"; rather, we would have to supply in addition a set of
parameters for each active symbol, characterizing some aspects of the symbol's internal
workings. It is an interesting question whether in each symbol there are certain core
neurons, which invariably fire when the symbol is activated. If such a core set of neurons
exists, we might refer to it as the "invariant core" of the symbol. It is tempting to assume
that each time you think of, say, a waterfall, some fixed neural process is repeated,
without doubt embellished in different ways depending on the context, but reliably
occurring. However, it is not clear that this must be so.




Brains and Thoughts                                                                     355

  Now what does a symbol do, when awakened? A low-level description would say,
"Many of its neurons fire." But this no longer interests us. The high-level description
should eliminate all reference to neurons, and concentrate exclusively on symbols. So a
high-level description of what makes a symbol active, as distinguished from dormant,
would be, "It sends out messages, or signals, whose purpose is to try to awaken, or
trigger, other symbols." Of course these messages would be carried as streams of nerve
impulses, by neurons-but to the extent that we can avoid such phraseology, we should,
for it represents a low-level way of looking at things, and we hope that we can get along
on purely a high level. In other words, we hope at thought processes can be thought of as
being sealed off from neural events in the same way that the behavior of a clock is sealed
off from the laws of quantum mechanics, or the biology of cells is sealed off from the
laws of quarks.
  But what is the advantage of this high-level picture? Why is it better to say, "Symbols A
and B triggered symbol C" than to say, "Neurons 183 through 612 excited neuron 75 and
caused it to fire"? This question was answered in the Ant Fugue: It is better because
symbols symbolize things, and neurons don't. Symbols are the hardware realizations of
concepts. Whereas group of neurons triggering another neuron corresponds to no outer
event, the triggering of some symbol by other symbols bears a relation to events in the
real world-or in an imaginary world. Symbols are related to each other by the messages
which they can send back and forth, in such a way that their triggering patterns are very
much like the large-scale events rich do happen in our world, or could happen in a world
similar to ours. essence, meaning arises here for the same reason as it did in the -system-
isomorphism; only here, the isomorphism is infinitely more complex, subtle, delicate,
versatile, and intensional.
  Incidentally, the requirement that symbols should be able to pass sophisticated
messages to and fro is probably sufficient to exclude neurons themselves from playing
the role of symbols. Since a neuron has only a single way of sending information out of
itself, and has no way of selectively selecting a signal now in one direction, now in
another, it simply does not have the kind of selective triggering power which a symbol
must have to act e an object in the real world. I n his book The Insect Societies, E. O.
Wilson makes a similar point about how messages propagate around inside ant colonies:

 [Mass communication] is defined as the transfer, among groups, of information that a
 single individual could not pass to another.'

It is not such a bad image, the brain as an ant colony!
       The next question-and an extremely important one it is, too concerns the
nature and "size" of the concepts which are represented in the tin by single
symbols. About the nature of symbols there are questions like this: Would there be
a symbol for the general notion of waterfalls, or would there be different symbols
for various specific waterfalls? Or would both of these alternatives be realized?
About the "size" of symbols, there are questions like this: Would there be a symbol
for an entire story? Or for a




Brains and Thoughts                                                                     356

melody? Or a joke? Or is it more likely that there would only be symbols for concepts
roughly the size of words, and that larger ideas, such as phrases or sentences, would be
represented by concurrent or sequential activation of various symbols?
  Let us consider the issue of the size of concepts represented by symbols. Most thoughts
expressed in sentences are made up out of basic, quasi-atomic components which we do
not usually analyze further. These are of word size, roughly-sometimes a little longer,
sometimes a little shorter. For instance, the noun "waterfall", the proper noun "Niagara
Falls", the past-tense suffix "-ed", the verb "to catch up with", and longer idiomatic
phrases are all close to atomic. These are typical elementary brush strokes which we use
in painting portraits of more complex concepts, such as the plot of a movie, the flavor of
a city, the nature of consciousness, etc. Such complex ideas are not single brush strokes..
It seems reasonable to think that the brush strokes of language are also brush strokes of
thought, and therefore that symbols represent concepts of about this size. Thus a symbol
would be roughly something for which you know a word or stock phrase, or with which
you associate a proper name. And the representation in the brain of a more complex idea,
such as a problem in a love affair, would be a very complicated sequence of activations
of various symbols by other symbols.

                                  Classes and Instances
  There is a general distinction concerning thinking: that between categories and
individuals, or classes and instances. (Two other terms sometimes used are "types" and
"tokens".) It might seem at first sight that a given symbol would inherently be either a
symbol for a class or a symbol for an instance-but that is an oversimplification. Actually,
most symbols may play either role, depending on the context of their activation. For
example, look at the list below:

(1) a publication
 (2) a newspaper
        (3) The San Francisco Chronicle
                (4) the May 18 edition of the Chronicle
                        (5) my copy of the May 18 edition of the Chronicle
                               (6) my copy of the May 18 edition of the Chronicle as
                                      it was when I first picked it up (as contrasted with
                                      my copy as it was a few days later: in my fireplace,
                                      burning)

  Here, lines 2 to 5 all play both roles. Thus, line 4 is an instance of of the general class of
line 3, and line 5 is an instance of line 4. Line 6 is a special kind of instance of a class: a
manifestation. The successive stages of an object during its life history are its
manifestations. It is interesting to wonder if the cows on a farm perceive the invariant
individual underneath all the manifestations of the jolly farmer who feeds then hay.




Brains and Thoughts                                                                        357

                               The Prototype Principle
The list above seems to be a hierarchy of generality-the top being a very road conceptual
category, the bottom some very humble particular thing located in space and time.
However, the idea that a "class" must always be enormously broad and abstract is far too
limited. The reason- is that our thought makes use of an ingenious principle, which might
be called the prototype principle:

                 The most specific event can serve as a general example
                                  of a class of events.

  Everyone knows that specific events have a vividness which imprints them i strongly on
the memory that they can later be used as models for other vents which are like them in
some way. Thus in each specific event, there is the germ of a whole class of similar
events. This idea that there is generality in the specific is of far-reaching importance.
  Now it is natural to ask: Do the symbols in the brain represent classes, r instances? Are
there certain symbols which represent only classes, while other symbols represent only
instances? Or can a single symbol serve duty either as a class symbol or instance symbol,
depending which parts of it are activated? The latter theory seems appealing; one might
think that a "light" activation of a symbol might represent a class, and that a deeper, or
more complex, activation would contain more detailed internal neural firing patterns, and
hence would represent an instance. But on second thought, its is crazy: it would imply,
for example, that by activating the symbol for publication" in a sufficiently complex way,
you would get the very complex symbol which represents a specific newspaper burning
in my fireplace. And very other possible manifestation of every other piece of printed
matter would be represented internally by some manner of activating the single symbol
for "publication". That seems much too heavy a burden to place on to single symbol
"publication". One must conclude, therefore, that finance symbols can exist side by side
with class symbols, and are not just lodes of activation of the latter.

                       The Splitting-off of Instances from Classes

  On the other hand, instance symbols often inherit many of their properties from the
classes to which those instances belong. If I tell you I went to see a Movie, you will begin
"minting" a fresh new instance symbol for that particular movie; but in the absence of
more information, the new instance symbol will have to lean rather heavily on your pre-
existing class symbol for movie". Unconsciously, you will rely on a host of
presuppositions about at movie-for example, that it lasted between one and three hours,
that it was shown in a local theater, that it told a story about some people, and so i. These
are built into the class symbol as expected links to other symbols e., potential triggering
relations), and are called default options. In any




Brains and Thoughts                                                                      358

freshly minted instance symbol, the default options can easily be overridden, but unless
this is explicitly done, they will remain in the instance symbol, inherited from its class
symbol. Until they are overridden, they provide some preliminary basis for you to think
about the new instance for example, the movie I went to see-by using the reasonable
guesses which are supplied by the "stereotype", or class symbol.
  A fresh and simple instance is like a child without its own ideas or experiences-it relies
entirely on its parents' experiences and opinions and just parrots them. But gradually, as it
interacts more and more with the rest of the world, the child acquires its own
idiosyncratic experiences and inevitably begins to split away from the parents.
Eventually, the child becomes a full-fledged adult. In the same way, a fresh instance can
split off from its parent class over a period of time, and become a class, or prototype, in
its own right.
  For a graphic illustration of such a splitting-off process, suppose that some Saturday
afternoon you turn on your car radio, and happen to tune in on a football game between
two "random" teams. At first you do not know the names of the players on either team.
All you register, when the announcer says, "Palindromi made the stop on the twenty-
seven yard line, and that brings up fourth down and six to go," is that some player
stopped some other player. Thus it is a case of activation of the class symbol "football
player", with some sort of coordinated activation of the symbol for tackling. But then as
Palindromi figures in a few more key plays, you begin building up a fresh instance
symbol for him in particular, using his name, perhaps, as a focal point. This symbol is
dependent, like a child, on the class symbol for "football player": most of your image of
Palindromi is supplied by your stereotype of a football player as contained in the
"football player" symbol. But gradually, as more information comes to you, the
"Palindromi" symbol becomes more autonomous, and relies less and less on concurrent
activation of its parent class symbol. This may happen in a few minutes, as Palindromi
makes a few good plays and stands out. His teammates may still all be represented by
activations of the class symbol, however. Eventually, perhaps after a few days, when you
have read some articles in the sports section of your paper, the umbilical cord is broken,
and Palindromi can stand on his own two feet. Now you know such things as his home
town and his major in college; you recognize his face; and so on. At this point,
Palindromi is no longer conceived of merely as a football player, but as a human being
who happens also to be a football player. "Palindromi" is an instance symbol which can
become active while its parent class symbol (football player) remains dormant.
  Once, the Palindromi symbol was a satellite orbiting around its mother symbol, like an
artificial satellite circling the Earth, which is so much bigger and more massive. Then
there came an intermediate stage, where one symbol was more important than the other,
but they could be seen as orbiting around each other-something like the Earth and the
Moon. Finally, the new symbol becomes quite autonomous; now it might easily serve as
a class symbol around which could start rotating new satellites-




Brains and Thoughts                                                                      359

symbols for other people who are less familiar but who have something in common with
Palindromi, and for whom he can serve as a temporary stereotype, until you acquire more
information, enabling the new symbols so to become autonomous.

               The Difficulty of Disentangling Symbols from Each Other

These stages of growth and eventual detachment of an instance from a ass will be
distinguishable from each other by the way in which the symbols involved are linked.
Sometimes it will no doubt be very difficult to 11 just where one symbol leaves off and
the other one begins. How "active" the one symbol, compared to the other? If one can be
activated independently of the other, then it would be quite sensible to call them
autonomous.
  We have used an astronomy metaphor above, and it is interesting that to problem of the
motion of planets is an extremely complex one-in fact the general problem of three
gravitationally interacting bodies (such as the Earth, Moon, and Sun) is far from solved,
even after several centuries of work. One situation in which it is possible to obtain good
approximate solutions, however, is when one body is much more massive than the other
two (here, the Sun); then it makes sense to consider that body as stationary, with the other
two rotating about it: on top of this can finally be added the interaction between the two
satellites. But this approximation depends on breaking up the system into the Sun, and a
"cluster": the Earth-Moon 'stem. This is an approximation, but it enables the system to be
understood quite deeply. So to what extent is this cluster a part of reality, and to hat
extent is it a mental fabrication, a human imposition of structure on me universe? This
problem of the "reality" of boundaries drawn between hat are perceived to be autonomous
or semi-autonomous clusters will create endless trouble when we relate it to symbols in
the brain.
  One greatly puzzling question is the simple issue of plurals. How do we visualize, say,
three dogs in a teacup? Or several people in an elevator? Do we begin with the class
symbol for "dog" and then rub three "copies" off of it? That is, do we manufacture three
fresh instance symbols using the class 'symbol "dog" as template? Or do we jointly
activate the symbols "three" and log"? By adding more or less detail to the scene being
imagined, either theory becomes hard to maintain. For instance, we certainly do not have
a separate instance symbol for each nose, mustache, grain of salt, etc., that we have ever
seen. We let class symbols take care of such numerous items, and when we pass people
on the street who have mustaches, we somehow just activate the "mustache" class
symbol, without minting fresh instance symbols, unless we scrutinize them carefully.
  On the other hand, once we begin to distinguish individuals, we cannot rely on a single
class symbol (e.g., "person") to timeshare itself among all the different people. Clearly
there must come into existence separate stance symbols for individual people. It would be
ridiculous to imagine




Brains and Thoughts                                                                     360

  that this feat could be accomplished by 'juggling"-that is, by the single class symbol
flitting back and forth between several different modes of activation (one for each
person).
  Between the extremes, there must be room for many sorts of intermediate cases. There
may be a whole hierarchy of ways of creating the class-instance distinction in the brain,
giving rise to symbols-and symbol organizations-of varying degrees of specificity. The
following different kinds of individual and joint activation of symbols might be
responsible for mental images of various degrees of specificity:

  (1) various different modes or depths of activation of a single class symbol:
  (2) simultaneous activation of several class symbols in some in some coordinated
       manner:
  (3) activation of a single instance symbol:
  (4) activation of a single instance symbol in conjunction with activation of several
       class symbols:
  (5) simultaneous activation of several instance symbols and several class symbols
       in some coordinated manner.

  This brings us right hack to the question: "When is a symbol a distinguishable
subsystem of the brain For instance, consider the second example-simultaneous
activation of several class symbols in some coordinated manner. This could easily be
what happens when "piano sonata" is the concept under consideration (the symbols for
"piano" and "sonata" being at least two of the activated symbols). But if this pair of
symbols gets activated in conjunction often enough, it is reasonable to assume that the
link between them will become strong enough that they will act as a unit, when activated
together in the proper way. So two or more symbols can act as one, under the proper
conditions, which means that the problem of enumerating the number of symbols in the
brain is trickier than one might guess.
  Sometimes conditions can arise where two previously unlinked symbols get activated
simultaneously and in a coordinated fashion. They may fit together so well that it seems
like an inevitable union, and a single new symbol is formed by the tight interaction of the
two old symbols. If this happens, would it be fair to say that the new symbol "always had
been there but never had been activated"-or should one say that it has been "created"?
  In case this sounds too abstract, let us take a concrete example: the Dialogue Crab
Canon. In the invention of this Dialogue, two existing symbols-that for "musical crab
canon", and that for "verbal dialogue “had to be activated simultaneously and in some
way forced to interact. Once this was done, 'the rest was quite inevitable: a new symbol-a
class symbol-was born from the interaction of these two, and from then on it was able to
be activated on its own. Now had it always been a dormant symbol in my brain? If so,
then it must have also been a dormant symbol in




Brains and Thoughts                                                                    361

the brain of every human who ever had its component symbols, even if it never was
awakened in them. This would mean that to enumerate the symbols in anyone's brain, one
would have to count all dormant symbols-all possible combinations and permutations of
all types of activations of all known symbols. This would even include those fantastic
creatures of software that one's brain invents when one is asleep-the strange mixtures of
ideas which wake up when their host goes to sleep ... The existence of these "potential
symbols" shows that it is really a huge oversimplification to imagine that the brain is a
well-defined collection of symbols in well-defined states of activation. It is much harder
than that to pin down a brain state on the symbol level.

                        Symbols -Software or Hardware?
  With the enormous and ever-growing repertoire of symbols that exist in each brain, you
might wonder whether there eventually comes a point when the brain is saturated-when
there is just no more room for a new symbol. This would come about, presumably, if
symbols never overlapped each other-if a given neuron never served a double function, so
that symbols would be like people getting into an elevator. "Warning: This brain has a
maximum capacity of 350,275 symbols!"
  This is not a necessary feature of the symbol model of brain function, however. In fact,
overlapping and completely tangled symbols are probably the rule, so that each neuron,
far from being a member of a unique symbol, is probably a functioning part of hundreds
of symbols. This gets a little disturbing, because if it is true, then might it not just as
easily be the case that each neuron is part of every single symbol? If that were so, then
there would be no localizability whatsoever of symbols-every symbol would be identified
with the whole of the brain. This would account for results like Lashley's cortex removal
in rats-but it would also mean abandonment of our original idea of breaking the brain up
into physically distinct subsystems. Our earlier characterization of symbols as "hardware
realizations of concepts" could at best be a great oversimplification. In fact, if every
symbol were made up of the same component neurons as every other symbol, then what
sense would it make to speak of distinct symbols at all? What would be the signature of a
given symbol's activation-that is, how could the activation of symbol A be distinguished
from the activation of symbol B? Wouldn't our whole theory go down the drain? And
even if there is not a total overlap of symbols, is our theory not more and more difficult to
maintain, the more that symbols do overlap? (One possible way of portraying
overlapping symbols is shown in Figure 68.)
  There is a way to keep a theory based on symbols even if physically, they overlap
considerably or totally. Consider the' surface of a pond, which can support many different
types of waves or ripples. The hardware namely the water itself-is the same in all cases,
but it possesses different possible modes of excitation. Such software excitations of the
same




Brains and Thoughts                                                                      362

FIGURE 68. In this schematic diagram, neurons are imagined as laid out as dots in one plane.
Two overlapping neural pathways are shown in different shades of gray. It may happen that two
independent "neural flashes" simultaneously race down these two pathways, passing through one
another like two ripples on a pond's surface (as in Fig. 52). This is illustrative of the idea of two
"active symbols" which share neurons and which may even be simultaneously activated. [From
John C. Eccles, Facing Reality (New York: Springer Verlag, 1970), p.21.]

hardware can all be distinguished from each other. By this analogy, I do not mean to go
so far as to suggest that all the different symbols are just different kinds of "waves"
propagating through a uniform neural medium which admits of no meaningful division
into physically distinct symbols. But it may be that in order to distinguish one symbol's
activation from that of another symbol, a process must be carried out which involves not
only locating the neurons which are firing, but also identifying very precise details of the
timing of the firing of those neurons. That is, which neuron preceded which other neuron,
and by how much? How many times a second was a particular neuron firing? Thus
perhaps several symbols can coexist in the same set of neurons by having different
characteristic neural firing patterns. The difference between a theory having physically
distinct symbols, and a theory having overlapping symbols which are distinguished from
each other by modes of excitation, is that the former gives hardware realizations of
concepts, while the latter gives partly hardware, partly software realizations of concepts.




Brains and Thoughts                                                                             363

                                Liftability of Intelligence
Thus we are left with two basic problems in the unraveling of thought processes, as they
take place in the brain. One is to explain how the A,-level traffic of neuron firings gives
rise to the high-level traffic of symbol activations. The other is to explain the high-level
traffic of symbol activation in its own terms-to make a theory which does not talk about
the ,v-level neural events. If this latter is possible-and it is a key assumption the basis of
all present research into Artificial Intelligence-then intelligence can be realized in other
types of hardware than brains. Then intelligence will have been shown to be a property
that can be "lifted" right out of e hardware in which it resides-or in other words,
intelligence will be a software property. This will mean that the phenomena of
consciousness and intelligence are indeed high-level in the same sense as most other
complex

FIGURE 69. The construction of an arch by workers of the termite Macrotermes belosus. Each
column is built up by the addition of pellets of soil and excrement. On the outer part of the left
column a worker is seen depositing a round fecal pellet. Other workers, having carried pellets in
their mandibles up the columns, are now placing them at the growing ends of ' columns. When a
column reaches a certain height the termites, evidently guided by odor, ;in to extend it at an angle
in the direction of a neighboring column. A completed arch is shown in the background.
[Drawing by Turid Holldobler; from E. 0. Wilson, The Insect Societies Cambridge, Mass.:
Harvard University Press, 1971), p. 230]




Brains and Thoughts                                                                            364

phenomena of nature: they have their own high-level laws which depend on, yet are
"liftable" out of, the lower levels. If, on the other hand, there is absolutely no way to
realize symbol-triggering patterns without having all the hardware of neurons (or
simulated neurons), this will imply that intelligence is a brain-bound phenomenon, and
much more difficult to unravel than one which owes its existence to a hierarchy of laws
on several different levels.
Here we come back to the mysterious collective behavior of ant colonies, which can build
huge and intricate nests, despite the fact that the roughly 100,000 neurons of an ant brain
almost certainly do not carry any. information about nest structure. How, then, does the
nest get created? Where does the information reside? In particular, ponder where the
information describing an arch such as is shown in Figure 69 can be found. Somehow, it
must be spread about in the colony, in the caste distribution, the age distribution-and
probably largely in the physical properties of the ant-body itself. That is, the interaction
between ants is determined just as much by their six-leggedness and their size and so on,
as by the information stored in their brain. Could there be an Artificial Ant Colony?

                          Can One Symbol Be Isolated?
Is it possible that one single symbol could be awakened in isolation from all others?
Probably not. Just as objects in the world always exist in a context of other objects, so
symbols are always connected to a constellation of other symbols. This does not
necessarily mean that symbols can never be disentangled from each other. To make a
rather simple analogy, males and females always arise in a species together: their roles
are completely intertwined, and yet this does not mean that a male cannot be
distinguished from a female. Each is reflected in the other, as the beads in Indra's net
reflect each other. The recursive intertwining of the functions F(n) and M(n) in Chapter V
does not prevent each function from having its own characteristics. The intertwining of F
and M could be mirrored in a pair of RTN's which call each other. From this we can jump
to a whole network of ATN's intertwined with each other-a heterarchy of interacting
recursive procedures. Here, the meshing is so inherent that no one ATN could be
activated in isolation; yet its activation may be completely distinctive, not confusable
with that of any other of the ATN's. It is not such a bad image, the brain as an ATN-
colony!
Likewise, symbols, with all their multiple links to each other, are meshed together and
yet ought to be able to be teased apart. This might involve identifying a neural network, a
network plus a mode of excitation-or possibly something of a completely different kind.
In any case, if symbols are part of reality, presumably there exists a natural way to chart
them out in a real brain. However, if some symbols were finally identified in a brain, this
would not mean that any one of them could be awakened in isolation.




Brains and Thoughts                                                                     365

        The fact that a symbol cannot be awakened in isolation does not diminish the
separate identity of the symbol; in fact, quite to the contrary: a symbol's identity lies
precisely in its ways of being connected (via potential triggering links) to other symbols.
The network by which symbols can potentially trigger each other constitutes the brain's
working model of the real universe, as well as of the alternate universes which it
considers (and which are every bit as important for the individual's survival in the real
world as the real world is).

                               The Symbols of Insects
Our facility for making instances out of classes and classes out of instances lies at the
basis of our intelligence, and it is one of the great differences between human thought and
the thought processes of other animals. Not that I have ever belonged to another species
and experienced at first hand how it feels to think their way-but from the outside it is
apparent that no other species forms general concepts as we do, or imagines hypothetical
worlds-variants on the world as it is, which aid in figuring out which future pathway to
choose. For instance, consider the celebrated "language of the bees"-information-laden
dances which are performed by worker bees returning to the hive, to inform other bees of
the location of nectar. While there may be in each bee a set of rudimentary symbols
which are activated by such a dance, there is no reason to believe that a bee has an
expandable vocabulary of symbols. Bees and other insects do not seem to have the power
to generalize-that is, to develop new class symbols from instances which we would
perceive as nearly identical.
A classic experiment with solitary wasps is reported in Dean Wooldridge's book,
Mechanical Man, from which I quote:

  When the time comes for egg laying, the wasp Sphex builds a burrow for the
  purpose and seeks out a cricket which she stings in such a way as to paralyze but not
  kill it. She drags the cricket into the burrow, lays her eggs alongside, closes the
  burrow, then flies away, never to return. In due course, the eggs hatch and the wasp
  grubs feed off the paralyzed cricket, which has not decayed, having been kept in the
  wasp equivalent of a deepfreeze. To the human mind, such an elaborately organized
  and seemingly purposeful routine conveys a convincing flavor of logic and
  thoughtfulness-until more details are examined. For example, the wasp's routine is to
  bring the paralyzed cricket to the burrow, leave it on the threshold, go inside to see
  that all is well, emerge, and then drag the cricket in. If the cricket is moved a few
  inches away while the wasp is inside making her preliminary inspection, the wasp,
  on emerging from the burrow, will bring the cricket back to the threshold, but not
  inside, and will then repeat the preparatory procedure of entering the burrow to see
  that everything is all right. If again the cricket is removed a few inches while the
  wasp is inside, once again she will move the cricket up to the threshold and reenter
  the burrow for a final check. The wasp never thinks of pulling the cricket straight in.
  On one occasion this procedure was repeated forty times, always with the same
  result.'




Brains and Thoughts                                                                    366

This seems to be completely hard-wired behavior. Now in the wasp brain, there may be
rudimentary symbols, capable of triggering each other; but there is nothing like the
human capacity to see several instances as instances of an as-yet-unformed class, and
then to make the class symbol; nor is there anything like the human ability to wonder,
"What if I did this-what would ensue in that hypothetical world\%" This type of thought
process requires an ability to manufacture instances and to manipulate them as if' they
were symbols standing for objects in a real situation, although that situation may not be
the case, and may never be the case.

                      Class Symbols and Imaginary Worlds
Let us reconsider the April Fool's joke about the borrowed car, and the images conjured
up in your mind during the telephone call. To begin with, you need to activate symbols
which represent a road, a car, a person in a car. Now the concept "road" is a very general
one, with perhaps several stock samples which you can unconsciously pull out of
dormant memory when the occasion arises. "Road" is a class, rather than an instance. As
you listen to the tale, you quickly activate symbols which are instances with gradually
increasing-specificity. For instance, when you learn that the road' was wet, this conjures
up a more specific image, though you realize that it is most likely quite different from the
actual road where the incident took place. But that is not important; what matters is
whether your symbol is sufficiently well suited for the story-that is, whether the symbols
which it can trigger are the right kind.
As the story progresses, you fill in more aspects of this road: there is a high bank against
which a car could smash. Now does this mean that you are activating the symbol for
"bank", or does it mean that you are setting some parameters in your symbol for "road
Undoubtedly both. That is, the network of neurons which represents "road" has many
different ways of firing, and you are selecting which subnetwork actually shall fire. At
the same time, you are activating the symbol for "bank", and this is probably instrumental
in the process of selecting the parameters for. "road", in that its neurons may send signals
to some of those in "road"-and vice versa. (In case this seems a little confusing, it is
because I am somewhat straddling levels of description-I am trying to set up an image of
the symbols, as well as of their component neurons.)
No less important than the nouns are the verbs, prepositions, etc: They, too, activate
symbols, which send messages back and forth to each other. There are characteristic
differences between the kinds of triggering patterns of symbols for verbs and symbols for
nouns, of course, which means that they may be physically somewhat differently
organized. For instance, nouns might have fairly localized symbols, while verbs and
prepositions might have symbols with many "tentacles" reaching all around the cortex; or
any number of other possibilities.
After the story is all over, you learn it was all untrue. The power of




Brains and Thoughts                                                                     367

"rubbing off" instances from classes, in the way that one makes rubbings from brasses in
churches, has enabled you to represent the situation, and has freed you from the need to
remain faithful to the real world. The fact that symbols can act as templates for other
symbols gives you some mental independence of reality: you can create artificial
universes, in which there can happen nonreal events with any amount of detail that you
care to imbue them with. But the class symbols themselves, from which all of this
richness springs, are deeply grounded in reality.
Usually symbols play isomorphic roles to events which seem like they could happen,
although sometimes symbols are activated which represent situations which could not
happen-for example, watches sizzling, tubas laying eggs, etc. The borderline between
what could and what could not happen is an extremely fuzzy one. As we imagine a
hypothetical event, we bring certain symbols into active states-and depending on how
well they interact (which is presumably reflected in our comfort in continuing the train of
thought), we say the event "could" or "could not" happen. Thus the terms "could" and
"could not" are extremely subjective. Actually, there is a good deal of agreement among
people about which events could or could not happen. This reflects the great amount of
mental structure which we all share-but there is a borderline area where the subjective
aspect of what kinds of hypothetical worlds we are willing to entertain is apparent. A
careful study of the kinds of imaginary events that people consider could and could not
happen would yield much insight into the triggering patterns of the symbols by which
people think.

                              Intuitive Laws of Physics
When the story has been completely told, you have built up quite an elaborate mental
model of a scene, and in this model all the objects obey physical law. This means that
physical law itself must be implicitly present in the triggering patterns of the symbols. Of
course, the phrase "physical law" here does not mean "the laws of physics as expounded
by a physicist", but rather the intuitive, chunked laws which all of us have to have in our
minds in order to survive.
A curious sidelight is that one can voluntarily manufacture mental sequences of events
which violate physical law, if one so desires. For instance, if I but suggest that you
imagine a scene with two cars approaching each other and then passing right through
each other, you won't have any trouble doing so. The intuitive physical laws can be
overridden by imaginary laws of physics; but how this overriding is done, how such
sequences of images are manufactured-indeed what any one visual image is-all of these
are deeply cloaked mysteries-inaccessible pieces of knowledge.
Needless to say, we have in our brains chunked laws not only of how inanimate objects
act, but also of how plants, animals, people and societies act-in other words, chunked
laws of biology, psychology, sociology, and so




Brains and Thoughts                                                                     368

on. All of the internal representations of such entities involve the inevitable feature of
chunked models: determinism is sacrificed for simplicity. Our representation of reality
ends up being able only to predict probabilities of ending up in certain parts of abstract
spaces of behavior-not to predict anything with the precision of physics.

                     Procedural and Declarative Knowledge
A distinction which is made in Artificial Intelligence is that between procedural and
declarative types of knowledge. A piece of knowledge is said to be declarative if it is
stored explicitly, so that not only the programmer but also the program can "read" it as if
it were in an encyclopedia or an almanac. This usually means that it is encoded locally,
not spread around. By contrast, procedural knowledge is not encoded as facts-only as
programs. A programmer may be able to peer in and say, "I see that because of these
procedures here, the program `knows' how to write English sentences"-but the program
itself may have no explicit awareness of how it writes those sentences. For instance, its
vocabulary may include none of the words "English", "sentence", and "write" at all! Thus
procedural knowledge is usually spread around in pieces, and you can't retrieve it, or
"key" on it. It is a global consequence of how the program works, not a local detail. In
other words, a piece of purely procedural knowledge is an epiphenomenon.
        In most people there coexists, along with a powerful procedural representation of
the grammar of their native language, a weaker declarative representation of it. The two
may easily be in conflict, so that a native speaker will often instruct a foreigner to say
things he himself would never say, but which agree with the declarative "book learning"
he acquired in school sometime. The intuitive or chunked laws of physics and other
disciplines mentioned earlier fall mainly on the procedural side; the knowledge that an
octopus has eight tentacles falls mainly on the declarative side.
        In between the declarative and procedural extremes, there are all possible shades.
Consider the recall of a melody. Is the melody stored in your brain, note by note? Could a
surgeon extract a winding neural filament from your brain, then stretch it straight, and
finally proceed to pinpoint along it the successively stored notes, almost as if it were a
piece of magnetic tape? If so, then melodies are stored declaratively. Or is the recall of a
melody mediated by the interaction of a large number of symbols, some of which
represent tonal relationships, others of which represent emotional qualities, others of
which represent rhythmic devices, and so on? If so, then melodies are stored
procedurally. In reality, there is probably a mixture of these extremes in the way a
melody is stored and recalled.
        It is interesting that, in pulling a melody out of memory, most people do not
discriminate as to key, so that they are as likely to sing "Happy Birthday" in the key of F-
sharp as in the key of C. This indicates that tone relationships, rather than absolute tones,
are stored. But there is no reason




Brains and Thoughts                                                                      369

that tone relationships could not be stored quite declaratively. On the other hand, some
melodies are very easy to memorize, whereas others are extremely elusive. If it were just
a matter of storing successive notes, any melody could be stored as easily as any other.
The fact that some melodies are catchy and others are not seems to indicate that the brain
has a certain repertoire of familiar patterns which are activated as the melody is heard.
So, to "play back" the melody, those patterns would have to be activated in the same
order. This returns us to the concept of symbols triggering one another, rather than a
simple linear sequence of declaratively stored notes or tone relationships.
        How does the brain know whether a piece of knowledge is stored declaratively?
For instance, suppose you are asked, "What is the population of Chicago?" Somehow the
number five million springs to mind, without your wondering, "Gee, how would I go
about counting them all?" Now suppose I ask you, "How many chairs are there in your
living room?" Here, the opposite happens-instead of trying to dredge the answer out of a
mental almanac, you immediately either go to the room and count the chairs, or you
manufacture the room in your head and count the chairs in the image of the room. The
questions were of a single type-"how many?"-yet one of them caused a piece of
declarative knowledge to be fetched, while the other one caused a procedural method of
finding the answer to be invoked. This is one example where it is clear that you have
knowledge about how you classify your own knowledge; and what is more, some of that
metaknowledge may itself be stored procedurally, so that it is used without your even
being aware of how it is done.

                                    Visual Imagery
One of the most remarkable and difficult-to-describe qualities of consciousness is visual
imagery. How do we create a visual image of our living room? Of a roaring mountain
brook? Of an orange? Even more mysterious, how do we manufacture images
unconsciously, images which guide our thoughts, giving them power and color and
depth? From what store are they fetched? What magic allows us to mesh two or three
images, hardly giving a thought as to how we should do it? Knowledge of how to do this
is among the most procedural of all, for we have almost no insight into what mental
imagery is.
         It may be that imagery is based on our ability to suppress motor activity. By this, I
mean the following. If you imagine an orange, there may occur in your cortex a set of
commands to pick it up, to smell it, to inspect it, and so on. Clearly these commands
cannot be carried out, because the orange is not there. But they can be sent along the
usual channels towards the cerebellum or other suborgans of the brain, until, at some
critical point, a "mental faucet" is closed, preventing them from actually being carried
out. Depending on how far down the line this "faucet" is situated, the images may be
more or less vivid and real-seeming. Anger can cause us to




Brains and Thoughts                                                                       370

imagine quite vividly picking up some object and throwing it, or kicking something; yet
we don't actually do so. On the other hand, we feel so "near" to actually doing so.
Probably the faucet catches the nerve impulses "at the last moment".         -
        Here is another way in which visualization points out the distinction between
accessible and inaccessible knowledge. Consider how you visualized the scene of the car
skidding on the mountain road. Undoubtedly you imagined the mountain as being much
larger than the car. Now did this happen because sometime long ago you had occasion to
note that "cars are not as big as mountains"; then you committed this statement to rote
memory: and in imagining the story, you retrieved this fact, and made use of it in
constructing your image? A most unlikely theory. Or did it happen instead as a
consequence of some introspectively inaccessible interactions of the symbols which were
activated in your brain? Obviously the latter seems far more likely. This knowledge that
cars are smaller than mountains is not a piece of rote memorization, but a piece of
knowledge which can be created by deduction. Therefore, most likely it is not stored in
any single symbol in your brain, but rather it can be produced as a result of the activation,
followed by the mutual interaction, of many symbols-for example, those for "compare",
"size", "car", "mountain", and probably, others. This means that the knowledge is stored
not explicitly, but implicitly, in a spread-about manner, rather than as a local "packet of
information". Such simple facts as relative sizes of objects have to be assembled, rather
than merely retrieved. Therefore, even in the case of a verbally accessible piece of
knowledge, there are complex inaccessible processes which mediate its coming to the
state of being ready to be said.
        We shall continue our exploration of the entities called "symbols" in different
Chapters. In Chapters XVIII and XIX, on Artificial Intelligence, we shall discuss some
possible ways of implementing active symbols in programs. And next Chapter, we shall
discuss some of the insights that our symbol-based model of brain activity give into the
comparison of brains.




Brains and Thoughts                                                                      371

                        English French German Suite
       By Lewis Carroll .. .

               ... et Frank L. Warrin .. . ..

                               . and Robert Scott

'Twas brillig, and the slithy toves
Did gyre and gimble in the wabe:
All mimsy were the borogoves,
And the mome raths outgrabe.

       I1 brilgue: les toves lubricilleux
       Se gyrent en vrillant dans le guave.
       Enmimes sont les gougebosqueux
       Et le momerade horsgrave.

               Es brillig war. Die schlichten Toven
               Wirrten and wimmelten in Waben;
               Und aller-mumsige Burggoven
               Die mohmen Rath' ausgraben.

"Beware the Jabberwock, my son!
The jaws that bite, the claws that catch!
Beware the Jubjub bird, and shun
The frumious Bandersnatch!"

       ((Garde-toi du Jaseroque, mon fits!
       La gueule qui mord; la griffe qui prend!
       Garde-toi de I'oiseau Jube, evite
       Le frumieux Band-a-prend!))

               ))Bewahre doch vor Jammerwoch!
               Die Zahne knirschen, Krallen kratzen!
               Bewahr' vor Jubjub-Vogel, vor
               Frumiosen Banderschnatzchen!)),

He took his vorpal sword in hand:
Long time the manxome foe he sought
So rested he by the Tumtum tree,
And stood awhile in thought.

       Son glaive vorpal en main, it va
       T-a la recherche du fauve manscant;
       Puis arrive a I'arbre Te-te,
       l y reste, reflechissant


English French German Suite                            372

              Er griff sein vorpals Schwertchen zu,
              Er suchte lang das manchsam' Ding;
              Dann, stehend unterm Tumtum Baum,
              Er an-zu-denken-fing.

And, as in uffish thought he stood,
The Jabberwock, with eyes of flame.
Came whiffling through the tulgey wood,
And burbled as it came!

       Pendant qu'il pense, tout uffuse,
       Le Jaseroque, a l'oeil flambant,
       Vient siblant par le bois tullegeais,
       Et burbule en venant.

              Als stand er tief in Andacht auf,
              Des Jammerwochen's Augen-feuer
              Durch turgen Wald mit Wiffek kam
              Fin burbelnd Ungeheuer!

One, two! One, two! And through and through
The vorpal blade went snicker-snack!
He left it dead, and with its head
He went galumphing back.

       Un deux, un deux, par le milieu,
       Le glaive vorpal fait pat-a-pan!
       La bete defaite, avec sa tete,
       Il rentre gallomphant.

              Eins, Zwei! Eins, Zwei! Und durch and durch
              Sein vorpals Schwert zerschnifer-schnuck,
              Da blieb es todt! Er, Kopf in Hand,
              Gelaumfig zog zuriick.

"And hast thou slain the Jabberwock?
Come to my arms, my beamish boy!
0 frabjous day! Callooh! Callay!"
He chortled in his joy.

       ((As-tu tue le Jaseroque?
       Viens a mon coeur, fils rayonnais!
       O jour frabbejais! Calleau! Callai!))
       Il cortule clans sa joie.

              ))Und schlugst Du ja den Jammerwoch?
              Umarme mich, mein Bohm'sches Kind!
              O Freuden-Tag! 0 Halloo-Schlag!((
              Er schortelt froh-gesinnt.
English French German Suite                                 373

      'Twas brillig, and the slithy toves
      Did gyre and gimble in the wabe:
      All mimsy were the borogoves,
      And the mome raths outgrabe.

             Il brilgue: les toves lubricilleux
             Se gyrent en vrillant dans le guave.
             Enmimes sont les gougebosqueux
             Et le momerade horsgrave.

                     Es brillig war. Die schlichten Toven
                     Wirrten and wimmelten in Waben:
                     Und aller-mumsige Burggoven
                     Die mohmen Rath' ausgraben.




English French German Suite                                 374


                                Minds and Thoughts

                      Can Minds Be Mapped onto Each Other?

        Now THAT WE have hypothesized the existence of very high-level active
subsystems of the brain (symbols), we may return to the matter of a possible
isomorphism, or partial isomorphism, between two brains. Instead of asking about an
isomorphism on the neural level (which surely does not exist), or on the macroscopic
suborgan level (which surely does exist but does not tell us very much), we ask about the
possibility of an isomorphism between brains on the symbol level: a correspondence
which not only maps symbols in one brain onto symbols in another brain, but also maps
triggering patterns onto triggering patterns. This means that corresponding symbols in the
two brains are linked in corresponding ways. This would be a true functional
isomorphism-the same type of isomorphism as we spoke of when trying to characterize
what it is that is invariant about all butterflies,
        It is clear from the outset that such an isomorphism does not exist between any
pair of human beings. If it did, they would be completely indistinguishable in their
thoughts; but in order for that to be true, they would have to have completely
indistinguishable memories, which would mean they would have to have led one and the
same life. Even identical twins do not approach, in the remotest degree, this ideal.
        How about a single individual\% When you look back over things which you
yourself wrote a few years ago, you think "How awful!" and smile with amusement at the
person you once were. What is worse is when you do the same thing with something you
wrote or said five minutes ago. When this happens, it shows that you do not fully
understand the person you were moments ago. The isomorphism from your brain now to
your brain then is imperfect. What, then, of the isomorphisms to other people, other
species ...
        The opposite side of the coin is shown by the power of the communication that
arises between the unlikeliest partners. Think of the barriers spanned when you read lines
of poetry penned in jail by Francois Villon, the French poet of the 1400's. Another human
being, in another era, captive in jail, speaking another language ... How can you ever
hope to have a sense of the connotations behind the facade of his words, translated into
English\% Yet a wealth of meaning comes through.
        Thus, on the one hand, we can drop all hopes of finding exactly isomorphic
software in humans, but on the other, it is clear that some people think more alike than
others do. It would seem an obvious conclu





FIGURE 70. A tiny portion of the author's "semantic network".





sion that there is some sort of partial software isomorphism connecting the brains of
people wbose style of thinking is similar-in particular, a correspondence of (1) the
repertoire of symbols, and (2) the triggering patterns of symbols

                       Comparing Different Semantic Networks

But what is a partial isomorphism? This is a most difficult question to answer. It is made
even more difficult by the fact that no one has found an adequate way to represent the
network of symbols and their triggering patterns. Sometimes a picture of a small part of
such a network of symbols is drawn, where each symbol is represented as a node into
which, and out of which, lead some arcs. The lines represent triggering relationships-in
some sense. Such figures attempt to capture something of the intuitively sensible notion
of "conceptual nearness". However, there are many different kinds of nearness, and
different ones are relevant in different contexts. A tiny portion of my own "semantic
network" is shown in Figure 70. The problem is that representing a complex
interdependency of many symbols cannot be carried out very easily with just a few lines
joining vertices.
        Another problem with such a diagram is that it is not accurate to think of a symbol
as simply "on" or "off". While this is true of neurons, it does not carry upwards, to
collections of them. In this respect, symbols are quite a bit more complicated than
neurons-as you might expect, since they are made up of many neurons. The messages that
are exchanged between symbols are more complex than the mere fact, "I am now
activated". That is more like the neuron-level messages. Each symbol can be activated in
many different ways, and the type of activation will be influential in determining which
other symbols it tries to activate. How these intertwining triggering relationships can be
represented in a pictorial manner-indeed, whether they can be at all-is not clear.
        But for the moment, suppose that issue had been solved. Suppose we now agree
that there are certain drawings of nodes, connected by links (let us say they come in
various colors, so that various types of conceptual nearness can be distinguished from
each other), which capture precisely the way in which symbols trigger other symbols.
Then under what conditions would we feel that two such drawings were isomorphic, or
nearly isomorphic? Since we are dealing with a visual representation of the network of
symbols, let us consider an analogous visual problem. How would you try to determine
whether two spiderwebs had been spun by spiders belonging to the same species? Would
you try to identify individual vertices which correspond exactly, thereby setting up an
exact map of one web onto the other, vertex by vertex, fiber by fiber, perhaps even angle
by angle? This would be a futile effort. Two webs are never exactly the same: yet there is
still some sort of "style", "form", what-have-you, that infallibly brands a given species'
web.
        In any network-like structure, such as a spiderweb, one can look at local
properties and global properties. Local properties require only a very





nearsighted observer-for example an observer who can only see one vertex at a time; and
global properties require only a sweeping vision, without attention to detail. Thus, the
overall shape of a spiderweb is a global property, whereas the average number of lines
meeting at a vertex is a local property. Suppose we agree that the most reasonable
criterion for calling two spiderwebs "isomorphic" is that they should have been spun by
spiders of the same species. Then it is interesting to ask which kind of observation-local
or global-tends to be a more reliable guide in determining whether two spiderwebs are
isomorphic. Without answering the question for spiderwebs, let us now return to the
question of the closeness-or isomorphicness, if you will-of two symbol networks.

                             Translations of "Jabberwocky"

Imagine native speakers of English, French, and German, all of whom have excellent
command of their respective native languages, and all of whom enjoy wordplay in their
own language. Would their symbol networks be similar on a local level, or on a global
level? Or is it meaningful to ask such a question? The question becomes concrete when
you look at the preceding translations of Lewis Carroll's famous "Jabberwocky".
        I chose this example because it demonstrates, perhaps better than an example in
ordinary prose, the problem of trying to find "the same node" in two different networks
which are, on some level of analysis, extremely nonisomorphic. In ordinary language, the
task of translation is more straightforward, since to each word or phrase in the original
language, there can usually be found a corresponding word or phrase in the new
language. By contrast, in a poem of this type, many "words" do not carry ordinary
meaning, but act purely as exciters of nearby symbols. However, what is nearby in one
language may be remote in another.
        Thus, in the brain of a native speaker of English, "slithy" probably activates such
symbols as "slimy", "slither", "slippery", "lithe", and "sly", to varying extents. Does
"lubricilleux" do the corresponding thing in the brain of a Frenchman? What indeed
would be "the corresponding thing"? Would it be to activate symbols which are the
ordinary translations of those words? What if there is no word, real or fabricated, which
will accomplish that? Or what if a word does exist, but is very intellectual-sounding and
Latinate ("lubricilleux"), rather than earthy and Anglo-Saxon ("slithy")? Perhaps
"huilasse" would be better than "lubricilleux"? Or does the Latin origin of the word
"lubricilleux" not make itself felt to a speaker of French in the way that it would if it were
an English word ("lubricilious", perhaps)?
        An interesting feature of the translation into French is the transposition into the
present tense. To keep it in the past would make some unnatural turns of phrase
necessary, and the present tense has a much fresher flavor in French than the past. The
translator sensed that this would be "more appropriate"-in some ill-defined yet
compelling senseand made the switch. Who can say whether remaining faithful to the
English tense would have been better?





        In the German version, the droll phrase "er an-zu-denken-fing" occurs; it does not
correspond to any English original. It is a playful reversal of words, whose flavor vaguely
resembles that of the English phrase "he out-to-ponder set", if I may hazard a reverse
translation. Most likely this funny turnabout of words was inspired by the similar playful
reversal in the English of one line earlier: "So rested he by the Tumtum tree". It
corresponds, yet doesn't correspond.
        Incidentally, why did the Tumturn tree get changed into an "arbre T6-t6" in
French? Figure it out for yourself.
        The word "manxome" in the original, whose "x" imbues it with many rich
overtones, is weakly rendered in German by "manchsam", which hack-translates into
English as "maniful". The French "manscant" also lacks the manifold overtones of
"manxome". There is no end to the interest of this kind of translation task.
        When confronted with such an example, one realizes that it is utterly impossible
to make an exact translation. Yet even in this pathologically difficult case of translation,
there seems to be some rough equivalence obtainable. Why is this so, if there really is no
isomorphism between the brains of people who will read the different versions? The
answer is that there is a kind of rough isomorphism, partly global, partly local, between
the brains of all the readers of these three poems.

                                           ASU's

An amusing geographical fantasy will give some intuition for this kind of quasi-
isomorphism. (Incidentally, this fantasy is somewhat similar to a geographical analogy
devised by M. Minsky in his article on "frames", which can be found in P. H. Winston's
book The Psychology of Computer Vision.) Imagine that you are given a strange atlas of
the USA, with all natural geological features premarked-such as rivers, mountains, lakes,
and so on-but with nary a printed word. Rivers are shown as blue lines, mountains b
color, and so on. Now you are told to convert it into a road atlas for a trip which you will
soon make. You must neatly fill in the names of all states, their boundaries, time zones,
then all counties, cities, towns, all freeways and highways and toll routes, all county
roads, all state and national parks, campgrounds, scenic areas, dams, airports, and so on
... All of this must be carried out down to the level that would appear in a detailed road
atlas. And it must be manufactured out of your own head. You are not allowed access to
any information which would help you for the duration of your task.
        You are told that it will pay off, in ways that will become clear at a later date, to
make your map as true as you can. Of course, you will begin by filling in large cities and
major roads, etc., which you know. And when you have exhausted your factual
knowledge of an area, it will be to your advantage to use your imagination to help you
reproduce at least the flavor of that area, if not its true geography, by making up fake
town names, fake populations, fake roads, fake parks, and so on. This arduous task will
take





months. To make things a little easier, you have a cartographer on hand to print
everything in neatly. The end product will be your personal map of the "Alternative
Structure of the Union"-your own personal "ASU".
       Your personal ASU will be very much like the USA in the area where you grew
up. Furthermore, wherever your travels have chanced to lead you, or wherever you have
perused maps with interest, your ASU will have spots of striking agreement with the
USA: a few small towns in North Dakota or Montana, perhaps, or the whole of
metropolitan New York, might be quite faithfully reproduced in your ASU.

                                   A Surprise Reversal

When your ASU is done, a surprise takes place. Magically, the country you have
designed comes into being, and you are transported there. A friendly committee presents
you with your favorite kind of-automobile, and explains that, "As a reward for your
designing efforts, you may now enjoy an all-expense-paid trip, at a leisurely pace, around
the good old A. S. of U. You may go wherever you want, do whatever you wish to do,
taking as long as you wish-compliments of the Geographical Society of the ASU. And-to
guide you around-here is a road atlas." To your surprise, you are given not the atlas
which you designed, but a regular road atlas of the USA.
        When you embark on your trip, all sorts of curious incidents will take place. A
road atlas is being used to guide you through a country which it only partially fits. As
long as you stick to major freeways, you will probably be able to cross the country
without gross confusions. But the moment you wander off into the byways of New
Mexico or rural Arkansas, there will be adventure in store for you. The locals will not
recognize any of the towns you're looking for, nor will they know the roads you're asking
about. They will only know the large cities you name, and even then the routes to those
cities will not be the same as are indicated on your map. It will happen occasionally that
some of the cities which are considered huge by the locals are nonexistent on your map of
the USA; or perhaps they exist, but their population according to the atlas is wrong by an
order of magnitude.

                               Centrality and Universality

What makes an ASU and the USA, which are so different in some ways, nevertheless so
similar? It is that their most important cities and routes of communication can be mapped
onto each other. The differences between them are found in the less frequently traveled
routes, the cities of smaller size, and so on. Notice that this cannot be characterized either
as a local or a global isomorphism. Some correspondences do extend down to the very
local level-for instance, in both New Yorks, the main street may be Fifth Avenue, and
there may be a Times Square in both as well-yet there may not be a single town that is
found in both Montanas. So the local-global





distinction is not relevant here. What is relevant is the centrality of the city, in terms of
economics, communication, transportation, etc. The more vital the city is, in one of these
ways, the more certain it will be to occur in both the ASU and the USA.
        In this geographic analogy, one aspect is very crucial: that there are certain
definite, absolute points of reference which will occur in nearly all ASU's: New York,
San Francisco, Chicago, and so on. From these it is then possible to orient oneself. In
other words, if we begin comparing my ASU with yours, I can use the known agreement
on big cities to establish points of reference with which I can communicate the location
of smaller cities in my ASU. And if I hypothesize a voyage from Kankakee to Fruto and
you don't know where those towns are, I can refer to something we have in common, and
thereby guide you. And if I talk about a voyage from Atlanta to Milwaukee, it may go
along different freeways or smaller roads, but the voyage itself can still be carried out in
both countries. And if you start describing a trip from Horsemilk to Janzo, I can plot out
what seems to me to be an analogous trip in my ASU, despite not having towns by those
names, as long as you constantly keep me oriented by describing your position with
respect to nearby larger towns which are found in my ASU as well as in yours.
        My roads will not be exactly the same as yours, but, with our separate maps, we
can each get from a particular part of the country to another. We can do this, thanks to the
external, predetermined geological facts mountain chains, streams, etc.-facts which were
available to us both as we worked on our maps. Without those external features, we
would have no possibility of reference points in common. For instance, if you had been
given only a map of France, and I had been given a map of Germany, and then we had
both filled them in in great detail, there would he no way to try to find "the same place"
in our fictitious lands. It is necessary to begin with identical external conditions-
otherwise nothing will match.
        Now that we have carried our geographical analogy quite far, we return to the
question of isomorphisms between brains. You might well wonder why this whole
question of brain isomorphisms has been stressed so much. What does it matter if two
brains are isomorphic, or quasi-isomorphic, or not isomorphic at all? The answer is that
we have an intuitive sense that, although other people differ from us in important ways,
they are still "the same" as we are in some deep and important ways. It would be
instructive to be able to pinpoint what this invariant core of human intelligence is, and
then to be able to describe the kinds of "embellishments" which can be added to it,
making each one of us a unique embodiment of this abstract and mysterious quality
called "intelligence".
        In our geographic analogy, cities and towns were the analogues of symbols, while
roads and highways were analogous to potential triggering paths. The fact that all ASU's
have some things in common, such as the East Coast, the West Coast, the Mississippi
River, the Great Lakes, the Rockies, and many major cities and roads is analogous to the
fact that we are all forced, by external realities, to construct certain class symbols and
trigger





ing paths in the same way. These core symbols are like the large cities, to which everyone
can make reference without ambiguity. (Incidentally, the fact that cities are localized
entities should in no way be taken as indicative that symbols in a brain are small, almost
point-like entities. They are merely symbolized in that manner in a network.)
        The fact is that a large proportion of every human's network of symbols is
universal. We simply take what is common to all of us so much for granted that it is hard
to see how much we have in common with other people. It takes the conscious effort of
imagining how much-or how little-we have in common with other types of entities, such
as stones, cars, restaurants, ants, and so forth, to make evident the large amount of
overlap that we have with randomly chosen people. What we notice about another person
immediately is not the standard overlap, because that is taken for granted as soon as we
recognize the humanity of the other person; rather, we look beyond the standard overlap
and generally find some major differences, as well as some unexpected, additional
overlap.
        Occasionally, you find that another person is missing some of what you thought
was the standard, minimal core-as if Chicago were missing from their ASU, which is
almost unimaginable. For instance, someone might not know what an elephant is, or who
is President, or that the earth is round. In such cases, their symbolic network is likely to
be so fundamentally different from your own that significant communication will be
difficult. On the other hand, perhaps this same person will share some specialized kind of
knowledge with you-such as expertise in the game of dominoes-so that you can
communicate well in a limited domain. This would be like meeting someone who comes
from the very same rural area of North Dakota as you do, so that your two ASU's
coincide in great detail over a very small region, which allows you to describe how to get
from one place to another very fluently.

              How Much Do Language and Culture Channel Thought?

If we now go back to comparing our own symbol network with those of a Frenchman and
a German, we can say that we expect them to have the standard core of class symbols,
despite the fact of different native languages. We do not expect to share highly
specialized networks with them, but we do not expect such sharing with a randomly
chosen person who shares our native language, either. The triggering patterns of people
with other languages will be somewhat different from our own, but still the major class
symbols, and the major routes between them, will be universally available, so that more
minor routes can be described with reference to them.
        Now each of our three people may in addition have some command of the
languages of the other two. What is it that marks the difference between true fluency, and
a mere ability to communicate? First of all, someone fluent in English uses most words at
roughly their- regular frequencies. A non-native speaker will have picked up some words
from





dictionaries, novels, or classes-words which at some time may have been prevalent or
preferable, but which are now far down in frequency-for example, "fetch" instead of
"get", "quite" instead of "very", etc. Though the meaning usually comes through, there is
an alien quality transmitted by the unusual choice of words.
         But suppose that a foreigner learns to use all words at roughly the normal
frequencies. Will that make his speech truly fluent? Probably not. Higher than the word
level, there is an association level, which is attached to the culture as a whole-its history,
geography, religion, children's stories, literature, technological level, and so on. For
instance, to be able to speak modern Hebrew absolutely fluently, you need to know the
Bible quite well in Hebrew, because the language draws on a stock of biblical phrases
and their connotations. Such an association level permeates each language very deeply.
Yet there is room for all sorts of variety inside fluency-otherwise the only truly fluent
speakers would be people whose thoughts were the most stereotyped possible!
         Although we should recognize the depth to which culture affects thought, we
should not overstress the role of language in molding thoughts. For instance, what we
might call two "chairs" might be perceived by a speaker of French as objects belonging to
two distinct types: "chaise" and "fauteuil" ("chair" and "armchair"). People whose native
language is French are more aware of that difference than we are-but then people who
grow up in a rural area are more aware of, say, the difference between a pickup and a
truck, than a city dweller is. A city dweller may call them both "trucks". It is not the
difference in native language, but the difference in culture (or subculture), that gives rise
to this perceptual difference.
         The relationships between the symbols of people with different native languages
have every reason to be quite similar, as far as the core is concerned, because everyone
lives in the same world. When you come down to more detailed aspects of the triggering
patterns, you will find that there is less in common. It would he like comparing rural
areas in Wisconsin in ASU's which had been made up by people who had never lived in
Wisconsin. This will be quite irrelevant, however, as long as there is sufficient agreement
on the major cities and major routes, so that there are common points of reference all
over the map.


                              Trips and Itineraries in ASU's

Without making it explicit, I have been using an image of what a "thought" is in the
ASU-analogy-namely, I have been implying that a thought corresponds to a trip. The
towns which are passed through represent the symbols which are excited. This is not a
perfect analogy, but it is quite strong. One problem with it is that when a thought recurs
in someone's mind sufficiently often, it can get chunked into a single concept. This would
correspond to quite a strange event in an ASU: a commonly taken trip would become, in
some strange fashion, a new town or city! If one is to continue to use the ASU-metaphor,
then, it is important to remember that





the cities represent not only the elementary symbols, such as those for "grass", "house",
and "car", but also symbols which get created as a result of the chunking ability of a
brain-symbols for such sophisticated concepts as "crab canon", "palindrome", or "ASU".
        Now if it is granted that the notion of taking a trip is a fair counterpart to the
notion of having a thought. then the following difficult issue comes up: virtually any
route leading from one city to a second, then to a third, and so on, can be imagined, as
long as one remembers that some intervening cities are also passed through. This would
correspond to the activation of an arbitrary sequence of symbols, one after another,
making allowance for some extra symbols-those which lie en route. Now if virtually any
sequence of symbols can be activated in any desired order, it may seem that a brain is an
indiscriminate system, which can absorb or produce any thought whatsoever. But we all
know that that is not so. In fact, there are certain kinds of thoughts which we call
knowledge, or beliefs, which play quite a different role from random fancies, or
humorously entertained absurdities. How can we characterize the difference between
dreams, passing thoughts, beliefs, and pieces of knowledge?

                   Possible, Potential, and Preposterous Pathways

There are some pathways-you can think of them as pathways either in an ASU or in a
brain-which are taken routinely in going from one place to another. There are other
pathways which can only be followed if one is led through them by the hand. These
pathways are "potential pathways", which would be followed only if special external
circumstances arose. The pathways which one relies on over and over again are pathways
which incorporate knowledge-and here I mean not only knowledge of facts (declarative
knowledge), but also knowledge of how-to's (procedural knowledge). These stable,
reliable pathways are what constitute knowledge. Pieces of knowledge merge gradually
with beliefs, which are also represented by reliable pathways, but perhaps ones which are
more susceptible to replacement if, so to speak, a bridge goes out, or there is heavy fog.
This leaves us with fancies, lies, falsities, absurdities, and other variants. These would
correspond to peculiar routes such as: New York City to Newark via Bangor, Maine and
Lubbock, Texas. They are indeed possible pathways, but ones which are not likely to be
stock routes, used in everyday voyages.
         A curious, and amusing, implication of this model is that all of the "aberrant"
kinds of thoughts listed above are composed, at rock bottom, completely out of beliefs or
pieces of knowledge. That is, any weird and snaky indirect route breaks up into a number
of non-weird, non-snaky direct stretches, and these short, straightforward symbol-
connecting routes represent simple thoughts that one can rely on-beliefs and pieces of
knowledge. On reflection, this is hardly surprising, however, since it is quite reasonable
that we should only be able to imagine fictitious things that are somehow grounded in the
realities we have experienced, no matter how





wildly they deviate from them. Dreams are perhaps just such random meanderings about
the ASU's of our minds. Locally, they make sense-but globally ...


                         Different Styles of Translating Novels

A poem like `Jabberwocky" is like an unreal journey around an ASU, hopping from one
state to another very quickly, following very curious routes. The translations convey this
aspect of the poem, rather than the precise sequence of symbols which are triggered,
although they do their best in that respect. In ordinary prose, such leaps and bounds are
not so common. However, similar problems of translation do occur. Suppose you are
translating a novel from Russian to English, and come across a sentence whose literal
translation is, "She had a bowl of borscht." Now perhaps many of your readers will have
no idea what borscht is. You could attempt to replace it by the "corresponding" item in
their culture-thus, your translation might run, "She had a bowl of Campbell's soup." Now
if you think this is a silly exaggeration, take a look at the first sentence of Dostoevsky's
novel Crime and Punishment in Russian and then in a few different English translations. I
happened to look at three different English paperback translations, and found the
following curious situation.
         The first sentence employs the street name "S. Pereulok" (as transliterated). What
is the meaning of this? A careful reader of Dostoevsky's work who knows Leningrad
(which used to be called "St. Petersburg"-or should I say "Petrograd"?) can discover by
doing some careful checking of the rest of the geography in the book (which incidentally
is also given only by its initials) that the street must be "Stoliarny Pereulok". Dostoevsky
probably wished to tell his story in a realistic way, yet not so realistically that people
would take literally the addresses at which crimes and other events were supposed to
have occurred. In any case, we have a translation problem; or to be more precise, we have
several translation problems, on several different levels.
         First of all, should we keep the initial so as to reproduce the aura of semi-mystery
which appears already in this first sentence of the book? We would get "S. Lane" ("lane"
being the standard translation of "pereulok"). None of the three translators took this tack.
However, one chose to write "S. Place". The translation of Crime and Punishment which
I read in high school took a similar option. I will never forget the disoriented feeling I
experienced when I began reading the novel and encountered those streets with only
letters for names. I had some sort of intangible malaise about the beginning of the book; I
was sure that I was missing something essential, and yet I didn't know what it was ... I
decided that all Russian novels were very weird.
         Now we could be frank with the reader (who, it may be assumed, probably won't
have the slightest idea whether the street is real or fictitious anyway!) and give him the
advantage of our modern scholarship, writing





"Stoliarny Lane" (or "Place"). This was the choice of translator number 2, who gave the
translation as "Stoliarny Place".
         What about number 3? This is the most interesting of all. This translation says
"Carpenter's Lane". And why not, indeed? After all, "stoliar" means "carpenter" and "ny"
is an adjectival ending. So now we might imagine ourselves in London, not Petrograd,
and in the midst of a situation invented by Dickens, not Dostoevsky. Is that what we
want-, Perhaps we should just read a novel by Dickens instead, with the justification that
it is "the corresponding work in English". When viewed on a sufficiently high level, it is
a "translation" of the Dostoevsky novel-in fact, the best possible one! Who needs
Dostoevsky?
         We have come all the way from attempts at great literal fidelity to the author's
style, to high-level translations of flavor. Now if this happens already in the first
sentence, can you imagine how it must go on in the rest of the book? What about the
point where a German landlady begins shouting in her German-style Russian\% How do
you translate broken Russian spoken with a German accent, into English?
         Then one may also consider the problems of how to translate slang and colloquial
modes of expression. Should one search for an "analogous" phrase, or should one settle
for a word-by-word translation? If you search for an analogous phrase, then you run the
risk of committing a "Campbell's soup" type of blunder; but if you translate every
idiomatic phrase word by word, then the English will sound alien. Perhaps this is
desirable, since the Russian culture is an alien one to speakers of English. But a speaker
of English who reads such a translation will constantly be experiencing, thanks to the
unusual turns of phrase, a sense-an artificial sense-of strangeness, which was not intended
by the author, and which is not experienced by readers of the Russian original.
         Problems such as these give one pause in considering such statements as this one,
made by Warren Weaver, one of the first advocates of translation by computer, in the late
1940's: "When I look at an article in Russian, I say, 'This is really written in English, but
it has been coded in some strange symbols. I will now proceed to decode."" Weaver's
remark simply cannot be taken literally; it must rather be considered a provocative way of
saying that there is an objectively describable meaning hidden in the symbols, or at least
something pretty close to objective; therefore, there would be no reason to suppose a
computer could not ferret it out, if sufficiently well programmed.

High-Level Comparisons between Programs

Weaver's statement is about translations between different natural languages. Let's
consider now the problem of translating between two computer languages. For instance,
suppose two people have written programs which run on different computers, and we
want to know if the two programs carry out the same task. How can we find out? We
must compare the programs. But on what level should this be done? Perhaps one program





mer wrote in a machine language, the other in a compiler language. Are two such
programs comparable? Certainly. But how to compare them? One way might be to
compile the compiler language program, producing a program in the machine language of
its home computer.
        Now we have two machine language programs. But there is another problem:
there are two computers, hence two different machine languages-and they may be
extremely different. One machine may have sixteen-bit words; the other thirty-six-bit
words. One machine may' have built-in stack-handling instructions (pushing and
popping), while the other lacks them. The differences between the hardware of the two
machines may make the two machine language programs seem incomparable-and yet we
suspect they are performing the same task, and we would Iike to see that at a glance. We
are obviously looking at the programs from much too close a distance.
        What we need to do is to step back, away from machine language, towards a
higher, more chunked view. From this vantage point, we hope we will be able to perceive
chunks of program which make each program seem rationally planned out on a global,
rather than a local, scale-that is, chunks which fit together in a way that allows one to
perceive the goals of the programmer. Let us assume that both programs were originally
written in high-level languages. Then some chunking has already been done for us. But
we will run into other troubles. There is a proliferation of such languages: Fortran, Algol,
LISP, APL, and many others. How can you compare a program written in APL with one
written in Algol: Certainly not by matching them up line by line. You will again chunk
these programs in your mind, looking for conceptual, functional units which correspond.
Thus, you are not comparing hardware, you are not comparing software-you are
comparing "etherware"-the pure concepts which lie back of the software. There is some
sort of abstract "conceptual skeleton" which must be lifted out of low levels before you
can carry out a meaningful comparison of two programs in different computer languges,
of two animals, or of two sentences in different natural languages.
        Now this brings us back to an earlier question which we asked about computers
and brains: How can we make sense of a low-level description of a computer or a brain?
Is there, in any reasonable sense, an objective way to pull a high-level description out of a
low-level one, in such complicated systems? In the case of a computer, a full display of
the contents of memory-a so-called memory dump-is easily available. Dumps were
commonly printed out in the early days of computing, when something went wrong with
a program. Then the programmer would have to go home and pore over the memory
dump for hours, trying to understand what each minuscule piece of memory represented.
In essence, the programmer would be doing the opposite of what a compiler does: he
would be translating from machine language into a higher-level language, a conceptual
language. In the end, the programmer would understand the goals of the program and
could describe it in high-level terms-for example, "This program translates novels front
Russian to English", or "This program composes an eight-voice fugue based on any
theme which is fed in".





                       High-Level Comparisons between Brains

Now our question must be investigated in the case of brains. In this case, we are asking,
"Are people's brains also capable of being 'read', on a high level? Is there some objective
description of the content of a brain?" In the Ant Fugue, the Anteater claimed to be able
to tell what Aunt Hillary was thinking about, by looking at the scurryings of her
component ants. Could some superbeing-a Neuroneater, perhaps-conceivably look down
on our neurons, chunk what it sees, and come up with an analysis of our thoughts?
        Certainly the answer must be yes, since we are all quite able to describe, in
chunked (i.e., non-neural) terms, the activity of our minds at any given time. This means
that we have a mechanism which allows us to chunk our own brain state to some rough
degree, and to give a functional description of it. To be more precise, we do not chunk all
of the brain state-we only chunk those portions of it which are active. However, if
someone asks us about a subject which is coded in a currently inactive area of our brain,
we can almost instantly gain access to the appropriate dormant area and come up with a
chunked description of it-that is, some belief on that subject. Note that we come back
with absolutely zero information on the neural level of that part of the brain: our
description is so chunked that we don't even have any idea what part of our brain it is a
description of. This can be contrasted with the programmer whose chunked description
comes from conscious analysis of every part of the memory dump.
        Now if a person can provide a chunked description of any part of his own brain,
why shouldn't an outsider too, given some nondestructive means of access to the same
brain, not only be able to chunk limited portions of the brain, but actually to give a
complete chunked description of it-in other words, a complete documentation of the
beliefs of the person whose brain is accessible? It is obvious that such a description
would have an astronomical size, but that is not of concern here. We are interested in the
question of whether, in principle, there exists a well-defined, highlevel description of a
brain, or whether, conversely, the neuron-level description-or something equally
physiological and intuitively unenlightening-is the best description that in principle
exists. Surely, to answer this question would be of the highest importance if we seek to
know whether we can ever understand ourselves.

Potential Beliefs, Potential Symbols

It is my contention that a chunked description is possible, but when we get it, all will not
suddenly be clear and light. The problem is that in order to pull a chunked description out
of the brain state, we need a language to describe our findings. Now the most appropriate
way to describe a brain, it would seem, would be to enumerate the kinds of thoughts it
could entertain, and the kinds of thoughts it could not entertain-or, perhaps, to enumerate
its beliefs and the things which it does not believe. If that is the





kind of goal we will be striving for in a chunked description, then it is easy to see what
kinds of troubles we will run up against.
         Suppose you wanted to enumerate all possible voyages that could be taken in an
ASU; there are infinitely many. How do you determine which ones are plausible, though?
Well, what does "plausible" mean? We will have precisely this kind of difficulty in trying
to establish what a "possible pathway" from symbol to symbol in a brain is. We can
imagine an upsidedown dog flying through the air with a cigar in its mouth-or a collision
between two giant fried eggs on a freeway-or any number of other ridiculous images. The
number of far-fetched pathways which can be followed in our brains is without bound,
just as is the number of insane itineraries that could be planned on an ASU. But just what
constitutes a "sane" itinerary, given an ASU? And just what constitutes a "reasonable"
thought, given a brain state? The brain state itself does not forbid anv pathway, because
for any pathway there are always circumstances which could force the following of that
pathway. The physical status of a brain, if read correctly, gives information telling not
which pathways could be followed, but rather how much resistance would be offered
along the way.
         Now in an ASU, there are many trips which could be taken along two or more
reasonable alternative routes. For example, the trip from San Francisco to New York
could go along either a northern route or a southern route. Each of them is quite
reasonable, but people tend to take them under different circumstances. Looking at a map
at a given moment in time does not tell you anything about which route will be preferable
at some remote time in the future-that depends on the external circumstances under which
the trip is to be taken. Likewise, the "reading" of a brain state will reveal that several
reasonable alternative pathways are often available, connecting a given set of symbols.
However, the trip among these symbols need not be imminent; it may be simply one of
billions of "potential" trips, all of which figure in the readout of the brain state. From this
follows an important conclusion: there is no information in the brain state itself which
tells which route will be chosen. The external circumstances will play a large determining
role in choosing the route.
         What does this imply? It implies that thoughts which clash totally may be
produced by a single brain, depending on the circumstances. And any high-level readout
of the brain state which is worth its salt must contain all such conflicting versions.
Actually this is quite obvious-that we all are bundles of contradictions, and we manage to
hang together by bringing out only one side of ourselves at a given time. The selection
cannot be predicted in advance, because the conditions which will force the selection are
not known in advance. What the brain state can provide, if properly read, is a conditional
description of the selection of routes.
         Consider, for instance, the Crab's plight, described in the Prelude. He can react in
various ways to the playing of a piece of music. Sometimes he will be nearly immune to
it, because he knows it so well. Other times, he will be quite excited by it, but this
reaction requires the right kind of triggering from the outside-for instance, the presence of
an enthusiastic listener, to





whom the work is new. Presumably, a high-level reading of the Crab's brain state would
reveal the potential thrill (and conditions which would induce it), as well as the potential
numbness (and conditions which would induce it). The brain state itself would not tell
which one would occur on the next hearing of the piece, however: it could only say, "If
such-\&-such conditions obtain, then a thrill will result; otherwise ..."
        Thus a chunked description of a brain state would give a catalogue of beliefs
which could be evoked conditionally, dependent on circumstances. Since not all possible
circumstances can be enumerated, one would have to settle for those which one thinks are
"reasonable". Furthermore, one would have to settle for a chunked description of the
circumstances themselves, since they obviously cannot-and should not-be specified down
to the atomic level! Therefore, one will not be able to make an exact, deterministic
prediction saying which beliefs will be pulled out of the brain state by a given chunked
circumstance. In summary, then, a chunked description of a brain state will consist of a
probabilistic catalogue, in which are listed those beliefs which are most likely to be
induced (and those symbols which are most likely to be activated) by various sets of
"reasonably likely" circumstances, themselves described on a chunked level. Trying to
chunk someone's beliefs without referring to context is precisely as silly as trying to
describe the range of a single person's "potential progeny" without referring to the mate.
        The same sorts of problems arise in enumerating all the symbols in a given
person's brain. There are potentially not only an infinite number of pathways in a brain,
but also an infinite number of symbols. As was pointed out, new concepts can always be
formed from old ones, and one could argue that the symbols which represent such new
concepts are merely dormant symbols in each individual, waiting to be awakened. They
may never get awakened in the person's lifetime, but it could be claimed that those
symbols are nonetheless always there, just waiting for the right circumstances to trigger
their synthesis. However, if the probability is very low, it would seem that "dormant"
would be a very unrealistic term to apply in the situation. To make this clear, try to
imagine all the "dormant dreams" which are sitting there inside your skull while you're
awake. Is it conceivable that there exists a decision procedure which could tell
"potentially dreamable themes" from "undreamable themes", given your brain State

                               Where Is the Sense of Self?

Looking back on what we have discussed, you might think to yourself, "These
speculations about brain and mind are all well and good, but what about the feelings
involved in consciousness, These symbols may trigger each other all they want, but
unless someone perceives the whole thing, there's no consciousness."
        This makes sense to our intuition on some level, but it does not make much sense
logically. For we would then be compelled to look for an





explanation of the mechanism which does the perceiving of all the active symbols, if it is
not covered by what we have described so far. Of course, a "soulist" would not have to
look any further-he would merely assert that the perceiver of all this neural action is the
soul, which cannot be described in physical terms, and that is that. However, we shall try
to give a "nonsoulist" explanation of where consciousness arises.
         Our alternative to the soulist explanation-and a disconcerting one it is, too- is to
stop at ohe symbol level and say, "This is it-this is what consciousness is. Consciousness
is that property of a system that arises whenever there exist symbols in the system which
obey triggering patterns somewhat like the ones described in the past several sections."
Put so starkly, this may seem inadequate. How does it account for the sense of "I", the
sense of' self?

                                        Subsystems

There is no reason to expect that "I", or "the self"', should not be represented by a
symbol. In fact, the symbol for the self is probably the most complex of all the symbols
in the brain. For this reason, I choose to put it on a new level of the hierarchy and call it a
subsystem, rather than a symbol. To be precise, by "subsystem", I mean a constellation of
symbols, each of which can be separately activated under the control of the subsystem
itself. The image I wish to convey of a subsystem is that it functions almost as an
independent "subbrain", equipped with its own repertoire of symbols which can trigger
each other internally. Of course, there is also much communication between the
subsystem and the "outside" world-that is, the rest of the brain. "Subsystem" is just
another name for an overgrown symbol, one which has gotten so complicated that it has
many subsymbols which interact among themselves. Thus, there is no strict level
distinction between symbols and subsystems.
        Because of the extensive links between a subsystem and the rest of the brain
(some of which will be described shortly), it would be very difficult to draw a sharp
boundary between the subsystem and the outside; but even if the border is fuzzy, the
subsystem is quite a real thing. The interesting thing about a subsystem is that, once
activated and left to its own devices, it can work on its own. Thus, two or more
subsystems of the brain of an individual may operate simultaneously. I have noticed this
happening on occasion in my own brain: sometimes I become aware that two different
melodies are running through my mind, competing for "my" attention. Somehow, each
melody is being manufactured, or "played", in a separate compartment of my brain. Each
of the systems responsible for drawing a melody out of my brain is presumably activating
a number of symbols, one after another, completely oblivious to the other system doing
the same thing. Then they both attempt to communicate with a third subsystem of my
brain-mv self'-symbol- and it is at that point that the "1" inside my brain gets wind of
what's going on: in other words, it starts picking up a chunked description of the activities
of those two subsystems.





                             Subsystems and Shared Code

Typical subsystems might be those that represent the people we know intimately. They
are represented in such a complex way in our brains that their symbols enlarge to the rank
of subsystem, becoming able to act autonomously, making use of some resources in our
brains for support. By this, I mean that a subsystem symbolizing a friend can activate
many of the symbols in my brain just as I can. For instance, I can fire up my subsystem
for a good friend and virtually feel myself in his shoes, running through thoughts which
he might have, activating symbols in sequences which reflect his thinking patterns more
accurately than my own. It could be said that my model of this friend, as embodied in a
subsystem of my brain, constitutes my own chunked description of his brain.
        Does this subsystem include, then, a symbol for every symbol which I think is in
his brain? That would be redundant. Probably the subsystem makes extensive use of
symbols already present in my brain. For instance, the symbol for "mountain" in my brain
can be borrowed by the subsystem, when it is activated. The way in which that symbol is
then used by the subsystem will not necessarily be identical to the way it is used by my
full brain. In particular, if I am talking with my friend about the Tien Shan mountain
range in Central Asia (neither of us having been there), and I know that a number of years
ago he had a wonderful hiking experience in the Alps, then my interpretation of his
remarks will be colored in part by my imported images of his earlier Alpine experience,
since I will be trying to imagine how he visualizes the area.
        In the vocabulary we have been building up in this Chapter, we could say that the
activation of" the "mountain" symbol in me is under control of my subsystem
representing him. The effect of this is to open up a different window onto to my
memories from the one which I normally use-namely, my "default option" switches from
the full range of my memories to the set of my memories of his memories. Needless to
say, my representations of his memories are only approximations to his actual memories,
which are complex modes of activation of the symbols in his brain, inaccessible to me.
        My representations of his memories are also complex modes of activation of my
own symbols-those for "primordial" concepts, such as grass, trees, snow, sky, clouds, and
so on. These are concepts which I must assume are represented in him "identically" to the
way they are in me. I must also assume a similar representation in him of even more
primordial notions: the experiences of gravity, breathing, fatigue, color, and so forth.
Less primordial but perhaps a nearly universal human quality is the enjoyment of
reaching a summit and seeing a view. Therefore, the intricate processes in my brain
which are responsible for this enjoyment can be taken over directly by the friend-
subsystem without much loss of fidelity.
        We could go on to attempt to describe how I understand an entire tale told by my
friend, a tale filled with many complexities of human relationships and mental
experiences. But our terminology would quickly become inadequate. There would be
tricky recursions connected with representa





tions in him of representations in me of representations in him of one thing and another.
If' mutual friends figured in the tale being told, I would unconsciously look for
compromises between my image of his representations of them, and my own images of
them. Pure recursion would simply be an inappropriate formalism for dealing with
symbol amalgams of this type. And I have barely scratched the surface!
We plainly lack the vocabulary today for describing the complex interactions that are
possible between symbols. So let us stop before we get bogged down.
        We should note, however, that computer systems are beginning to run into some
of the some kinds of complexity, and therefore some of these notions have been given
names. For instance, my "mountain" symbol is analogous to what in computer jargon is
called shared (or reentrant) codecode which can be used by two or more separate
timesharing programs running on a single computer. The fact that activation of one
symbol can have different results when it is part of different subsystems can be explained
by saying that its code is being processed by different interpreters. Thus, the triggering
patterns in the "mountain" symbol are not absolute; they are relative to the system within
which the symbol is activated.
        The reality of such "subbrains" may seem doubtful to some. Perhaps the
following quote from M. C. Escher, as he discusses how he creates his periodic plane-
filling drawings, will help to make clear what kind of phenomenon I am referring to:

      While drawing I sometimes feel as if' I were a spiritualist medium, controlled
      by the creatures which I am conjuring up. It is as if they themselves decide on
      the shape in which they choose to appear. They take little account of my
      critical opinion during their birth and I cannot exert much influence on the
      measure of their development. They are usually very difficult and obstinate
      creatures

        Here is a perfect example of the near-autonomy of certain subsystems of the
brain, once they are activated. Escher's subsystems seemed to him almost to be able to
override his esthetic judgment. Of course, this opinion must be taken with a grain of salt,
since those powerful subsystems came into being as a result of his many years of training
and submission to precisely the forces that molded his esthetic sensitivities. In short, it is
wrong to divorce the subsystems in Escher's brain from Escher himself or from his
esthetic judgment. They constitute a vital part of his esthetic sense, where "he" is the
complete being of the artist.

The Self-Symbol and Consciousness

A very important side effect of the self-subsystem is that it can play the role of "soul", in
the following sense: in communicating constantly with the rest of the subsystems and
symbols in the brain, it keeps track of what symbols are active, and in what way. This
means that it has to have symbols for mental activity-in other words, symbols for
symbols, and symbols for the
actions of symbols.





        Of course, this does not elevate consciousness or awareness to any "magical",
nonphysical level. Awareness here is a direct effect of the complex hardware and
software we have described. Still, despite its earthly origin, this way of describing
awareness-as the monitoring of brain activity by a subsystem of the brain itself-seems to
resemble the nearly indescribable sensation which we all know and call "consciousness".
Certainly one can see that the complexity here is enough that many unexpected effects
could be created. For instance, it is quite plausible that a computer program with this kind
of structure would make statements about itself which would have a great deal of
resemblance to statements which people commonly make about themselves. This
includes insisting that it has free will, that it is not explicable as a "sum of its parts", and
so on. (On this subject, see the article "Matter, Mind, and Models" by M. Minsky in his
book Semantic Information Processing.)
        What kind of guarantee is there that a subsystem, such as I have here postulated,
which represents the self, actually exists in our brains? Could a whole complex network
of symbols such as has been described above evolve without a self-symbol evolving,
How could these symbols and their activities play out "isomorphic" mental events to real
events in the surrounding universe, if there were no symbol for the host organism, All the
stimuli coming into the system are centered on one small mass in space. It would be quite
a glaring hole in a brain's symbolic structure not to have a symbol for the physical object
in which it is housed, and which plays a larger role in the events it mirrors than any other
object. In fact, upon reflection, it seems that the only way one could make sense of the
world surrounding a localized animate object is to understand the role of that object in
relation to the other objects around it. This necessitates the existence of a selfsymbol; and
the step from symbol to subsystem is merely a reflection of the importance of the self-
symbol', and is not a qualitative change.

                              Our First Encounter with Lucas

The Oxford philosopher J. R. Lucas (not connected with the Lucas numbers described
earlier) wrote a remarkable article in 1961, entitled "Minds, Machines, and Gödel". His
views are quite opposite to mine, and yet he manages to mix many of the same
ingredients together in coming up with his opinions. The following excerpt is quite
relevant to what we have just been discussing:

   At one's first and simplest attempts to philosophize, one becomes entangled in questions of
   whether when one knows something one knows that one knows it, and what, when one is
   thinking of oneself, is being thought about, and what is doing the thinking. After one has
   been puzzled and bruised by this problem for a long time, one learns not to press these
   questions: the concept of a conscious being is, implicitly, realized to be different from that
   of an unconscious object. In saying that a conscious being knows something, we are saying
   not onh that he knows it, but that he knows that he knows it, and that he knows that he
   knows that he knows it, and so on, as long as we care to pose the





  question: there is, we recognize, an infinity here, but it is not an infinite regress in the had
  sense, for it is the questions that peter out, as being pointless, rather than the answers. The
  questions are felt to be pointless because the concept contains within itself the idea of
  being able to go on answering such questions indefinitely. Although conscious beings have
  the power of going on, we do not wish to exhibit this simply as a succession of tasks they
  are able to perform, nor do we see the mind as an infinite sequence of selves and super-
  selves and super-super-selves. Rather, we insist that a conscious being is a unity, and
  though we talk about parts of the mind, we (to so only as a metaphor, and will not allow it
  to be taken literally.
        The paradoxes of consciousness arise because a conscious being can be aware of itself
  as well as of other things, and yet cannot really be construed as being divisible into parts. It
  means that a conscious being can deal with Gödelian questions in a was in which a
  machine cannot, because a conscious being can both consider itself and its perform a rice
  and vet not be other than that which did the performance. A machine can be made in a
  manner of speaking to "consider" its performance, but it cannot take this "into account"
  without thereby becoming a different machine, namely the old machine with a "new part"
  added. Btu it is inherent in our idea of a conscious mind that it can reflect upon itself and
  criticize its own performances, and no extra part is required to (to this: it is already
  complete, and has no Achilles' heel.
        The thesis thus begins to become more of a matter of conceptual analysis than
  mathematical discovery. This is borne out by considering another argument put forward by
  Turing. So far, we have constructed only fairly simple and predictable artifacts. When we
  increase the complexity of our machines, there may, perhaps, be surprises in store for us.
  He draws a parallel with a fission pile. Below a certain "critical" size, nothing much
  happens: but above the critical size, the sparks begin to fly. So too, perhaps, with brains
  and machines. Most brains and all machines are, at present, sub-critical"-they react to
  incoming stimuli in a stodgy and uninteresting way, have no ideas of their own, can
  produce only stock responses-but a few brains at present, and possibly some machines in
  the future, are super-critical, and scintillate on their own account. Turing is suggesting that
  it is only a matter of complexity, and that above a certain level of complexity a qualitative
  difference appears, so that "super-critical" machines will be quite unlike the simple ones
  hitherto envisaged.
        This may be so. Complexity often does introduce qualitative differences. Although it
  sounds implausible, it might turn out that above a certain level of complexity, a machine
  ceased to be predictable, even in principle, and started doing things on its own account, or,
  to use a very revealing phrase, it might begin to have a mind of its own. It might begin to
  have a mind of its own. It would begin to have a mind of its own when it was no longer
  entirely predictable and entirely docile, but was capable of doing things which we
  recognized as intelligent, and not just mistakes or random shots, but which we had not
  programmed into it. But then it would cease to be a machine, within the meaning of the act.
  What is at stake in the mechanist debate is not how minds are, or might be, brought into
  being, but how they operate. It is essential for the mechanist thesis that the mechanical
  model of the mind shall operate according to "mechanical principles," that is, that we can
  understand the operation of the whole in terms of the operations of its parts, and the
  operation of each part either shall be determined by its initial state and the construction of
  the machine, or shall be a random choice between a determinate number of determinate
  operations. If the mechanist produces a machine which is so complicated that this ceases to
  hold good of it, then it is no longer a





   machine for the purposes of our discussion, no matter how it was constructed. We should
   say, rather, that he had created a mind, in the same sort of sense as we procreate people at
   present. There would then be two ways of bringing new minds into the world, the
   traditional way, by begetting children born of women, and a new way by constructing very,
   very complicated systems of, say, valves and relays. When talking of the second way. we
   should take care to stress that although what was created looked like a machine, it was not
   one really, because it was not just the total of its parts. One could not tell what it was going
   to do merely by knowing the way in which it was built up and the initial state of its parts:
   one could not even tell the limits of what it could do, for even when presented with a
   Gödel-type question, it got the answer right. In fact we should say briefly that any system
   which was not floored by the Gödel question was eo ipso not a Turing machine, i.e. not a
   machine within the meaning of the act .3

   In reading this passage, my mind constantly boggles at the rapid succession of topics,
allusions, connotations, confusions, and conclusions. We jump from a Carrollian paradox
to Gödel to Turing to Artificial Intelligence to holism and reductionism, all in the span of
two brief pages. About Lucas one can say that he is nothing if not stimulating. In the
following Chapters, we shall come back to many of the topics touched on so tantalizingly
and fleetingly in this odd passage.





